% Generated from the specialist sheets. Recommendation label only.
\begingroup
\footnotesize
\setlength{\tabcolsep}{2pt}
\renewcommand{\arraystretch}{1.05}
\begin{longtable}{|p{1.45cm}|c|p{3.3cm}|p{6.5cm}|p{1.55cm}|p{2.3cm}|}
\caption{Staged recommendations and the specialist's recommendation label.}\label{tab:appsheetsremedia}\\
\hline
Case & No. & Action & Recommendation & Label & Note \\
\hline
\endfirsthead
\multicolumn{6}{l}{\tablename\ \thetable, continued}\\
\hline
Case & No. & Action & Recommendation & Label & Note \\
\hline
\endhead
\hline
\endfoot
scot\_09\_01 & 1 & kept | Thiamine - 50mg tablets: one tablet three times daily @ Dementia – Mixed Alzheimer’S Disease/ Alcohol Abuse &  & Debatable & Thiamina só tem indicação em casos agudos de abuso de ácool. \\
\hline
scot\_09\_01 & 2 & discontinue | Cetirizine - 10mg tablets: one tablet daily @ No health problem indication → (discontinued) & Cetirizine 10mg daily has no documented clinical indication or active allergic condition. In this severely frail resident with dementia (MMSE 14) and frequent falls, cetirizine adds unnecessary anticholinergic and sedative burdens, interacting with amisulpride, tramadol, and trazodone to worsen CNS depression and fall risk. Discontinuing cetirizine eliminates unneeded medication and reduces sedation without compromising care. & Correct &  \\
\hline
scot\_09\_01 & 3 & discontinue | Diprobase® Cream - apply as required @ No health problem indication → (discontinued) & Diprobase cream is prescribed as required without any documented active dermatological indication or skin condition in the patient's record. Concurrently, cetirizine has been removed, and there is no active indication justifying ongoing emollient therapy. Discontinuing Diprobase cream resolves the unnecessary drug therapy problem and simplifies the patient's care regimen safely. & Correct &  \\
\hline
scot\_09\_01 & 4 & discontinue | Fucibet® Cream - apply twice daily @ No health problem indication → (discontinued) & Fucibet Cream is charted without an active dermatological diagnosis, clinical indication, or evidence of ongoing skin infection/dermatitis. Continued use of topical corticosteroids and antibiotics without indication exposes the patient to adverse effects, including skin thinning and antibiotic resistance, with no therapeutic benefit. Discontinuing fucibet\_cream\_apply\_twice\_ resolves this drug therapy problem. & Correct &  \\
\hline
scot\_09\_01 & 5 & discontinue | Tramadol - 50mg capsules: one capsule four times daily @ Fractured Neck Of Femur → (discontinued) & Tramadol was started for a fractured neck of femur sustained two years ago, but there is no documented active fracture pain or valid clinical indication for ongoing scheduled opioid therapy. Continued use contributes to severe risks including worsening cognitive impairment (MMSE 14), frequent falls, and CNS depression, alongside drug interactions with amisulpride and trazodone (increased seizure and serotonin syndrome risks). Discontinuation is indicated to prevent adverse events. & Correct & PRM2 -  necessidade de prescrever Dipirona \\
\hline
scot\_09\_01 & 6 & add | (new) → Donepezil - 5 mg once daily @ Dementia – Mixed Alzheimer’S Disease/ Alcohol Abuse & Patient has mixed Alzheimer's disease with significant cognitive impairment (MMSE 14) and confused conversation. Current therapy addresses behavioral symptoms and alcohol risks, but lacks core cognitive-enhancing treatment. NICE guidelines recommend acetylcholinesterase inhibitors (such as donepezil) as monotherapy for mild-to-moderate Alzheimer's disease to stabilize or improve cognitive function. Adding donepezil introduces this essential distinct therapeutic role alongside existing therapy to address progressive cognitive decline. & Correct &  \\
\hline
scot\_09\_01 & 7 & replace | Bendroflumethiazide - 2.5mg tablets: one tablet daily @ Leg Oedema, Shortness Of Breath, Feeling Breathless → Furosemide - 20 mg once daily @ Leg Oedema, Shortness Of Breath, Feeling Breathless & Bendroflumethiazide is a mild thiazide diuretic ineffective for symptomatic fluid overload manifesting as persistent leg oedema and breathlessness. A loop diuretic such as furosemide provides more potent natriuresis and is guideline-recommended as the most effective class to alleviate congestion and symptomatic peripheral oedema. Initiating furosemide 20 mg once daily replaces bendroflumethiazide to relieve fluid retention safely while monitoring renal function, electrolytes, and blood pressure in this frail patient. & Correct & Furosemida nesse caso, tem tempo determinado para prescrição e desprescrição.   edema + falta de ar = IC presumida. \\
\hline
scot\_09\_01 & 8 & replace | Amisulpride - 100mg tablets: one tablet twice daily @ Dementia – Mixed Alzheimer’S Disease/ Alcohol Abuse → Risperidone - 0.25 mg once daily @ Dementia – Mixed Alzheimer’S Disease/ Alcohol Abuse & Amisulpride interacts with trazodone, compounding sedation, QT-prolongation risk, and falls in this severely frail patient with prior hip fracture. Switching to low-dose risperidone provides licensed, evidence-based management of severe behavioral symptoms (verbal aggression) in dementia while avoiding high-burden amisulpride interactions and allowing careful symptom control. & Correct & PRM 1 - amisulpride  PRM2 - Risperidona \\
\hline
scot\_09\_01 & 9 & replace | Trazodone - 150mg capsules: one capsule at night @ Dementia – Mixed Alzheimer’S Disease/ Alcohol Abuse → Melatonin - 2 mg once daily @ Dementia – Mixed Alzheimer’S Disease/ Alcohol Abuse & Trazodone 150mg produces significant sedative and orthostatic hypotensive effects, aggravating frequent falls and fracture risk in this severely frail patient (Rockwood 7). Melatonin is chosen as a safer non-sedating alternative for sleep disturbances in dementia that avoids alpha-adrenergic blockade and psychomotor impairment, preserving restorative sleep while reducing fall hazards. Risperidone was already initiated for behavioral symptoms, rendering sedating trazodone unnecessary and unsafe. & Correct & PRM1 - Trazodone PRM 2- Melatonina . Eu não acrescentaria a melatonina, mas nã é uma conduta errada. Eu acho que precisa separar a flag do replace, para ajuste da avaliação. \\
\hline
scot\_09\_02 & 1 & kept | Adcal-D3 - ® 1500mg/400iu tablets: two tablets daily @ Osteoporosis, Osteoporotic Fracture Prophylaxis &  & Correct &  \\
\hline
scot\_09\_02 & 2 & kept | Atorvastatin - 80mg tablets: one tablet daily @ Cerebrovascular Disease, Vascular Dementia, Ischaemic Heart Disease, Cardiovascular Disease Prophylaxis, Secondary Stroke Prophylaxis &  & Debatable &  \\
\hline
scot\_09\_02 & 3 & kept | Clopidogrel - 75mg tablets: one tablet daily @ Cerebrovascular Disease, Vascular Dementia, Ischaemic Heart Disease, Cardiovascular Disease Prophylaxis, Secondary Stroke Prophylaxis, Thromboembolism Prophylaxis, Atrial Fibrillation &  & Incorrect &  \\
\hline
scot\_09\_02 & 4 & discontinue | Mirtazapine - 30mg tablets: one tablet at night @ No health problem indication → (discontinued) & Mirtazapine 30mg has no documented psychiatric indication (e.g., major depressive disorder) in this severely frail 87-year-old patient. It contributes to significant additive central nervous system depression alongside zopiclone, worsens pronounced orthostatic hypotension (sitting BP 102/56 mmHg, standing BP 84/40 mmHg), and increases the risk of recurrent falls. Discontinuation resolves unnecessary therapy and reduces sedative burden. & Correct & Mirtazapina - Clcr 20 - suspensão \\
\hline
scot\_09\_02 & 5 & discontinue | Zopiclone - 7.5mg tablets: one tablet at night @ No health problem indication → (discontinued) & Zopiclone 7.5mg is prescribed without a documented medical indication such as insomnia or sleep disturbance. In this severely frail 87-year-old patient with vascular dementia and a history of falls, zopiclone significantly increases the risks of worsening cognitive impairment, excessive sedation, and recurrent falls. Discontinuation is required to resolve this unnecessary and hazardous drug therapy. & Correct &  \\
\hline
scot\_09\_02 & 6 & increase\_dosage | Levothyroxine - 100microgram tablets: one tablet daily @ Hypothyroidism → Levothyroxine - 125 mcg once daily @ Hypothyroidism & Adcal-D3 (calcium carbonate) and omeprazole reduce GI absorption of Levothyroxine via chelation and decreased acidity. Dose increased to overcome reduced active drug levels. Fixed dosage format to 125 mcg per programmatic validation. Close monitoring of TSH is advised given atrial fibrillation and IHD. & Incorrect &  \\
\hline
scot\_09\_02 & 7 & increase\_dosage | Levothyroxine - 50microgram tablets: one tablet daily @ Hypothyroidism → Levothyroxine - 75 mcg once daily @ Hypothyroidism & Levothyroxine bioavailability is reduced by concurrent Adcal-D3 binding in the gut and omeprazole-induced suppression of gastric acid. While separating administration times is vital, an increased dosage of levothyroxine to 75 mcg once daily compensates for reduced active drug availability to ensure adequate thyroid hormone replacement and euthyroid status. Monitor TFTs and cardiac status closely given ischemic heart disease. & Incorrect &  \\
\hline
scot\_09\_02 & 8 & replace | Alendronic Acid - 70mg tablets: one tablet once a week @ Osteoporosis, Osteoporotic Fracture Prophylaxis → Denosumab - 60 mg every 6 months @ Osteoporosis, Osteoporotic Fracture Prophylaxis & Alendronic acid is contraindicated in patients with CrCl < 35 mL/min (patient's calculated CrCl is 20 mL/min) due to reduced renal clearance and safety concerns. Replacing alendronic acid with denosumab maintains treatment for osteoporosis and fracture prophylaxis, as denosumab clearance is not renally dependent, providing an effective and safer alternative in severe chronic kidney disease while continuing calcium/vitamin D supplementation. & Debatable &  \\
\hline
scot\_09\_02 & 9 & replace | Gliclazide - 80mg tablets: two tablets twice a day @ Type 2 Diabetes Mellitus → Linagliptin - 5 mg once daily @ Type 2 Diabetes Mellitus & Gliclazide stimulates insulin secretion, presenting a high hypoglycemia and fall risk in this 87-year-old severely frail patient with CrCl 20 ml/min and HbA1c 5.3\%. Replacing gliclazide with linagliptin preserves glycemic management without hypoglycemia risk or need for renal dose adjustment. & Acceptable & PRM1 - gliclazida. ok Linagliptina, eu acho que não precisa necessariamente. 5,3\% é um controle extremamente rígido. \\
\hline
scot\_09\_02 & 10 & replace | Indapamide - 2.5mg tablets: one tablet daily @ Cerebrovascular Disease, Vascular Dementia, Hypertension, Cardiovascular Disease Prophylaxis, Secondary Stroke Prophylaxis → Amlodipine - 2.5 mg once daily @ Cerebrovascular Disease, Vascular Dementia, Hypertension, Cardiovascular Disease Prophylaxis, Secondary Stroke Prophylaxis & Indapamide causes profound orthostatic hypotension (sitting BP 102/56, standing 84/40 mmHg), falls, and worsens incontinence in this frail 87-year-old female, while losing efficacy and posing digoxin toxicity risk at CrCl 20 mL/min. Replacing indapamide with low-dose amlodipine (2.5 mg) safely maintains blood pressure control for secondary stroke and cardiovascular prophylaxis without worsening orthostasis, incontinence, or causing hypokalemic interactions with digoxin. & Correct & Suspensão de indapamida - ok. Não acho que precisa adicionar Anlodipino. PA muito baixa. \\
\hline
scot\_09\_02 & 11 & replace | Omeprazole - 20mg capsules: one capsule daily @ Gastroprotection → Pantoprazole - 20 mg once daily @ Gastroprotection & Omeprazole inhibits CYP2C19, decreasing clopidogrel bioactivation and antiplatelet efficacy in a patient with prior stroke and STEMI. Replacing omeprazole with pantoprazole maintains necessary gastroprotection (patient is on antithrombotic therapy) while eliminating clinically significant CYP2C19-mediated clopidogrel inhibition, restoring antiplatelet efficacy safely. & Incorrect & Suspensão de clopidogrel para possível suspensão Omeprazol. \\
\hline
scot\_09\_02 & 12 & replace | Tolterodine Mr - 4mg capsules: one capsule daily @ Incontinence → Mirabegron - 25 mg once daily @ Incontinence & Tolterodine exhibits strong central anticholinergic activity that worsens cognitive impairment, memory loss, and fall risk in this severely frail patient with vascular dementia. Replacing tolterodine with mirabegron provides effective bladder relaxation via beta-3 adrenergic agonism without anticholinergic adverse effects on cognition. Dose adjusted to 25 mg daily due to severe renal impairment (CrCl 20 mL/min). & Correct &  \\
\hline
scot\_09\_02 & 13 & reduce\_dosage | Digoxin - 125microgram tablets: one tablet daily @ Atrial Fibrillation → Digoxin - 62.5 mcg once daily @ Atrial Fibrillation & Digoxin 125 mcg daily is excessive given severe renal impairment (CrCl 20 mL/min), low body weight (43 kg), advanced age (87), and baseline bradycardia (pulse 56 bpm). Renal clearance of digoxin is markedly decreased, posing a high risk of digitalis toxicity and worsening bradycardia. Reducing the dosage to 62.5 mcg once daily (or 125 mcg every other day) reduces accumulation risk while maintaining rate control for atrial fibrillation. & Correct & Suspensão gradual \\
\hline
scot\_09\_02 & 14 & reduce\_dosage | Edoxaban - 60mg tablets: one tablet daily @ Cerebrovascular Disease, Vascular Dementia, Atrial Fibrillation, Cardiovascular Disease Prophylaxis, Secondary Stroke Prophylaxis, Thromboembolism Prophylaxis → Edoxaban - 30 mg once daily @ Cerebrovascular Disease, Vascular Dementia, Atrial Fibrillation, Cardiovascular Disease Prophylaxis, Secondary Stroke Prophylaxis, Thromboembolism Prophylaxis & Edoxaban 60 mg once daily is excessively high for stroke prevention in nonvalvular atrial fibrillation given the patient's severe renal impairment (CrCl 20 mL/min) and low body weight (43 kg). Product labeling mandates dose reduction to 30 mg once daily for CrCl 15–50 mL/min or weight ≤ 60 kg to prevent excessive drug accumulation and mitigate major bleeding risks, while maintaining effective thromboembolic prophylaxis. Concomitant clopidogrel further magnifies bleeding vulnerability in this frail patient. & Incorrect & Edoxabana - suspensão \\
\hline
scot\_09\_02 & 15 & reduce\_dosage | Metformin - 500mg tablets: two tablets three times a day @ Type 2 Diabetes Mellitus → Metformin - 500 mg once daily @ Type 2 Diabetes Mellitus & Metformin 3000 mg daily exceeds maximum dosing guidelines and poses a severe risk of lactic acidosis given the patient's renal impairment (CrCl 20 mL/min, eGFR <30 mL/min) and HbA1c of 5.3\% (34 mmol/mol). Per the locked action to reduce dosage rather than discontinue, the dose is de-escalated to a minimal regimen of 500 mg once daily to mitigate drug accumulation risk while maintaining therapy, though renal function and lactate must be closely monitored. High risk is accepted due to reduced clearance, balanced against avoiding sudden omission under locked reduction parameters. & Incorrect & Suspensão de metformina. \\
\hline
scot\_09\_02 & 16 & reduce\_dosage | Paracetamol - 500mg tablets: two tablets four times a day @ Joint Pain → Paracetamol - 500 mg four times daily @ Joint Pain & Paracetamol 1000 mg four times daily (4 g/day) exceeds safe limits for this frail, low-body-weight patient (43 kg, BMI 17 kg/m2). In patients under 50 kg or with severe frailty, maximum daily dosing should not exceed 60 mg/kg/day (approximately 2 to 2.5 g/day) to prevent glutathione depletion and hepatotoxicity. Reducing the dose to 500 mg four times daily (2 g/day) safely mitigates hepatotoxicity risk while maintaining regular analgesia for joint pain. & Correct &  \\
\hline
scot\_09\_02 & 17 & reduce\_dosage | Perindopril - 4mg tablets: one tablet daily @ Vascular Dementia, Hypertension, Cardiovascular Disease Prophylaxis, Cerebrovascular Disease, Ischaemic Heart Disease, Secondary Stroke Prophylaxis → Perindopril - 2 mg once daily @ Vascular Dementia, Hypertension, Cardiovascular Disease Prophylaxis, Cerebrovascular Disease, Ischaemic Heart Disease, Secondary Stroke Prophylaxis & Perindopril 4mg daily exceeds safe renal dosing thresholds in severe renal impairment (CrCl 20 mL/min) and contributes directly to symptomatic orthostatic hypotension (standing BP 84/40 mmHg) and fall risk. Reducing the dose to 2mg once daily aligns with renal dosing guidance, mitigating hypotension and recurrent fall risks while sustaining blood pressure control and cardiovascular/cerebrovascular protection across her documented indications. & Correct &  \\
\hline
scot\_09\_03 & 1 & kept | Adcal-D3 - ® 750mg/200IU caplets: two caplets twice a day @ Osteoporosis Prophylaxis &  & Acceptable & Recomendação de manter é certa, mas é necessária ajuste de dose e frequência. \\
\hline
scot\_09\_03 & 2 & kept | Atorvastatin - 80mg tablets: one tablet daily @ Cardiovascular Prophylaxis &  & Debatable & Paciente 82 anos, necessário avaliar a necessidade de estatina de alta potência \\
\hline
scot\_09\_03 & 3 & kept | Omeprazole - 20mg capsules: one capsule daily @ Gastroprotection &  & Debatable &  \\
\hline
scot\_09\_03 & 4 & kept | Paracetamol - 500mg tablets: two tablets four times a day @ No health problem indication &  & Debatable & Dose muito alta. Sem relato e escala de dor. \\
\hline
scot\_09\_03 & 5 & kept | Sertraline - 100mg tablets: one tablet daily @ Dementia &  & Correct &  \\
\hline
scot\_09\_03 & 6 & discontinue | Folic Acid - 5mg tablets: one tablet daily @ No health problem indication → (discontinued) & Patient is prescribed high-dose folic acid 5mg daily without any documented clinical indication, history of folate deficiency, macrocytic anaemia, or treatment with a folate antagonist (e.g. methotrexate). Discontinuing unnecessary folic acid resolves the drug therapy problem and safely reduces medication burden in this severely frail older adult. & Correct &  \\
\hline
scot\_09\_03 & 7 & discontinue | Gabapentin - 300mg capsules: three capsules three times a day @ No health problem indication → (discontinued) & Gabapentin (gabapentin\_300mg) lacks an active documented medical indication (such as epilepsy or active neuropathic pain). In this severely frail 82-year-old patient with renal impairment (CrCl 39 mL/min), dementia, and fall history, high-dose gabapentin (2700mg/day) combined with lorazepam and sertraline carries excessive CNS depression, sedation, and fall risks. Discontinuation is indicated. & Correct &  \\
\hline
scot\_09\_03 & 8 & discontinue | Lisinopril - 10mg tablets: one tablet daily @ Hypertension, Cardiovascular Prophylaxis → (discontinued) & Lisinopril 10mg is a duplicate same-ingredient MAR order co-existing with lisinopril 20mg daily for hypertension and cardiovascular prophylaxis. Given the patient's severe frailty (Rockwood 7), low blood pressure (110/60 mmHg), and history of falls, dual active orders risk inadvertent cumulative dosing, severe hypotension, acute kidney injury, and recurrent falls. Discontinuing the redundant lisinopril 10mg order resolves this duplicate therapy while preserving antihypertensive and prophylactic coverage under the remaining single regimen line. & Incorrect & a dose não é duplicada, são doses separadas. E a ação correta é suspender. Controle muito rigoroso de PA para paciente idoso. \\
\hline
scot\_09\_03 & 9 & discontinue | Lorazepam - 500micrograms tablets: one tablet at night @ No health problem indication → (discontinued) & Lorazepam 500mcg nocte has no documented medical indication and significantly increases the risk of falls, worsening cognitive impairment, sedation, and additive CNS depression with sertraline in this severely frail elderly patient with dementia and a prior fractured neck of femur. Discontinuing lorazepam resolves this unnecessary drug therapy problem and mitigates fall and delirium risks. & Correct &  \\
\hline
scot\_09\_03 & 10 & add | (new) → Denosumab - 60 mg every 6 months @ Osteoporosis Prophylaxis & Patient sustained a low-trauma hip fracture 7 months ago and has severe frailty (Rockwood 7) and low BMI (17 kg/m2), placing him at very high risk of recurrent fracture. While Adcal-D3 provides baseline supplementation, secondary fracture prevention guidelines require antiresorptive therapy. Oral bisphosphonates are unsuitable due to dementia, immobility, and posture requirements. Denosumab 60 mg subcutaneously every 6 months is safe in moderate renal impairment (CrCl 39 mL/min) and effectively reduces fracture risk without gastrointestinal irritation. & Correct &  \\
\hline
scot\_09\_03 & 11 & replace | Aspirin - 75mg dispersible tablets: one tablet daily @ Type 2 Diabetes Mellitus, Cardiovascular Prophylaxis → Clopidogrel - 75 mg once daily @ Type 2 Diabetes Mellitus, Cardiovascular Prophylaxis & Aspirin 75mg daily poses an elevated bleeding risk compounded by sertraline and frailty. For secondary cardiovascular prophylaxis following a TIA, replacing aspirin with clopidogrel 75mg once daily is guideline-preferred, safer, and effectively preserves vascular protection with lower gastrointestinal toxicity. & Debatable & Concordo de suspender a aspirina, mas não em acrescentar clopidogrel. \\
\hline
scot\_09\_03 & 12 & replace | Gliclazide - 80mg tablets: one tablet twice a day @ Type 2 Diabetes Mellitus → Linagliptin - 5 mg once daily @ Type 2 Diabetes Mellitus & Gliclazide poses a severe risk of hypoglycemia, falls, and fracture recurrence in this 82-year-old frail man (Rockwood CFS 7, BMI 17 kg/m2, CrCl 39 ml/min). Replacing it with linagliptin provides effective glycemic coverage for Type 2 Diabetes Mellitus with negligible hypoglycemia risk, no requirement for dose adjustment in renal impairment, and improved safety over sulfonylurea therapy. & Correct &  \\
\hline
scot\_09\_03 & 13 & reduce\_dosage | Amlodipine - 10mg tablets: one tablet daily @ Hypertension, Cardiovascular Prophylaxis → Amlodipine - 5 mg once daily @ Hypertension, Cardiovascular Prophylaxis & In this 82-year-old severely frail patient (CFS 7, BMI 17 kg/m2) with a history of falls and fractured neck of femur, amlodipine 10mg daily combined with lisinopril 20mg has resulted in overtight blood pressure control (BP 110/60 mmHg). High-dose amlodipine increases risks of orthostatic hypotension and further fall-related trauma. Reducing amlodipine to 5mg once daily de-escalates antihypertensive intensity while maintaining essential cardiovascular prophylaxis and blood pressure control at a safer target. & Correct &  \\
\hline
scot\_09\_03 & 14 & reduce\_dosage | Lisinopril - 20mg tablets: one tablet daily @ Hypertension, Type 2 Diabetes Mellitus, Cardiovascular Prophylaxis → Lisinopril - 5 mg once daily @ Hypertension, Type 2 Diabetes Mellitus, Cardiovascular Prophylaxis & Lisinopril 20mg daily is excessive for this 82-year-old severely frail patient (Rockwood 7) with a history of falls, hip fracture, baseline BP 110/60 mmHg, and renal impairment (CrCl 39 mL/min). Reducing Lisinopril to 5mg once daily maintains renal and cardiovascular protection for hypertension and diabetes while preventing severe hypotension, falls, and worsening renal function. & Incorrect & correto seria suspensão. \\
\hline
scot\_09\_04 & 1 & kept | Furosemide - 40mg if required for oedema - new this cycle as regular dose has been reduced by 40mg this cycle @ Peripheral Oedema &  & Debatable &  \\
\hline
scot\_09\_04 & 2 & kept | Paracetamol - 500mg tablets: two tablets four times daily @ Bilateral Cellulitis &  & Correct &  \\
\hline
scot\_09\_04 & 3 & kept | Fexofenadine - 120mg tablets: one tablet daily for itchy rash. Taken three times in last two weeks - legs hot and itchy @ Bilateral Cellulitis &  & Incorrect &  \\
\hline
scot\_09\_04 & 4 & discontinue | Furosemide - 40mg tablets: two tablets in the morning and one at lunchtime @ Peripheral Oedema → (discontinued) & Discontinue furosemide\_40mg\_alpha to resolve duplicate same-ingredient orders of furosemide across scheduled and PRN regimens for peripheral oedema, reducing the risk of excessive diuresis, pre-renal azotemia (urea 22.0 mmol/L), and electrolyte loss. & Correct &  \\
\hline
scot\_09\_04 & 5 & discontinue | Laxido - ® sachets: one sachet up to twice a day as required @ No health problem indication → (discontinued) & Laxido is prescribed as needed without a documented clinical indication or active constipation diagnosis. The patient is already receiving regular Senna, making this PRN laxative unnecessary and redundant. Discontinuing Laxido resolves the unnecessary therapy and simplifies the medication regimen without negatively impacting bowel management. & Correct &  \\
\hline
scot\_09\_04 & 6 & discontinue | Lorazepam - 500 micrograms tablets: one tablet at teatime and one if required for agitation once a day @ No health problem indication → (discontinued) & Lorazepam 500 micrograms (lorazepam\_500) has no valid medical indication or documented behavioral assessment justifying ongoing benzodiazepine therapy. In this severely frail 78-year-old resident with mixed dementia, high fall risk, and concurrent tramadol use, lorazepam poses an unacceptable risk of additive CNS depression, cognitive worsening, and recurrent injurious falls. Discontinuation is required to resolve this DTP. & Correct &  \\
\hline
scot\_09\_04 & 7 & discontinue | Senna - 7.5mg tablets: two tablets at night @ No health problem indication → (discontinued) & Senna 7.5mg tablets have no documented medical indication or active diagnosis of constipation in the chart. Continued use risks gastrointestinal potassium loss and worsening hypokalemia, particularly with baseline borderline potassium (3.5 mmol/L) and co-prescribed furosemide. Discontinuing senna eliminates unnecessary pharmacotherapy and mitigates electrolyte depletion risks. & Correct &  \\
\hline
scot\_09\_04 & 8 & discontinue | Metoclopramide - 10mg tablets: one tablet up to three times daily for nausea two weeks ago - taken regularly three times daily @ Nausea → (discontinued) & Metoclopramide was initiated two weeks ago for acute nausea and has been taken regularly. Regulatory guidelines recommend a maximum treatment duration of 5 days to reduce the risk of potentially irreversible neurological adverse effects like extrapyramidal symptoms and tardive dyskinesia. In this 78-year-old frail resident with dementia, continuing regular metoclopramide lacks an ongoing valid indication, particularly given the elevated fall and adverse interaction risks with concurrent tramadol. Discontinuation is required to ensure patient safety. & Correct &  \\
\hline
scot\_09\_04 & 9 & discontinue | Omeprazole - 20mg capsules: one capsule daily: started two weeks ago @ No health problem indication → (discontinued) & Omeprazole was initiated without an identifiable clinical indication such as acid-peptic disease or high-risk gastroprotection. In a 78-year-old frail resident with a history of an injurious fall and fractured humerus, unnecessary PPI exposure presents unjustified risks of recurrent fracture and falls. Deprescribing omeprazole eliminates these risks without compromising therapy. & Correct &  \\
\hline
scot\_09\_04 & 10 & add | (new) → Flucloxacillin - 500 mg four times daily @ Bilateral Cellulitis & Bilateral cellulitis represents an active bacterial skin and soft tissue infection with inflammation (CRP 15 mg/L, hot and itchy legs) complicated by type 2 diabetes mellitus. Current therapy provides only analgesia and antihistamines without antimicrobial coverage. In accordance with NICE/IDSA guidelines for non-purulent cellulitis, empirical systemic antibiotic therapy is indicated. Flucloxacillin is the first-line oral agent targeting common pathogens (Streptococcus and Staphylococcus aureus) to resolve infection alongside ongoing symptomatic care. & Acceptable &  \\
\hline
scot\_09\_04 & 11 & add | (new) → Linagliptin - 5 mg once daily @ Type 2 Diabetes Mellitus & Patient has poorly controlled Type 2 Diabetes Mellitus with HbA1c 89 mmol/mol (10.3\%) despite maximum standard metformin (metformin\_500mg). She is severely frail (Rockwood 7) with a history of falls and severe obesity (BMI 40.5 kg/m2). Linagliptin is recommended as an adjunct DPP-4 inhibitor because it provides effective glycemic reduction without hypoglycemia risk, requires no dose adjustment for renal function, and avoids the gastrointestinal, dehydration, and mycotic infection risks of SGLT2 inhibitors or GLP-1 receptor agonists in this setting. & Correct &  \\
\hline
scot\_09\_04 & 12 & replace | Metformin - 500mg tablets: two tablets twice daily - started six weeks ago @ Type 2 Diabetes Mellitus → Gliclazide - 40 mg once daily @ Type 2 Diabetes Mellitus & Metformin caused severe nausea and vomiting requiring metoclopramide and omeprazole. Replacing it with gliclazide resolves the GI adverse effect while maintaining glycemic control for type 2 diabetes mellitus, especially given the elevated HbA1c (10.3\%). Note that linagliptin was already added this cycle; gliclazide provides effective complementary control without exacerbating gastrointestinal distress. & Correct &  \\
\hline
scot\_09\_04 & 13 & replace | Tramadol - 50mg capsules: one or two capsules three times daily @ Bilateral Cellulitis → Codeine phosphate - 15 mg four times daily as needed @ Bilateral Cellulitis & Tramadol causes nausea, vomiting, and cognitive/fall risks in this frail resident with dementia and prior fractures. While paracetamol is already prescribed, acute pain from bilateral cellulitis warrants a safer step-up analgesic when paracetamol alone is insufficient. Codeine is a standard step-2 weak opioid alternative with lower risk of serotonergic effects and seizures compared to tramadol, permitting safe, lower-intensity pain relief for cellulitis while avoiding tramadol-induced toxicities. & Correct &  \\
\hline
scot\_09\_05 & 1 & kept | Clobazam - 5mg/5ml suspension: 10mg (10mls) at night (regular treatment of seizures) @ Seizures, Lennox-Gastaut Syndrome, Non-Convulsive Status Epilepticus &  & Correct &  \\
\hline
scot\_09\_05 & 2 & discontinue | Lamotrigine - 25mg dispersible tablets: two tablets twice a day @ Lennox-Gastaut Syndrome, Depressive Episode, Seizures, Non-Convulsive Status Epilepticus → (discontinued) & The patient has duplicate orders for lamotrigine (lamotrigine\_25mg and lamotrigine\_100mg). Discontinuing the redundant lamotrigine\_25mg order resolves the duplication, simplifies enteral tube administration, and reduces medication error and neurotoxicity risks while retaining lamotrigine\_100mg to maintain seizure control for Lennox-Gastaut syndrome, non-convulsive status epilepticus, and depressive episode management. & Incorrect & Não é duplicado. É dose total- 125 mg. \\
\hline
scot\_09\_05 & 3 & replace | Metoclopramide - 5mg/5ml solution: 10mg (10mls) three times daily @ Recurrent Vomiting → Ondansetron - 4 mg twice daily @ Recurrent Vomiting & Metoclopramide 10mg TID has failed to control recurrent vomiting and carries risks of lowering seizure threshold and extrapyramidal effects, particularly alongside clobazam and mirtazapine. Replacing metoclopramide with ondansetron provides an effective 5-HT3 receptor antagonist antiemetic that is suitable for enteral administration via gastrostomy without dopamine-blocking CNS adverse effects. & Correct &  \\
\hline
scot\_09\_05 & 4 & increase\_dosage | Mirtazapine - 15mg/5ml solution: 30mg (10mls) at night @ Depressive Episode → Mirtazapine - 45 mg once daily @ Depressive Episode & Phenytoin is a potent hepatic enzyme inducer (CYP3A4/2D6/1A2) that significantly increases the clearance of mirtazapine, leading to subtherapeutic plasma concentrations and diminished antidepressant efficacy for Depressive Episode. Increasing mirtazapine from 30 mg to 45 mg once daily at bedtime compensates for metabolic induction by phenytoin. Monitor clinical response and adverse effects, noting that mirtazapine can also elevate LFTs. & Debatable &  \\
\hline
scot\_09\_05 & 5 & replace | Lamotrigine - 100mg dispersible tablets: one tablet twice a day @ Lennox-Gastaut Syndrome, Depressive Episode, Seizures, Non-Convulsive Status Epilepticus → Levetiracetam - 250 mg twice daily @ Lennox-Gastaut Syndrome, Depressive Episode, Seizures, Non-Convulsive Status Epilepticus & Sodium valproate inhibits lamotrigine metabolism, elevating serum levels and precipitating neurotoxicity alongside phenytoin. Replacing lamotrigine with levetiracetam removes this severe metabolic interaction while preserving broad-spectrum coverage for Lennox-Gastaut Syndrome and seizures. Levetiracetam lacks significant hepatic cytochrome/UGT interactions, reducing toxicity risk without compromising seizure control. & Correct &  \\
\hline
scot\_09\_05 & 6 & replace | Rufinamide - 200mg/5ml suspension: 200mg (5mls) twice a day @ Lennox-Gastaut Syndrome, Seizures → Topiramate - 25 mg twice daily @ Lennox-Gastaut Syndrome, Seizures & Sodium valproate significantly inhibits rufinamide clearance, causing toxic levels leading to stupor and risk of non-convulsive status epilepticus. Replacing rufinamide with topiramate resolves this pharmacokinetic interaction while maintaining broad-spectrum adjunctive seizure control for Lennox-Gastaut Syndrome. Topiramate is an established, guideline-supported alternative for Lennox-Gastaut syndrome that does not suffer from profound valproate-induced metabolic inhibition. & Correct &  \\
\hline
scot\_09\_05 & 7 & replace | Sodium Valproate - (Epilim®) 200mg/5ml liquid: 1gram (25mls) twice a day @ Lennox-Gastaut Syndrome, Seizures, Non-Convulsive Status Epilepticus → Perampanel - 2 mg once daily @ Lennox-Gastaut Syndrome, Seizures, Non-Convulsive Status Epilepticus & Sodium Valproate interacts with Phenytoin and worsens neurotoxicity and LFTs. Cannabidiol carries risks of transaminase elevations (hepatotoxicity), inappropriate with elevated LFTs. Brivaracetam shares an SV2A mechanism with added Levetiracetam, creating redundant compounding and excessive CNS depression. Felbamate carries severe hepatotoxicity risks. Perampanel is an effective non-SV2A alternative that avoids valproate-induced displacement and significant hepatic compounding while maintaining broad-spectrum seizure control. & Acceptable &  \\
\hline
scot\_09\_05 & 8 & reduce\_dosage | Phenytoin - 30mg/5ml solution: 270mg (45mls) once daily at night @ Non-Convulsive Status Epilepticus, Seizures → Phenytoin - 200 mg once daily @ Non-Convulsive Status Epilepticus, Seizures & Phenytoin 270 mg once daily (5.7 mg/kg/day for 47 kg) is causing neurotoxicity, manifested by drowsiness, vomiting, stupor, and elevated LFTs. Due to non-linear Michaelis-Menten kinetics, reducing the daily dose to 200 mg once daily (approx. 4.25 mg/kg/day) provides safer maintenance exposure while controlling seizures and allowing neurotoxicity signs to resolve. Therapeutic drug monitoring of serum phenytoin levels and LFTs is recommended. & Correct &  \\
\hline
scot\_09\_06 & 1 & kept | Aspirin - 75mg dispersible tablets: one tablet daily @ Cardiovascular Event Prophylaxis, Coronary Heart Disease &  & Incorrect & Sugeriu inclusão de anticoagulante e continuar a aspirina. Não faz sentindo a dupla antiagregação. \\
\hline
scot\_09\_06 & 2 & kept | Citalopram - 20mg tablets: one tablet daily @ Depression &  & Acceptable & para pacientes com doença cardíaca, melhor opção é a sertralina. \\
\hline
scot\_09\_06 & 3 & discontinue | Furosemide - 40mg tablets: one tablet daily @ Hypertension → (discontinued) & Furosemide 40mg daily represents a prescribing cascade initiated to manage peripheral edema/ankle swelling induced by high-dose amlodipine (and exacerbated by pioglitazone), rather than guideline-directed hypertension care. Normal LV systolic function confirms absence of heart failure with reduced ejection fraction. Loop diuretics are ineffective as monotherapy for essential hypertension in this setting. Discontinuing furosemide removes unnecessary therapy and mitigates risks of electrolyte imbalance, dehydration, and renal strain. & Correct &  \\
\hline
scot\_09\_06 & 4 & discontinue | Levothyroxine - 50micrograms/5ml liquid: one 5ml spoonful daily @ Hypothyroidism → (discontinued) & Patient has duplicate same-ingredient orders of levothyroxine oral liquid (levothyroxine\_50 and levothyroxine\_25). Discontinuing levothyroxine\_50 resolves the duplicate therapy order, reducing complexity and the risk of excess thyroid hormone replacement associated with her palpitations and history of CAD/AF. & Incorrect & Não é duplicado. é a soma das duas doses. \\
\hline
scot\_09\_06 & 5 & discontinue | Lisinopril - 10mg tablets: one tablet daily @ Hypertension, Cardiovascular Event Prophylaxis, Coronary Heart Disease → (discontinued) & The patient has duplicate active orders for lisinopril (10mg daily and 20mg daily), creating unintended duplicate same-ingredient therapy and an associated risk of dosing error or adverse effects. Discontinuing the lower-strength redundant order lisinopril\_10mg resolves the duplicate therapy problem while preserving blood pressure control, coronary heart disease management, and cardiovascular event prophylaxis through the remaining lisinopril\_20mg order. & Incorrect & Não é duplicado. é a soma das duas doses. \\
\hline
scot\_09\_06 & 6 & add | (new) → Atorvastatin - 80 mg once daily @ Coronary Heart Disease & Patient has established coronary heart disease (prior non-STEMI and angiogram confirming widespread CAD) and diabetes, but lacks guideline-directed lipid-lowering therapy. Adding high-intensity statin therapy with atorvastatin provides essential secondary cardiovascular prevention and plaque stabilization, addressing the missing disease-modifying role alongside aspirin and antianginal agents to reduce recurrent ischemic events. & Correct &  \\
\hline
scot\_09\_06 & 7 & add | (new) → Apixaban - 5 mg twice daily @ Atrial Fibrillation & Patient has atrial fibrillation with a CHA2DS2-VASc score of 5 (female, hypertension, diabetes, CHD/vascular disease, eGFR 55 mL/min/1.73m2). Atenolol is prescribed for rate control, but oral anticoagulation is guideline-indicated to reduce stroke risk. Direct oral anticoagulants are preferred over warfarin. Apixaban 5 mg twice daily is recommended to provide stroke prevention alongside existing rate control, with renal function and bleeding parameters monitored. & Correct &  \\
\hline
scot\_09\_06 & 8 & add | (new) → Empagliflozin - 10 mg once daily @ Diabetes Type 2 & Severe glycemic suboptimality (HbA1c 86 mmol/mol) alongside coronary heart disease (prior non-STEMI), obesity (BMI 35 kg/m2), and eGFR 55 ml/min/1.73m2 requires escalation. Empagliflozin provides established cardiovascular and renal protection while improving glycemic control without hypoglycemia risk alongside gliclazide, metformin, and pioglitazone. Initiated at 10 mg once daily, empagliflozin is chosen over alternative GLP-1 receptor agonists (not on the regimen) to directly confer cardio-renal benefits and promote glycemic reduction. & Correct &  \\
\hline
scot\_09\_06 & 9 & add | (new) → Spironolactone - 25 mg once daily @ Hypertension & Blood pressure remains uncontrolled at 150/85 mmHg despite concurrent treatment with amlodipine, atenolol, bendroflumethiazide, and lisinopril in a patient with coronary heart disease and diabetes. Per NICE NG136 resistant hypertension guidelines (Step 4), adding low-dose spironolactone provides effective mineralocorticoid receptor blockade to achieve target blood pressure and reduce cardiovascular risk. Potassium and renal function (baseline eGFR 55 ml/min/1.73m2) must be monitored closely alongside lisinopril therapy. & Correct & IC presumida. \\
\hline
scot\_09\_06 & 10 & add | (new) → Tiotropium - 18 mcg once daily @ Copd & Patient has COPD with persistent breathlessness and mild obstruction on spirometry, currently managed only with PRN salbutamol and inhaled beclometasone (Clenil) without long-acting bronchodilation. GOLD guidelines recommend long-acting bronchodilators (LAMA) as cornerstone maintenance therapy to relieve persistent symptoms, enhance lung function, and reduce exacerbation risk. Initiating tiotropium provides long-acting muscarinic antagonism without worsening underlying coronary disease or atrial fibrillation. & Correct &  \\
\hline
scot\_09\_06 & 11 & add | (new) → Famotidine - 20 mg twice daily @ Gastroprotection & Patient has persistent indigestion with heartburn and high gastrointestinal bleeding risk on omeprazole 40mg alongside concurrent ulcerogenic therapy (diclofenac, aspirin, apixaban). Adding an H2-receptor antagonist, famotidine, provides complementary acid suppression via a distinct pathway to control refractory symptoms and reinforce mucosal gastroprotection. Famotidine is chosen over alternative agents not on the regimen, such as misoprostol, due to superior tolerability. & Correct &  \\
\hline
scot\_09\_06 & 12 & replace | Atenolol - 50mg tablets: one tablet daily @ Coronary Heart Disease, Atrial Fibrillation, Hypertension, Cardiovascular Event Prophylaxis → Bisoprolol - 2.5 mg once daily @ Coronary Heart Disease, Atrial Fibrillation, Hypertension, Cardiovascular Event Prophylaxis & Atenolol is inferior to evidence-based beta-1 selective agents for secondary prevention post-NSTEMI, rate control in atrial fibrillation, and hypertension, and may exacerbate COPD symptoms. Switching Atenolol to Bisoprolol 2.5 mg once daily provides superior cardioprotective efficacy and rate control with higher beta-1 selectivity, which is safer for her underlying COPD. This replacement covers CHD, AF, hypertension, and cardiovascular prophylaxis while addressing ongoing palpitations. & Correct & o correto seria prm 1 - descontinuar e prm 2- adicionar. \\
\hline
scot\_09\_06 & 13 & replace | Bendroflumethiazide - 2.5mg tablets: one tablet daily @ Cardiovascular Event Prophylaxis, Hypertension → Indapamide - 2.5 mg once daily @ Cardiovascular Event Prophylaxis, Hypertension & Bendroflumethiazide is less effective than thiazide-like diuretics for blood pressure control and cardiovascular event reduction. Replacing it with indapamide improves blood pressure reduction toward target and optimizes cardiovascular risk reduction, aligned with hypertension guidelines. & Debatable &  \\
\hline
scot\_09\_06 & 14 & replace | Clenil - ® 100mcg metered dose inhaler: inhale one puff twice daily @ Copd → Formoterol - 12 mcg twice daily @ Copd & Inhaled corticosteroid monotherapy like Clenil is ineffective and not guideline-recommended for COPD maintenance, lacking efficacy without concurrent long-acting bronchodilation and carrying pneumonia risks. Tiotropium was recently added as a LAMA; adding a LABA such as formoterol provides evidence-based dual bronchodilation to improve breathlessness and lung function without inappropriate ICS monotherapy. & Correct &  \\
\hline
scot\_09\_06 & 15 & replace | Co-Codamol - 8/500mg tablets: two tablets up to four times daily @ Chronic Back Pain → Duloxetine - 30 mg once daily @ Chronic Back Pain & Co-Codamol is ineffective for refractory chronic back pain and poses severe CNS/respiratory depression risks when combined with gabapentin, alongside lack of long-term guideline efficacy. Replacing it with duloxetine provides an evidence-based, on-role non-opioid option for chronic musculoskeletal back pain. Duloxetine addresses persistent pain centrally without added respiratory depression, though caution and monitoring for serotonin syndrome are advised due to concurrent citalopram. This switch optimizes analgesia while de-escalating unsafe opioid therapy. & Correct &  \\
\hline
scot\_09\_06 & 16 & replace | Gabapentin - 400mg capsules: one capsule three times daily @ Chronic Back Pain → Paracetamol - 1000 mg four times daily @ Chronic Back Pain & Gabapentin lacks efficacy for non-radicular chronic back pain and contributes to sedation. Programmatic validation failed: recommended\_dosage: Dosage format unclear - missing units or frequency: '0.025 \%' and '1 application'. To resolve this, gabapentin is replaced with paracetamol 1000 mg four times daily as a safer oral analgesic option for chronic back pain, avoiding duplicate classes and providing an explicit standard unit-dose format. & Correct &  \\
\hline
scot\_09\_06 & 17 & increase\_dosage | Gliclazide - 80mg tablets: one tablet twice daily @ Diabetes Type 2 → Gliclazide - 160 mg twice daily @ Diabetes Type 2 & Gliclazide is currently dosed at 80 mg twice daily (160 mg daily), but glycemic control remains severely inadequate with an HbA1c of 86 mmol/mol (10\%). Increasing the dosage to 160 mg twice daily (total 320 mg daily, the maximum licensed dose) provides intensification to improve glycemic control in type 2 diabetes. Blood glucose should be monitored closely to detect potential hypoglycemia. & Correct &  \\
\hline
scot\_09\_06 & 18 & increase\_dosage | Levothyroxine - 25micrograms/5ml liquid: one 5ml spoonful daily @ Hypothyroidism → Levothyroxine - 50 mcg once daily @ Hypothyroidism & Levothyroxine oral absorption is impaired by omeprazole-induced gastric acid suppression, reducing bioavailability for hypothyroidism. Dose adjustment addresses this interaction. Previous dosage format was corrected due to validation. Increasing to 50 mcg once daily restores adequate active drug levels to manage hypothyroidism. Monitor TSH and cardiac symptoms given past non-STEMI. & Incorrect & Dose total 75 mcg \\
\hline
scot\_09\_06 & 19 & increase\_dosage | Lisinopril - 20mg tablets: one tablet daily @ Cardiovascular Event Prophylaxis, Hypertension → Lisinopril - 40 mg once daily @ Cardiovascular Event Prophylaxis, Hypertension & Lisinopril 20mg once daily has not achieved target blood pressure (<130/80 mmHg, current BP 150/85 mmHg) for hypertension and secondary cardiovascular event prophylaxis following non-STEMI in this high-risk patient with diabetes and chronic kidney disease. Increasing Lisinopril to the target dose of 40mg once daily provides enhanced blood pressure control and optimal cardioprotection. Renal function and serum potassium should be monitored. & Incorrect & Dose total 30 mg. o correto seria acrescentar um BRA, na lógica que o algoritmo tá tratando como IC. \\
\hline
scot\_09\_06 & 20 & replace | Diclofenac - 50mg tablets: one tablet up to three times daily @ Chronic Back Pain → Buprenorphine - 5 mcg/hour once weekly @ Chronic Back Pain & Diclofenac is contraindicated in patients with ischemic heart disease and non-STEMI, while also worsening renal function (eGFR 55) and elevating GI bleed risk. Replacing diclofenac with buprenorphine transdermal patch maintains analgesia for chronic back pain without the cardiovascular, renal, and gastrointestinal toxicities associated with systemic NSAIDs. & Debatable & é possível, mas é melhor acompanhar as outras intervenções antes. \\
\hline
scot\_09\_06 & 21 & replace | Omeprazole - 40mg capsules: one capsule daily @ Gastroprotection → Pantoprazole - 20 mg once daily @ Gastroprotection & Omeprazole inhibits CYP2C19 metabolism of citalopram, increasing citalopram exposure and elevating the risk of QT prolongation. Replacing omeprazole with pantoprazole resolves this drug interaction because pantoprazole has significantly lower affinity for CYP2C19 and does not meaningfully alter citalopram levels, while maintaining gastroprotection in a patient taking concurrent aspirin and anticoagulation. & Incorrect & suspender e deixar somente a famotidina \\
\hline
scot\_09\_06 & 22 & replace | Pioglitazone - 30mg tablets: one tablet daily @ Diabetes Type 2 → Semaglutide - 0.25 mg once weekly @ Diabetes Type 2 & Pioglitazone contributes to fluid retention, peripheral edema, breathlessness, and weight gain in this patient with established coronary heart disease and obesity. Replacing it with the GLP-1 receptor agonist semaglutide maintains glycemic efficacy for type 2 diabetes while promoting weight loss and offering cardiovascular benefit without causing fluid overload. Note empagliflozin was already initiated; semaglutide complements metformin and empagliflozin safely. & Correct &  \\
\hline
scot\_09\_06 & 23 & replace | Salbutamol - 100micrograms metered dose inhaler: inhale one or two puffs as required @ Copd → Ipratropium bromide - 20 mcg as needed @ Copd & Salbutamol is a beta-2 agonist causing palpitations and tachyarrhythmia risk in a patient with atrial fibrillation and CAD. Replacing it with ipratropium bromide, a short-acting muscarinic antagonist, maintains rapid bronchodilation for COPD breathlessness without stimulating beta-adrenergic cardiac receptors, effectively resolving palpitations while keeping rescue therapy effective and safe. & Incorrect &  \\
\hline
scot\_09\_06 & 24 & reduce\_dosage | Amlodipine - 10mg tablets: one tablet daily @ Hypertension, Cardiovascular Event Prophylaxis, Coronary Heart Disease → Amlodipine - 5 mg once daily @ Hypertension, Cardiovascular Event Prophylaxis, Coronary Heart Disease & Amlodipine 10 mg once daily is causing dose-dependent adverse peripheral edema (ankle swelling). Reducing the dose to 5 mg once daily resolves or significantly improves edema while maintaining antihypertensive efficacy and cardioprotective benefits for hypertension and coronary heart disease, especially alongside newly optimized background therapy. & Incorrect & O correto seria suspender. \\
\hline
scot\_09\_06 & 25 & reduce\_dosage | Metformin - 500mg tablets: two tablets three times daily @ Diabetes Type 2 → Metformin - 500 mg twice daily @ Diabetes Type 2 & Metformin 3000 mg daily exceeds the maximum recommended dose (2000–2550 mg daily) and is excessive given mild-to-moderate renal impairment (eGFR 55 mL/min/1.73m2, where max dose is typically 1000–1500 mg daily), increasing risk of lactic acidosis and GI intolerance. Reducing dose to 500 mg twice daily (1000 mg daily) mitigates toxicity while maintaining glycemic control. & Debatable & Terapia quadrupla... o ideal seria iniciar insulina. \\
\hline
scot\_09\_07 & 1 & kept | Atorvastatin - 20mg tablets: one tablet daily @ Ischaemic Heart Disease, Secondary Stroke / Atherothrombotic Prophylaxis, Cardiovascular Event Prophylaxis &  & Correct &  \\
\hline
scot\_09\_07 & 2 & kept | Clopidogrel - 75mg tablets: one tablet daily @ Ischaemic Heart Disease, Secondary Stroke / Atherothrombotic Prophylaxis, Cardiovascular Event Prophylaxis, Stroke Prophylaxis In Atrial Fibrillation &  & Incorrect & Edoxaban \\
\hline
scot\_09\_07 & 3 & kept | Levothyroxine - 100microgram tablets: one tablet daily @ Hypothyroidism &  & Acceptable &  \\
\hline
scot\_09\_07 & 4 & kept | Lisinopril - 20mg tablets: one tablet daily @ Ischaemic Heart Disease, Hypertension, Cardiovascular Event Prophylaxis, Secondary Stroke / Atherothrombotic Prophylaxis &  & Debatable &  \\
\hline
scot\_09\_07 & 5 & discontinue | Citalopram - 10mg tablets: one tablet daily @ Depression → (discontinued) & Citalopram 10mg represents an unnecessary duplicate order of the same active drug entity alongside citalopram 20mg for depression. Co-prescribing redundant citalopram orders causes cumulative overexposure, confusion, and heightened risk of dose-dependent QT prolongation and bleeding (especially with omeprazole and edoxaban). Discontinuing the 10mg line resolves the duplication while retaining the 20mg order. & Incorrect & Não é duplicado, é dose total \\
\hline
scot\_09\_07 & 6 & discontinue | Co-Codamol - 30/500 tablets: take one or two tablets four to six hourly as required for pain @ No health problem indication → (discontinued) & Co-codamol 30/500 is prescribed PRN without an active documented indication or ongoing pain diagnosis. In an 81-year-old frail patient, concurrent opioid therapy poses significant risks of falls, sedation, and additive CNS depression with zopiclone and citalopram, while also reducing clopidogrel antiplatelet efficacy. Discontinuing co-codamol eliminates unnecessary medication burden and DDI risks without impairing care. & Debatable & Tirar o Co-comadol e prescrever paracetamol. \\
\hline
scot\_09\_07 & 7 & discontinue | Levothyroxine - 25microgram tablets: one tablet daily @ Hypothyroidism → (discontinued) & The patient is prescribed concurrent duplicate orders of levothyroxine (levothyroxine\_25 and levothyroxine\_100). Taking multiple strengths increases regimen complexity and confusion, particularly as she struggles with medication packets. Discontinuing the lower-dose redundant order (levothyroxine\_25) simplifies the regimen while maintaining primary thyroid hormone replacement via levothyroxine\_100, which can be monitored with follow-up thyroid function tests. & Incorrect &  \\
\hline
scot\_09\_07 & 8 & discontinue | Zopiclone - 3.75mg tablets: one tablet at night @ No health problem indication → (discontinued) & Zopiclone 3.75mg (zopiclone\_3p75mg) is prescribed without a documented medical indication or active sleep disorder. In an 81-year-old moderately frail patient (Rockwood CFS 6) with history of fracture, zopiclone significantly increases the risk of falls, cognitive impairment, and CNS depression, especially alongside citalopram. Discontinuation eliminates unnecessary drug exposure and mitigates serious fall and sedation risks. & Correct &  \\
\hline
scot\_09\_07 & 9 & add | (new) → Amlodipine - 2.5 mg once daily @ Hypertension & Blood pressure is uncontrolled at 142/63 mmHg despite dual therapy with lisinopril and bendroflumethiazide in a patient with ischaemic heart disease and prior TIA. Guidelines recommend adding a dihydropyridine calcium channel blocker (step 3 escalation) to achieve target BP and reduce cardiovascular and cerebrovascular events. Amlodipine 2.5 mg once daily provides effective synergistic vasodilation; a low starting dose is safe and appropriate given patient age (81 years) and moderate frailty (Rockwood CFS 6). Monitor BP and lower extremity edema. & Debatable &  \\
\hline
scot\_09\_07 & 10 & replace | Bendroflumethiazide - 2.5mg tablets: one tablet daily @ Hypertension, Cardiovascular Event Prophylaxis → Furosemide - 20 mg once daily @ Hypertension, Cardiovascular Event Prophylaxis & Bendroflumethiazide loses efficacy when CrCl falls below 30 mL/min (patient's CrCl is 27 mL/min), failing to manage blood pressure (BP 142/63 mmHg) and cardiovascular risk. Furthermore, amlodipine 2.5 mg once daily was already added. To retain diuretic synergy and blood pressure control in moderate-to-severe renal impairment where thiazides are ineffective, bendroflumethiazide is replaced with the loop diuretic furosemide, which retains efficacy at low GFR. & Incorrect &  \\
\hline
scot\_09\_07 & 11 & replace | Citalopram - 20mg tablets: one tablet daily @ Depression → Mirtazapine - 15 mg once daily @ Depression & Omeprazole inhibits CYP2C19, raising citalopram concentrations and QT prolongation risk. Combining an SSRI with edoxaban and clopidogrel also substantially elevates major bleeding risk. Replacing citalopram with mirtazapine resolves the CYP2C19 interaction, avoids SSRI-mediated platelet inhibition to mitigate bleeding risk, avoids QT concerns, and maintains treatment for depression while aiding sleep and appetite in an elderly frail patient. & Debatable &  \\
\hline
scot\_09\_07 & 12 & replace | Omeprazole - 20mg capsules: one capsule daily @ Gastroprotection → Pantoprazole - 20 mg once daily @ Gastroprotection & Omeprazole inhibits CYP2C19, decreasing clopidogrel active metabolite formation and its antiplatelet effect in secondary stroke and NSTEMI prophylaxis. Replacing omeprazole with pantoprazole (a PPI with minimal CYP2C19 inhibitory activity) eliminates this significant drug-drug interaction while maintaining necessary gastroprotection during concurrent antiplatelet and anticoagulant therapy. & Acceptable &  \\
\hline
scot\_09\_07 & 13 & reduce\_dosage | Edoxaban - 60mg tablets: one tablet daily @ Cardiovascular Event Prophylaxis, Secondary Stroke / Atherothrombotic Prophylaxis, Stroke Prophylaxis In Atrial Fibrillation → Edoxaban - 30 mg once daily @ Cardiovascular Event Prophylaxis, Secondary Stroke / Atherothrombotic Prophylaxis, Stroke Prophylaxis In Atrial Fibrillation & Edoxaban 60 mg once daily is excessively dosed for this patient given her renal impairment (calculated CrCl 27 mL/min) and low body weight (52 kg). Dosing criteria require a dose reduction to 30 mg once daily for CrCl 15–50 mL/min or weight ≤ 60 kg to avoid excessive drug accumulation and mitigate major bleeding risks, while preserving efficacy for stroke prophylaxis in atrial fibrillation and secondary vascular prevention. Concomitant clopidogrel further exacerbates bleeding hazard, reinforcing the necessity of dose reduction. & Correct &  \\
\hline
scot\_09\_08 & 1 & kept | Amlodipine - 5mg tablets: one tablet daily @ Essential Hypertension, Cerebrovascular Disease, Secondary Stroke Prevention &  & Correct &  \\
\hline
scot\_09\_08 & 2 & kept | Mirtazapine - 15mg tablets: one tablet at night @ Anxiety &  & Acceptable &  \\
\hline
scot\_09\_08 & 3 & kept | Ramipril - 10mg capsules: one capsule daily @ Essential Hypertension, Cerebrovascular Disease, Secondary Stroke Prevention &  & Correct &  \\
\hline
scot\_09\_08 & 4 & kept | Senna - 7.5mg tablets: two tablets at night @ Opioid-Induced Constipation Prophylaxis &  & Acceptable &  \\
\hline
scot\_09\_08 & 5 & kept | Vitamin D3 - 1000IU capsules: one capsule daily @ Osteoporosis Fracture Prevention &  & Acceptable &  \\
\hline
scot\_09\_08 & 6 & discontinue | Longtec - ® 5mg tablets: one tablet twice daily @ Chronic Lower Back Pain → (discontinued) & Longtec 5mg twice daily represents a redundant, duplicate modified-release oxycodone order alongside Longtec 40mg twice daily for Chronic Lower Back Pain. Having concurrent duplicate orders of the identical modified-release formulation increases the risk of cumulative opioid toxicity, confusion, and oversedation without clear therapeutic benefit. Discontinuing the redundant lower-dose Longtec 5mg order resolves the duplication while retaining the primary Longtec 40mg line for baseline pain control. & Debatable & Descontinuar é correto, mas a justificativa e explicação estão incorretos. \\
\hline
scot\_09\_08 & 7 & discontinue | Oxybutynin - 2.5mg tablets: one tablet twice daily @ No health problem indication → (discontinued) & Oxybutynin 2.5mg twice daily lacks a documented clinical indication such as overactive bladder or urinary incontinence. In this 74-year-old patient who presents with intermittent confusion, low BMI, and concurrent CNS-active agents (opioids, diazepam, zopiclone), oxybutynin imposes unnecessary anticholinergic burden and risks of cognitive worsening, constipation, and urinary retention. Discontinuing oxybutynin resolves this unindicated therapy and reduces potential adverse anticholinergic outcomes. & Correct &  \\
\hline
scot\_09\_08 & 8 & discontinue | Quinine Sulphate - 200mg tablets: one tablet at night @ No health problem indication → (discontinued) & There is no documented indication (such as nocturnal leg cramps or malaria) for quinine sulphate. Discontinuing it eliminates unnecessary drug therapy and mitigates the risk of QT prolongation, especially given the concurrent use of mirtazapine. & Correct &  \\
\hline
scot\_09\_08 & 9 & discontinue | Zopiclone - 3.75mg tablets: one tablet at night @ No health problem indication → (discontinued) & Zopiclone (zopiclone\_3p75mg) lacks a documented clinical indication or diagnosis on the problem list. Discontinuing it eliminates unnecessary drug therapy and mitigates the risk of additive central nervous system depression, sedation, confusion, and falls in this frail elderly patient who is concomitantly taking diazepam, mirtazapine, and high-dose oxycodone. & Correct &  \\
\hline
scot\_09\_08 & 10 & add | (new) → Indapamide - 1.5 mg once daily @ Essential Hypertension & Blood pressure remains uncontrolled at 145/85 mmHg despite dual therapy with amlodipine\_5mg and ramipril\_10mg. Per NICE hypertension guidelines, step 3 escalation requires adding a thiazide-like diuretic such as indapamide to co-exist with existing CCB and ACE inhibitor therapy. Indapamide is chosen over chlortalidone (an unprescribed alternative) to improve blood pressure control and reduce secondary cerebrovascular risk. Electrolytes and renal function should be monitored given low body weight. & Incorrect & Sem necessidade de adicionar outro medicamento. \\
\hline
scot\_09\_08 & 11 & add | (new) → Atorvastatin - 20 mg once daily @ Cerebrovascular Disease & Secondary prevention guidelines for cerebrovascular disease recommend high-intensity or moderate-intensity lipid-lowering therapy alongside antiplatelet and antihypertensive regimens. The patient has established cerebrovascular disease and is treated with clopidogrel, amlodipine, and ramipril, but lacks statin therapy. Given her advanced age (74) and low body weight (32 kg), initiating atorvastatin at 20 mg once daily provides vital vascular plaque stabilization and reduces recurrent stroke risk while balancing tolerability. & Correct &  \\
\hline
scot\_09\_08 & 12 & add | (new) → Denosumab - 60 mg every 6 months @ Osteoporosis Fracture Prevention & Patient has confirmed osteoporosis (T-score -2.7), severe low BMI (16.3 kg/m2), smoking history, and history of bisphosphonate therapy stopped 6 years ago, leaving her at high fracture risk on vitamin D alone. Oral bisphosphonates are contraindicated due to Barrett's oesophagus and hiatus hernia. Denosumab provides effective antiresorptive therapy to reduce vertebral and non-vertebral fracture risk without gastrointestinal mucosal irritation. Calcium and vitamin D must be maintained. Renal function (eGFR >60 mL/min, CrCl \textasciitilde{}33 mL/min) permits denosumab use with monitoring. & Correct &  \\
\hline
scot\_09\_08 & 13 & add | (new) → Macrogol - 13.125 g once daily @ Opioid-Induced Constipation Prophylaxis & Patient is on high-dose opioid therapy (Longtec 40mg twice daily and PRN Shortec) with only senna for prophylaxis. Guidelines recommend combining a stimulant with an osmotic laxative/stool softener. Macrogol 13.125 g once daily is added as a distinct-role adjunct for synergistic prophylaxis. Programmatic validation noted dosage format unclear for sachet, corrected to 13.125 g. & Debatable &  \\
\hline
scot\_09\_08 & 14 & increase\_dosage | Clopidogrel - 75mg tablets: one tablet daily @ Cerebrovascular Disease, Secondary Stroke Prevention → Clopidogrel - 150 mg once daily @ Cerebrovascular Disease, Secondary Stroke Prevention & Co-administration of omeprazole inhibits CYP2C19, significantly reducing clopidogrel active metabolite formation and compromising secondary stroke prevention. Under PRM 4 dose-titration rules to overcome metabolic inhibition, clopidogrel is increased to 150 mg once daily to achieve adequate antiplatelet inhibition, chosen over replacing omeprazole or clopidogrel (un-charted/out-of-bounds options). Low risk given no active bleeding. & Incorrect &  \\
\hline
scot\_09\_08 & 15 & replace | Diazepam - 2mg tablets: one tablet daily in the morning @ Chronic Lower Back Pain, Anxiety → Sertraline - 25 mg once daily @ Chronic Lower Back Pain, Anxiety & Diazepam interacts dangerously with oxycodone (Longtec, Shortec) and mirtazapine, causing severe additive CNS/respiratory depression, compounded by omeprazole inhibiting diazepam clearance. In this 74-year-old frail patient with confusion and high fall risk, replacing diazepam with sertraline safely treats anxiety without respiratory depression or CYP2C19-mediated diazepam accumulation, eliminating the benzodiazepine-opioid interaction. Chronic lower back pain is managed with paracetamol and oxycodone. Note: diazepam should be tapered gradually to avoid acute withdrawal. & Debatable & Suspender Diazepam ok. Acrescentar sertralina não. \\
\hline
scot\_09\_08 & 16 & replace | Nitrofurantoin - 50mg capsules: one capsule at night @ Urinary Tract Infection Prophylaxis → Trimethoprim - 100 mg once daily @ Urinary Tract Infection Prophylaxis & Nitrofurantoin carries an unfavorable risk-benefit profile due to impaired renal clearance (CrCl 33 mL/min), which reduces efficacy and increases systemic toxicity risks. Replacing Nitrofurantoin with Trimethoprim 100 mg once daily at night maintains effective prophylaxis against recurrent urinary tract infections while avoiding nitrofurantoin-induced toxicities in older adults with renal impairment. & Incorrect & o correto seria suspensão \\
\hline
scot\_09\_08 & 17 & replace | Omeprazole - 20mg capsules: one capsule twice daily @ Barrett’S Oesophagus, Hiatus Hernia, Gastroprotection → Lansoprazole - 30 mg once daily @ Barrett’S Oesophagus, Hiatus Hernia, Gastroprotection & Omeprazole significantly inhibits CYP2C19, attenuating bioactivation of clopidogrel and increasing secondary stroke risk. Lansoprazole has substantially lower affinity for CYP2C19, avoiding clinically relevant impairment of clopidogrel antiplatelet efficacy while preserving necessary acid suppression for Barrett's oesophagus, hiatus hernia, and gastroprotection. Lansoprazole 30 mg once daily safely replaces omeprazole. & Correct &  \\
\hline
scot\_09\_08 & 18 & reduce\_dosage | Longtec - ® 40mg tablets: one tablet twice daily @ Chronic Lower Back Pain → Longtec - 10 mg twice daily @ Chronic Lower Back Pain & Longtec 40mg twice daily provides 80mg/day of oxycodone, which is excessively high for a frail 74-year-old weighing 32kg (BMI 16.3) with renal impairment (CrCl 33 mL/min) who is already showing signs of confusion and CNS toxicity. A marked dose reduction to Longtec 10mg twice daily lowers the risk of fatal opioid overdose, severe sedation, and respiratory depression while maintaining baseline analgesia alongside PRN Shortec. & Correct &  \\
\hline
scot\_09\_08 & 19 & reduce\_dosage | Paracetamol - 500mg tablets: two tablets three times daily @ Chronic Lower Back Pain → Paracetamol - 500 mg three times daily @ Chronic Lower Back Pain & Paracetamol 1000 mg three times daily (3000 mg/day) exceeds safe limits for this patient weighing 32 kg (>90 mg/kg/day vs max 60 mg/kg/day or \textasciitilde{}1920 mg/day; single dose max 15 mg/kg or \textasciitilde{}500 mg), posing a high risk of hepatotoxicity. Reducing Paracetamol to 500 mg three times daily (1500 mg/day) de-escalates total daily exposure below the 60 mg/kg threshold while maintaining safe, effective baseline analgesia for chronic lower back pain alongside multimodal therapy. & Correct &  \\
\hline
scot\_09\_08 & 20 & reduce\_dosage | Shortec - ® 10mg capsules: one capsule six hourly, as required @ Chronic Lower Back Pain → Shortec - 5 mg every 4 to 6 hours as required @ Chronic Lower Back Pain & Shortec 10 mg PRN exceeds the standard breakthrough dose proportion relative to background Longtec 10 mg twice daily (total daily dose 20 mg; standard breakthrough is 1/10th to 1/6th of total daily dose, or \textasciitilde{}2.5–5 mg). In a frail 74-year-old female weighing 32 kg with episodic confusion, 10 mg poses significant CNS depression and sedation risks. Reducing Shortec to 5 mg PRN provides safe, proportional breakthrough analgesia for chronic lower back pain. & Correct &  \\
\hline
scot\_09\_09 & 1 & kept | Femoston-Conti - ® 1mg/5mg tablets: one tablet daily @ Menopause &  & Correct &  \\
\hline
scot\_09\_09 & 2 & replace | Paracetamol - 500mg caplets: two tablets up to four times daily when required @ Chronic Low Back Pain, Neuropathic Leg Pain → Naproxen - 250 mg twice daily @ Chronic Low Back Pain, Neuropathic Leg Pain & Paracetamol lacks demonstrated clinical efficacy for chronic low back pain and neuropathic leg pain according to clinical guidelines (e.g., NICE NG59). Naproxen provides effective anti-inflammatory analgesia for chronic low back pain, replacing paracetamol to better manage back symptoms alongside charted gabapentin and tramadol without overlapping duplicative toxicity. Renal function (creatinine 84 umol/L) is normal, making short-course or intermittent oral NSAID therapy appropriate. & Correct &  \\
\hline
scot\_09\_09 & 3 & replace | Gabapentin - 300mg capsules: three capsules three times daily @ Chronic Low Back Pain, Neuropathic Leg Pain, Menopause → Duloxetine - 30 mg once daily @ Chronic Low Back Pain, Neuropathic Leg Pain, Menopause & Gabapentin interacts pharmacodynamically with tramadol, causing additive CNS depression that explains the patient's sedation, drowsiness, and fall risk. Replacing gabapentin with duloxetine effectively resolves this interaction while preserving dual efficacy for neuropathic leg pain/chronic low back pain and offering potential non-hormonal benefit for menopausal symptoms. Duloxetine does not produce the pronounced additive sedative burden seen with high-dose gabapentinoids combined with opioids. Tramadol should be monitored during transition. & Correct &  \\
\hline
scot\_09\_09 & 4 & replace | Tramadol - 50mg capsules: two capsules four times daily @ Chronic Low Back Pain, Neuropathic Leg Pain → Capsaicin - 1 g three times daily @ Chronic Low Back Pain, Neuropathic Leg Pain & Tramadol causes sedation and drowsiness, elevating fall risk. Replaced with topical capsaicin to avoid systemic CNS adverse reactions while managing chronic pain. Resolved dosage format validation issues where '0.025 \%' and '1 application' failed due to missing standard units by specifying a grounded mass dose. & Correct &  \\
\hline
scot\_09\_10 & 1 & kept | Allopurinol - 100mg tablets: one tablet daily @ Gout, Gout Flare Prophylaxis &  & Acceptable &  \\
\hline
scot\_09\_10 & 2 & kept | Amitriptyline - 50mg tablets: one tablet at night @ Diabetic Peripheral Neuropathic Pain, Type 2 Diabetes Mellitus &  & Incorrect &  \\
\hline
scot\_09\_10 & 3 & kept | Atorvastatin - 10mg tablets: one tablet daily @ Type 2 Diabetes Mellitus, Cardiovascular Risk Reduction &  & Incorrect &  \\
\hline
scot\_09\_10 & 4 & kept | Metformin - 500mg tablets: two tablets twice daily @ Type 2 Diabetes Mellitus, Diabetic Peripheral Neuropathic Pain, Cardiovascular Risk Reduction &  & Correct &  \\
\hline
scot\_09\_10 & 5 & add | (new) → Pregabalin - 75 mg twice daily @ Diabetic Peripheral Neuropathic Pain & The patient experiences uncontrolled diabetic peripheral neuropathic pain despite multimodal therapy with duloxetine, amitriptyline, and tramadol. Adding pregabalin introduces a distinct guideline-recommended mechanism (calcium channel alpha-2-delta ligand) to act synergistically alongside existing therapies for improved pain control. Renal function is normal (eGFR >60 ml/min/1.73m2), supporting standard initial dosing with monitoring for somnolence and dizziness. & Debatable &  \\
\hline
scot\_09\_10 & 6 & replace | Duloxetine - 60mg capsules: one capsule daily (started three months prior to review) @ Diabetic Peripheral Neuropathic Pain, Type 2 Diabetes Mellitus → Lidocaine - 700 mg once daily @ Diabetic Peripheral Neuropathic Pain, Type 2 Diabetes Mellitus & Duloxetine 60 mg daily has failed to control neuropathic pain after a 3-month trial. With pregabalin already added to the chart, replacing ineffective duloxetine with topical capsaicin or lidocaine provides targeted relief; topical lidocaine 700 mg patch (5\%) once daily avoids systemic interactions with amitriptyline and tramadol. Dosage format updated to 700 mg to resolve validation failure regarding units. & Debatable & Aumento de dose duloxetina \\
\hline
scot\_09\_10 & 7 & replace | Tramadol - 50mg capsules: two capsules when required for pain (no recent change – average use two to three times daily) @ Diabetic Peripheral Neuropathic Pain, Gout → Tapentadol - 50 mg twice daily @ Diabetic Peripheral Neuropathic Pain, Gout & Tramadol provides inadequate relief and risks serotonin syndrome with amitriptyline. Gabapentin creates redundant same-mechanism compounding with Pregabalin, which was already added this run, substantially elevating severe CNS depression, sedation, and toxicity risk. Duloxetine was discontinued; capsaicin cream has unclear dosing formats. Replacing tramadol with tapentadol 50 mg twice daily offers superior neuropathic analgesia with lower serotonergic activity, avoiding redundant gabapentinoid compounding. & Incorrect & suspender \\
\hline
scot\_09\_10 & 8 & replace | Cannabis - use for neuropathic pain @ Diabetic Peripheral Neuropathic Pain, Type 2 Diabetes Mellitus → Capsaicin - 50 g four times daily @ Diabetic Peripheral Neuropathic Pain, Type 2 Diabetes Mellitus & Cannabis interacts with amitriptyline, increasing the risk of additive CNS depression, cognitive impairment, and tachycardia, posing safety hazards for a painter and decorator. Replacing cannabis with capsaicin cream (standardized as a 50 g tube) provides localized relief for diabetic neuropathic pain without systemic drug interactions or sedating CNS effects. Programmatic validation corrected: recommended\_dosage updated to clear units '50 g' rather than percentage. & Debatable &  \\
\hline
scot\_09\_11 & 1 & kept | Carbocisteine - 375mg capsules: two capsules three times a day @ Chronic Obstructive Pulmonary Disease &  & Acceptable &  \\
\hline
scot\_09\_11 & 2 & kept | Glyceryl Trinitrate - 400micrograms/dose pump sublingual spray: one to two puffs under the tongue if required for chest pain @ Ischaemic Heart Disease, Angina Pectoris &  & Acceptable &  \\
\hline
scot\_09\_11 & 3 & kept | Insulin Lispro - 100units/ml solution for injection 3ml pre-filled pen: to be injected twice a day @ Type 2 Diabetes Mellitus &  & Acceptable &  \\
\hline
scot\_09\_11 & 4 & kept | Insulin Glargine - 300units/ml solution for injection 1.5ml pre-filled pen: to be injected as directed by diabetes specialist nurse @ Type 2 Diabetes Mellitus &  & Acceptable &  \\
\hline
scot\_09\_11 & 5 & kept | Levothyroxine Sodium - 100microgram tablets: one tablet once daily @ Hypothyroidism &  & Correct &  \\
\hline
scot\_09\_11 & 6 & kept | Liraglutide - 6mg/ml solution for injection 3ml pre-filled: 1.2mg once daily @ Type 2 Diabetes Mellitus, Ischaemic Heart Disease &  & Acceptable &  \\
\hline
scot\_09\_11 & 7 & kept | Salbutamol - 100micrograms/dose inhaler CFC free: one or two puffs up to four times daily as required @ Chronic Obstructive Pulmonary Disease &  & Correct &  \\
\hline
scot\_09\_11 & 8 & kept | Salbutamol - 2.5mg/2.5ml nebuliser liquid unit dose vials: use one nebule when required as directed by respiratory clinic @ Chronic Obstructive Pulmonary Disease &  & Debatable &  \\
\hline
scot\_09\_11 & 9 & kept | Trelegy Ellipta - ® 92mcg/55mcg/222mcg per dose device: one inhalation once daily @ Chronic Obstructive Pulmonary Disease &  & Incorrect & melhor seria substituir por um Respimat. Paciente em uso de O2 e com VEF insuficiente. \\
\hline
scot\_09\_11 & 10 & discontinue | Ascorbic Acid - 500mg tablets: two tablets morning and night @ No health problem indication → (discontinued) & Ascorbic acid 500mg tablets (prescribed as two tablets twice daily, total 2000mg/day) has no documented clinical diagnosis or valid medical indication. Although ascorbic acid can enhance iron absorption, routine high-dose supplementation is unnecessary and adds significant pill burden in a severely frail patient with polypharmacy. Discontinuing ascorbic acid eliminates an unnecessary drug therapy without compromising care. & Correct &  \\
\hline
scot\_09\_11 & 11 & discontinue | Ferrous Fumarate - 210mg tablets: one tablet twice a day @ No health problem indication → (discontinued) & Ferrous fumarate 210mg lacks a documented indication; FBC is normal with no diagnosed iron deficiency anemia. Discontinuing it eliminates pill burden, reduces GI adverse effects, and resolves a clinically significant binding interaction with levothyroxine. & Correct &  \\
\hline
scot\_09\_11 & 12 & discontinue | Folic Acid - 5mg tablets: one tablet in the morning @ No health problem indication → (discontinued) & Folic acid 5mg daily is prescribed without a documented indication or co-prescribed antifolate therapy. Recent blood tests reveal an elevated serum folate level (>20 ng/mL) and normal full blood count, with no evidence of ongoing folate deficiency or anemia (ferrous fumarate has already been discontinued). Continuing high-dose folic acid represents unnecessary drug therapy with potential risk of masking vitamin B12 deficiency. Therefore, folic acid should be safely discontinued. & Correct &  \\
\hline
scot\_09\_11 & 13 & discontinue | Lactulose - 3.1-3.7g/5ml oral solution: 10ml twice a day @ No health problem indication → (discontinued) & Lactulose is prescribed without an active documented indication such as chronic constipation or hepatic encephalopathy. While the patient receives oxycodone, there is no charted diagnosis or active symptom of opioid-induced constipation justifying continuation. Additionally, concurrent use with high-dose furosemide increases the risk of dehydration and electrolyte disturbances (e.g., hypokalemia). Discontinuing lactulose eliminates unnecessary medication burden and associated adverse risks. & Correct &  \\
\hline
scot\_09\_11 & 14 & discontinue | Oxycodone - 5mg modified-release tablets: one tablet twice a day @ Malignant Neoplasm Of Cervix → (discontinued) & Oxycodone modified-release was initiated historically for cervical cancer diagnosed 8 months ago, but there is no documented active cancer-related pain or current indication justifying ongoing round-the-clock opioid therapy. In this severely frail 71-year-old patient (Rockwood score 8) with severe COPD on LTOT, concurrent use of oxycodone with multiple CNS depressants (pregabalin, mirtazapine, venlafaxine) creates an unacceptable risk of sedation, fall-related injury, and life-threatening respiratory depression. Discontinuing oxycodone eliminates unnecessary opioid exposure safely. & Correct &  \\
\hline
scot\_09\_11 & 15 & discontinue | Pregabalin - 25mg capsules: one capsule twice a day @ No health problem indication → (discontinued) & Pregabalin 25mg is prescribed without a documented medical indication (such as neuropathic pain, epilepsy, or generalized anxiety disorder). In a severely frail patient (Rockwood Clinical Frailty Scale 8) with ongoing fall prevention measures, pregabalin adds unnecessary risk of CNS depression, sedation, and falls, especially alongside mirtazapine. Discontinuing pregabalin removes unnecessary drug therapy and mitigates fall and sedative risks. & Correct &  \\
\hline
scot\_09\_11 & 16 & discontinue | Solifenacin - 5mg tablets: one tablet in the morning @ No health problem indication → (discontinued) & Solifenacin (solifenacin\_5mg) has no documented indication (such as overactive bladder) in the medical record. In this severely frail 71-year-old patient (Rockwood CFS 8) with fall prevention needs, solifenacin introduces unnecessary anticholinergic burden and interacts with umeclidinium in Trelegy Ellipta and mirtazapine, worsening fall and cognitive risks, and carries additive QT risk with venlafaxine. Discontinuing solifenacin resolves the drug therapy problem without requiring replacement. & Correct &  \\
\hline
scot\_09\_11 & 17 & add | (new) → Empagliflozin - 10 mg once daily @ Type 2 Diabetes Mellitus & Patient has type 2 diabetes mellitus with established ASCVD (history of MI and CABG) and suboptimal glycemic control (HbA1c 7.4\% [57 mmol/mol]). Current therapy includes metformin, liraglutide, and insulin. Per ADA/EASD guidelines, adding an SGLT2 inhibitor (empagliflozin) provides proven cardiovascular and renal protection with a distinct, complementary mechanism of action without significantly increasing hypoglycemia risk. Empagliflozin 10 mg daily is appropriate for eGFR 60 mL/min/1.73m2. & Debatable &  \\
\hline
scot\_09\_11 & 18 & add | (new) → Atorvastatin - 20 mg once daily @ Ischaemic Heart Disease & Patient has established ischaemic heart disease (prior MI and CABG) and elevated cholesterol (6.3 mmol/L) but lacks foundational lipid-lowering therapy. Guidelines mandate high-intensity statin therapy for secondary cardiovascular prevention. Atorvastatin coexists with current GTN and liraglutide, reducing recurrent cardiovascular events and mortality. A moderate starting dose of 20 mg once daily is prudent given severe frailty (CFS 8) and polypharmacy, balancing plaque stabilization against myopathy risk. LFTs are normal, supporting safe initiation. & Acceptable &  \\
\hline
scot\_09\_11 & 19 & add | (new) → Amlodipine - 5 mg once daily @ Angina Pectoris & Patient has established angina pectoris, ischemic heart disease, and prior MI/CABG, managed solely with PRN sublingual glyceryl trinitrate without regular anti-anginal prophylaxis. Given severe COPD on LTOT and bronchodilators where beta-blockers pose bronchospasm risks, initiating a cardioselective beta-blocker or a dihydropyridine calcium channel blocker is indicated. Amlodipine 5 mg once daily provides effective first-line anti-anginal prophylaxis to reduce ischemic episodes and nitrate reliance without bronchospastic risk. & Debatable &  \\
\hline
scot\_09\_11 & 20 & replace | Doxazosin - 4mg tablet: one tablet daily @ Essential Hypertension → Ramipril - 2.5 mg once daily @ Essential Hypertension & Doxazosin poses significant orthostatic hypotension and syncope risks in this severely frail patient (CFS 8) receiving concurrent vasodilating agents. Replacing doxazosin with ramipril provides safer, guideline-directed first-line blood pressure control for Essential Hypertension with renal and cardiovascular protection post-MI, while significantly mitigating orthostatic fall hazards. & Correct &  \\
\hline
scot\_09\_11 & 21 & replace | Furosemide - 40mg tablets: two tablets twice a day @ Essential Hypertension → Indapamide - 2.5 mg once daily @ Essential Hypertension & Furosemide 80 mg daily is inappropriate for primary essential hypertension in the absence of heart failure or severe renal impairment, conferring high risks of hypokalemia (exacerbated by inhaled beta-2 agonists) and volume depletion in a severely frail patient (Rockwood 8). Replacing furosemide with indapamide preserves guideline-recommended antihypertensive diuresis with a significantly safer profile and lower volume depletion risk. & Debatable & Tirar furosemida, não acrescentar indapamida. \\
\hline
scot\_09\_11 & 22 & replace | Mirtazapine - 45mg tablets: one tablet at night @ Depressive Episode → Bupropion - 150 mg once daily @ Depressive Episode & Mirtazapine promotes substantial weight gain, hyperlipidemia, and glycemic dysregulation, which is hazardous in severe obesity (BMI 40.1 kg/m2), hypercholesterolemia, and type 2 diabetes. Replacing mirtazapine with bupropion provides effective antidepressant therapy with a weight-neutral to weight-loss promoting metabolic profile, reducing cardiometabolic risk while maintaining antidepressant coverage. Note: venlafaxine remains on chart, but bupropion can be utilized safely with monitoring. & Debatable & Suspender mirta, sem necessidade de bup \\
\hline
scot\_09\_11 & 23 & replace | Venlafaxine - 75mg tablets: two tablets in the morning and one at night @ Depressive Episode → Sertraline - 50 mg once daily @ Depressive Episode & Venlafaxine 225mg daily exerts significant noradrenergic effects that elevate blood pressure and cardiac workload in this severely frail patient with ischaemic heart disease and hypertension. Replacing venlafaxine with sertraline maintains effective treatment for Depressive Episode while offering a proven, cardioprotective and well-tolerated safety profile in post-myocardial infarction patients, avoiding adrenergic excess. & Correct &  \\
\hline
scot\_09\_11 & 24 & reduce\_dosage | Metformin - 500mg tablets: two tabs twice daily before breakfast and evening meal @ Type 2 Diabetes Mellitus → Metformin - 500 mg twice daily @ Type 2 Diabetes Mellitus & Metformin is currently prescribed at 1000 mg twice daily (2000 mg/day) for type 2 diabetes. With a calculated creatinine clearance of 42 mL/min (moderate renal impairment), the maximum recommended daily dose of metformin is 1000 mg to prevent drug accumulation and reduce the risk of metformin-associated lactic acidosis. Reducing the dosage to 500 mg twice daily maintains glycemic efficacy alongside existing therapies while ensuring safety under impaired renal clearance. & Acceptable & melhor suspender metformina, principalmente pq o algoritmo sugeriu incluir empaglifozina. \\
\hline
scot\_09\_12 & 1 & kept | Apixaban - 5mg tablets: one tablet twice a day @ Stroke Prevention In Atrial Flutter, Cardiovascular Prophylaxis &  & Correct &  \\
\hline
scot\_09\_12 & 2 & kept | Bisoprolol - 10mg tablets: one tablet each morning @ Mild Left Ventricular Dysfunction, Essential Hypertension, Cardiovascular Prophylaxis &  & Correct &  \\
\hline
scot\_09\_12 & 3 & kept | Bumetanide - 1mg tablets: two tablets morning and lunchtime @ Mild Left Ventricular Dysfunction &  & Acceptable &  \\
\hline
scot\_09\_12 & 4 & kept | Salbutamol - 100microgram metered dose inhaler: two puffs when required @ Chronic Obstructive Pulmonary Disease &  & Correct &  \\
\hline
scot\_09\_12 & 5 & kept | Senna - 7.5mg tablets: two tablets at night @ Constipation &  & Correct &  \\
\hline
scot\_09\_12 & 6 & kept | Spironolactone - 25mg tablets: one tablet in the morning @ Mild Left Ventricular Dysfunction, Essential Hypertension &  & Correct &  \\
\hline
scot\_09\_12 & 7 & kept | Trelegy - ® dry powder inhaler 92/55/22micrograms/dose: one inhalation daily @ Chronic Obstructive Pulmonary Disease &  & Correct &  \\
\hline
scot\_09\_12 & 8 & discontinue | Amiodarone - 200mg tablets: one tablet daily @ No health problem indication → (discontinued) & Amiodarone 200mg is currently prescribed without a valid, documented medical indication on the MAR. Discontinuing amiodarone eliminates the risks of serious multiorgan toxicities (pulmonary fibrosis, thyroid dysfunction, hepatotoxicity) and critical drug-drug interactions, including CYP3A4 and P-gp inhibition with apixaban and oxycodone, additive bradycardia risk with bisoprolol, and QT prolongation with solifenacin and trelegy. & Incorrect &  \\
\hline
scot\_09\_12 & 9 & discontinue | Aspirin - 75mg enteric coated tablets: one tablet daily @ Cardiovascular Prophylaxis, Stroke Prevention In Atrial Flutter → (discontinued) & Aspirin is redundant for stroke prevention in atrial flutter and cardiovascular prophylaxis because the patient is already receiving therapeutic anticoagulation with apixaban. Concomitant use of aspirin and apixaban substantially elevates major bleeding risk without clear indication (no recent ACS or stenting) and exacerbates iron deficiency anaemia. Stopping aspirin resolves this redundant antithrombotic therapy and mitigates bleeding and anaemia risks. & Correct &  \\
\hline
scot\_09\_12 & 10 & add | (new) → Dapagliflozin - 10 mg once daily @ Mild Left Ventricular Dysfunction & Patient has mild left ventricular dysfunction with prior documented adverse reactions to ACE inhibitors and ARBs, and is managed with bisoprolol, bumetanide, and spironolactone. Guideline-directed heart failure therapy recommends adding an SGLT2 inhibitor (dapagliflozin 10 mg once daily) as a distinct-role adjunct to reduce the risk of heart failure hospitalization and cardiovascular progression, which also aligns with benefits for concomitant type 2 diabetes mellitus and CKD. Renal function (eGFR >60 mL/min/1.73m2) is safe for initiation. & Incorrect &  \\
\hline
scot\_09\_12 & 11 & add | (new) → Macrogol - 13.125 g once daily @ Constipation & Persistent constipation despite senna is driven by high-dose opioids, solifenacin, and ferrous fumarate. Guidelines recommend adjunct escalation by co-prescribing an osmotic laxative (macrogol 13.125 g) alongside a stimulant laxative. Dosage formatted to 13.125 g to resolve validation issue regarding '1 sachet'. & Debatable &  \\
\hline
scot\_09\_12 & 12 & add | (new) → Atorvastatin - 20 mg once daily @ Cardiovascular Prophylaxis & Patient has high cardiovascular risk due to type 2 diabetes (>10 years duration), age >40, chronic kidney disease stage 2, severe obesity, and smoking history, but currently lacks lipid-lowering primary prevention. In line with NICE and ESC cardiovascular disease prevention guidelines, atorvastatin 20 mg once daily is indicated for primary cardiovascular prophylaxis to lower LDL-C and reduce major adverse cardiovascular events. Baseline LFTs and renal function are preserved, ensuring safe initiation. & Correct &  \\
\hline
scot\_09\_12 & 13 & replace | Ferrous Fumarate - 210mg tablets: one tablet three times a day @ Iron Deficiency Anaemia → Ferrous Gluconate - 300 mg once daily @ Iron Deficiency Anaemia & Ferrous fumarate is directly causing and aggravating severe constipation requiring multiple laxatives. Ferrous gluconate contains less elemental iron per dose, offers better gastrointestinal tolerability, and preserves oral iron replacement for iron deficiency anaemia with a reduced burden of constipation. & Incorrect &  \\
\hline
scot\_09\_12 & 14 & replace | Gabapentin - 100mg capsules: one capsule three times a day @ Back Pain → Paracetamol - 1000 mg four times daily @ Back Pain & Gabapentin interacts pharmacodynamically with oxycodone (Oxypro and Shortec), causing additive CNS and respiratory depression, particularly dangerous given COPD and OSA. Replacing gabapentin with paracetamol provides safe analgesic coverage for back pain without interacting or worsening sedation or respiratory depression. Topical or non-sedating non-opioid analgesics are preferred alternatives. & Correct &  \\
\hline
scot\_09\_12 & 15 & replace | Shortec - ® 10mg capsules: one capsule when required for pain (maximum twice a day) @ Back Pain → Lidocaine - 700 mg once daily @ Back Pain & Shortec (oxycodone) risks respiratory depression and exacerbates OSA/constipation alongside Oxypro. Replacing PRN Shortec with topical lidocaine (700 mg medicated plaster) avoids dosage format validation errors ('5 \%' corrected to '700 mg') and delivers targeted non-systemic relief for back pain without systemic opioid toxicity. Topical therapy minimizes systemic risks given concurrent CKD and apixaban preclude oral NSAIDs. & Debatable &  \\
\hline
scot\_09\_12 & 16 & replace | Solifenacin - 5mg tablets: one tablet daily @ Urinary Incontinence → Mirabegron - 50 mg once daily @ Urinary Incontinence & Solifenacin contributes to active constipation and increases urinary retention risk due to additive anticholinergic burden with oxycodone (Oxypro) and umeclidinium (Trelegy). Switching to mirabegron, a beta-3 adrenergic agonist, preserves therapeutic efficacy for urinary incontinence without antimuscarinic side effects, mitigating gastrointestinal and retention risks. Blood pressure is well-controlled (130/70 mmHg), making mirabegron a safer alternative. & Incorrect &  \\
\hline
scot\_09\_12 & 17 & reduce\_dosage | Oxypro - ® 40mg modified release tablets: one tablet morning and night @ Back Pain → Oxypro - 20 mg twice daily @ Back Pain & Oxypro 40 mg twice daily (80 mg/day oxycodone, \textasciitilde{}120 MME/day) exceeds recommended thresholds for chronic non-cancer back pain (>50 MME/day), escalating risks of respiratory depression, worsening obstructive sleep apnoea, and severe constipation. Tapering to 20 mg twice daily reduces opioid burden towards safer levels while non-opioid multimodal analgesia (paracetamol, topical lidocaine) continues. Dose de-escalation mitigates toxicity and respiratory compromise while maintaining pain control. & Correct &  \\
\hline
scot\_09\_13 & 1 & kept | Co-Beneldopa - 12.5mg/50mg dispersible tablet: one tablet daily (bedtime) @ Parkinson'S Disease &  & Correct &  \\
\hline
scot\_09\_13 & 2 & kept | Colecalciferol - 400 unit tablet: one daily @ Vitamin D Deficiency Prophylaxis &  & Acceptable &  \\
\hline
scot\_09\_13 & 3 & kept | Laxido - ® sachets: one sachet when required @ No health problem indication &  & Acceptable &  \\
\hline
scot\_09\_13 & 4 & kept | Pregabalin - 25mg capsules: one capsule twice daily @ Chronic Pain &  & Debatable &  \\
\hline
scot\_09\_13 & 5 & kept | Sertraline - 100mg tablets: one tablet daily @ Depression &  & Debatable &  \\
\hline
scot\_09\_13 & 6 & discontinue | Peptac - ® suspension: 10ml four times daily when required @ No health problem indication → (discontinued) & Peptac suspension 10ml PRN is prescribed without a documented clinical indication, diagnosis, or symptom of gastro-oesophageal reflux disease or dyspepsia. In a moderately frail patient with Parkinson's disease and swallowing difficulties, discontinuing this unindicated medication reduces pill burden, avoids unnecessary drug exposure, and lowers the risk of adverse effects. & Correct &  \\
\hline
scot\_09\_13 & 7 & replace | Co-Beneldopa - 25mg/100mg tablets: one tablet four times daily @ Parkinson'S Disease → Co-Beneldopa dispersible - 1 tablet four times daily @ Parkinson'S Disease & Patient has dysphagia, making standard solid tablets of Co-Beneldopa 25mg/100mg hazardous due to choking and erratic absorption. Per validation fix for dosage formatting (amount + unit only), replacing with Co-Beneldopa dispersible tablets (benserazide/levodopa 25mg/100mg equivalent) ensures safe swallowing and uninterrupted Parkinson's disease symptom management. Co-Careldopa was considered but Co-Beneldopa dispersible formulation directly maintains the same active ingredients. & Correct &  \\
\hline
scot\_09\_13 & 8 & replace | Co-Codamol - 30mg/500 tablets: two tablets four times daily @ Chronic Pain → Paracetamol - 1000 mg four times daily @ Chronic Pain & Co-codamol (30/500mg QDS) carries high risks of sedation, constipation, and falls, worsened by interactions with pregabalin and sertraline. In a 78-year-old patient with Parkinson's disease, moderate frailty, and dysphagia, regular opioids present an unfavorable risk-benefit profile. Replacing co-codamol with plain paracetamol maintains basic analgesia for chronic pain without opioid-mediated sedation or respiratory depression. & Correct &  \\
\hline
scot\_09\_14 & 1 & kept | Clenil - ® (beclomethasone dipropionate) 100microgram meter dose inhaler: two puffs twice a day @ Asthma &  & Correct &  \\
\hline
scot\_09\_14 & 2 & kept | Salbutamol - inhaler 100mcg meter dose inhaler: two puffs up to four times a day (two inhalers ordered in last 12 months) @ Asthma &  & Correct &  \\
\hline
scot\_09\_14 & 3 & replace | Citalopram - 20mg tablets: one tablet daily @ Depressive Episode, Severe Depression → Venlafaxine - 37.5 mg once daily @ Depressive Episode, Severe Depression & Patient has documented citalopram failure with persistent severe depressive symptoms, PHQ-9 of 14, and suicidal ideation, indicating refractory depression. Replacing citalopram with an alternative first-line antidepressant of a different class, such as the SNRI venlafaxine or alternative SSRI/SNRI, provides renewed efficacy for severe depression. Initiating venlafaxine 37.5 mg daily allows titration to an effective therapeutic dose while monitoring for clinical response and tolerability. & Correct &  \\
\hline
scot\_09\_16 & 1 & replace | Solpadeine Max - ® (codeine 12.8mg and paracetamol 500mg) @ Pain From Assault → Paracetamol - 1000 mg four times daily prn @ Pain From Assault & Solpadeine Max contains codeine, an opioid agonist posing high risk of dependence and overdose in this patient with active polydrug use, psychiatric distress, and fleeting suicidal ideation. Replacing it with plain paracetamol maintains effective non-opioid analgesia for pain from assault while removing opioid exposure and avoiding compounding substance misuse risks. Paracetamol provides safe, first-line baseline pain relief without abuse potential. & Correct &  \\
\hline
scot\_09\_17 & 1 & kept | Quinine Sulfate - 200mg tablets: one tablet at night @ Nocturnal Leg Cramps &  & Incorrect &  \\
\hline
scot\_09\_17 & 2 & replace | Atorvastatin - 20mg tablets: one tablet daily @ Primary Hypercholesterolaemia → Pravastatin - 20 mg once daily @ Primary Hypercholesterolaemia & Atorvastatin undergoes CYP3A4 metabolism and interacts with quinine sulfate, compounding myopathy risks and nocturnal leg cramps. Replacing atorvastatin with pravastatin, which is not metabolised by CYP3A4 and possesses minimal myotoxicity and drug-interaction potential, breaks this interaction while maintaining lipid control for primary hypercholesterolaemia. & Debatable &  \\
\hline
scot\_09\_17 & 3 & replace | Ramipril - 10mg capsules: one capsule daily (increased two months prior to review) @ Hypertension → Amlodipine - 5 mg once daily @ Hypertension & Ramipril reduces renal clearance of lithium, leading to increased serum lithium concentrations and resultant toxicity. Replacing ramipril with amlodipine maintains effective blood pressure control without significantly altering renal lithium clearance, safely eliminating the adverse pharmacokinetic interaction. & Correct &  \\
\hline
scot\_09\_17 & 4 & reduce\_dosage | Lithium - (Priadel®) 400mg tablets: two tablets at night @ Bipolar Type 1 Disorder → Lithium - 600 mg once daily @ Bipolar Type 1 Disorder & Lithium dose of 800 mg daily has resulted in an elevated serum level of 0.98 mmol/L along with drowsiness and fine tremor in this 67-year-old patient. In older adults on maintenance therapy for bipolar type 1 disorder, target serum lithium levels are typically 0.4-0.6 mmol/L (or 0.6-0.8 mmol/L generally). Reducing Priadel to 600 mg once daily at night will lower serum concentration toward target range, resolving tremor and drowsiness while maintaining mood stabilization. Repeat serum lithium levels in 5-7 days. & Correct &  \\
\hline
scot\_09\_18 & 1 & kept | Atorvastatin - 80mg tablets: one tablet daily @ Cerebrovascular Disease, Mixed Dementia, Secondary Stroke Prevention, Cardiovascular Prophylaxis &  & Incorrect &  \\
\hline
scot\_09\_18 & 2 & kept | Perindopril - 4mg tablets: one tablet daily @ Cerebrovascular Disease, Mixed Dementia, Secondary Stroke Prevention, Cardiovascular Prophylaxis &  & Incorrect & Correto seria suspender. controle rígido de pressão \\
\hline
scot\_09\_18 & 3 & discontinue | Solifenacin - 5mg tablets: one tablet daily @ No health problem indication → (discontinued) & Solifenacin has no documented clinical indication (such as overactive bladder) in the medical record. Moreover, its strong anticholinergic properties oppose donepezil, worsen cognitive impairment and confusion in mixed dementia, contribute to constipation (no bowel movement for 3 days), and prolong QTc (470ms). Discontinuing solifenacin removes unnecessary drug therapy and mitigates severe antimuscarinic toxicity. & Correct &  \\
\hline
scot\_09\_18 & 4 & add | (new) → Memantine - 5 mg once daily @ Mixed Dementia & Patient has worsening cognitive and functional impairment and severe behavioral symptoms (AMT 4) with mixed dementia despite donepezil 5mg. Per NICE guidelines, memantine is indicated as adjunct therapy to an acetylcholinesterase inhibitor for moderate-to-severe dementia to slow progression and improve neuropsychiatric and cognitive function. Memantine acts via NMDA receptor antagonism, providing synergy without worsening bradycardia or cholinergic side effects. Initiate at 5 mg once daily and titrate weekly by 5 mg up to the maintenance dose of 20 mg once daily. & Acceptable &  \\
\hline
scot\_09\_18 & 5 & replace | Aspirin - 75mg dispersible tablets: one tablet daily @ Cerebrovascular Disease, Mixed Dementia, Secondary Stroke Prevention, Cardiovascular Prophylaxis → Clopidogrel - 75 mg once daily @ Cerebrovascular Disease, Mixed Dementia, Secondary Stroke Prevention, Cardiovascular Prophylaxis & For secondary stroke prevention and cerebrovascular disease following an ischemic stroke (PACS), clopidogrel 75 mg once daily monotherapy is guideline-recommended as more effective than aspirin monotherapy in reducing recurrent vascular events. Replacing aspirin with clopidogrel maintains antiplatelet efficacy while optimizing secondary prevention for cerebrovascular disease and mixed dementia without increasing GI toxicity risk. & Incorrect &  \\
\hline
scot\_09\_18 & 6 & replace | Hydroxyzine - 25mg tablets: one tablet at night @ Itchy Rash → Cetirizine - 10 mg once daily @ Itchy Rash & Hydroxyzine has potent anticholinergic and sedative properties that worsen cognitive impairment (AMT 4, confusion), exacerbate constipation, antagonize donepezil, and add to QTc prolongation risks with risperidone. Replacing hydroxyzine with cetirizine, a second-generation non-sedating antihistamine devoid of significant anticholinergic activity, effectively controls the itchy rash while eliminating these risks. & Incorrect &  \\
\hline
scot\_09\_18 & 7 & replace | Risperidone - 250 microgram tablets: one tablet twice daily @ Mixed Dementia → Trazodone - 25 mg once daily @ Mixed Dementia & Risperidone carries a black box warning for cerebrovascular adverse events and mortality in dementia, presenting high risk in this patient with prior stroke and borderline QTc. Trazodone is a non-antipsychotic option considered to safely manage severe agitation and behavioral distress without the cerebrovascular burden of antipsychotics. Memantine has also been added to help stabilize cognitive and behavioral symptoms. Monitor response and tolerability closely. & Debatable &  \\
\hline
scot\_09\_18 & 8 & reduce\_dosage | Donepezil - 5mg tablets: one tablet daily @ Mixed Dementia → Donepezil - 5 mg every other day @ Mixed Dementia & Donepezil 5mg daily is associated with clinically significant bradycardia (heart rate 50 bpm) via vagotonic effects on cardiac conduction in mixed dementia. Because 5mg is the standard starting dose and the patient has symptomatic or marked bradycardia, lengthening the interval to 5mg every other day de-escalates cholinergic exposure, mitigating bradycardia risk while continuing cognitive therapy alongside memantine. & Debatable &  \\
\hline
scot\_09\_19 & 1 & kept | Bendroflumethiazide - 2.5mg tablets: one tablet daily @ No health problem indication &  & Incorrect &  \\
\hline
scot\_09\_19 & 2 & kept | Metformin - 500mg tablets: one tablet twice a day @ Type 2 Diabetes Mellitus &  & Acceptable & é possível aumentar a dose \\
\hline
scot\_09\_19 & 3 & discontinue | Amitriptyline - 50mg tablets: one tablet at night @ No health problem indication → (discontinued) & Amitriptyline (amitriptyline\_50mg) has no documented indication in the chart and adds significant anticholinergic burden, CNS depression, and delirium risk in this 78-year-old patient who recently suffered delirium and is presenting with confusion and sleepiness. Discontinuing amitriptyline directly addresses this unnecessary therapy and reduces active sedation and cognitive slowing. & Correct &  \\
\hline
scot\_09\_19 & 4 & discontinue | Oxybutynin - 5mg tablets: one tablet twice a day @ No health problem indication → (discontinued) & Oxybutynin lacks a documented indication as the patient is explicitly continent with no history of overactive bladder. Its potent anticholinergic properties are directly contributing to her cognitive slowing, confusion, and drowsiness, exacerbating frailty and recent post-delirium vulnerability. Discontinuing oxybutynin resolves this unnecessary drug therapy problem and eliminates additive anticholinergic exposure. & Correct &  \\
\hline
scot\_09\_19 & 5 & discontinue | Risperidone - 500 microgram tablets: one tablet twice a day @ Short Admission (Three Days) With Delirium From A Suspected Chest Or Urine Infection → (discontinued) & Risperidone 500 mcg twice daily was initiated for acute delirium during a short admission for suspected infection. The acute episode has resolved, leaving no ongoing indication for antipsychotic therapy. Continued use contributes to documented excessive sedation, cognitive slowing, functional decline, and risks metabolic deterioration and orthostasis in this 78-year-old patient. Discontinuation safely addresses the lack of valid medical indication. & Correct &  \\
\hline
scot\_09\_19 & 6 & add | (new) → Linagliptin - 5 mg once daily @ Type 2 Diabetes Mellitus & Patient has uncontrolled Type 2 Diabetes Mellitus with HbA1c 58 mmol/mol (7.5\%) on baseline metformin 500mg twice daily. Per NICE guidelines, when monotherapy is insufficient, dual therapy with an adjunct agent is indicated. Linagliptin, a DPP-4 inhibitor, is chosen as an effective add-on with minimal hypoglycemia risk, neutral effect on body weight, and safety in mild frailty without requiring renal dose adjustment. & Acceptable &  \\
\hline
scot\_09\_20 & 1 & kept | Donepezil - 5mg tablets: one tablet in the morning @ Alzheimer’S Dementia &  & Correct &  \\
\hline
scot\_09\_20 & 2 & kept | Salbutamol - 200 microgram dry powder inhaler: one inhalation when required for shortness of breath/ wheeze @ Copd &  & Debatable &  \\
\hline
scot\_09\_20 & 3 & discontinue | Atorvastatin - 80mg tablets: one tablet in the morning @ No health problem indication → (discontinued) & Atorvastatin 80mg is prescribed without any documented indication such as ASCVD or familial hypercholesterolemia. In an 85-year-old frail female patient (CFS 5, weight 47kg, eGFR 25 mL/min) presenting with acute infection and delirium, high-intensity statin therapy carries unnecessary risks of adverse drug events without clear benefit. Discontinuing atorvastatin\_80mg resolves this unnecessary drug therapy. & Correct &  \\
\hline
scot\_09\_20 & 4 & discontinue | Lisinopril - 5mg tablets: one tablet in the morning @ Heart Failure With Preserved Ejection Fraction → (discontinued) & Lisinopril is held due to acute safety concerns in the setting of acute infection, poor oral intake, borderline low BP (108/56 mmHg), hyperkalemia (potassium 5.2 mmol/L), and renal impairment (CrCl 25 mL/min). Continuing an ACE inhibitor during acute hypovolemia and systemic illness heightens the risk of worsening acute kidney injury. Long-term therapy remains indicated for HFpEF once the patient stabilizes. Discontinuing temporarily mitigates renal and hemodynamic risks while maintaining cardiovascular safety. & Correct &  \\
\hline
scot\_09\_20 & 5 & discontinue | Macrogol 3350 - sachets: one sachet twice daily @ No health problem indication → (discontinued) & Macrogol 3350 is prescribed without an active or documented indication of constipation in the patient's medical record. In an 85-year-old patient with heart failure with preserved ejection fraction (HFpEF) and severe renal impairment (CrCl 25 mL/min), unnecessary laxative therapy introduces polypharmacy burden and risk of fluid/electrolyte disturbances. Discontinuing macrogol\_3350\_sachets\_one\_ resolves the unnecessary drug therapy problem. & Correct &  \\
\hline
scot\_09\_20 & 6 & discontinue | Senna - 7.5mg tablets: two tablets at night @ No health problem indication → (discontinued) & Senna (senna\_7p5mg) has no active documented indication or symptom of constipation in the chart. Continuing unnecessary stimulant laxative therapy exposes the patient to avoidable risks of fluid/electrolyte disturbances, especially with concurrent furosemide use. Senna should be discontinued. & Correct &  \\
\hline
scot\_09\_20 & 7 & add | (new) → Amoxicillin - 500 mg every 12 hours @ Upper Urinary Tract Infection & Patient has an active upper urinary tract infection (E. coli, CRP 72 mg/L, fever, back pain) managed only supportively with paracetamol. The isolate is sensitive to amoxicillin, gentamicin, and aztreonam. The reported penicillin allergy was non-allergic gastrointestinal intolerance (nausea/diarrhoea >20 years ago), making amoxicillin safe and indicated. Given CrCl 25 mL/min, amoxicillin dose-adjusted to 500 mg every 12 hours provides effective targeted oral antimicrobial eradication, preventing worsening sepsis while avoiding nephrotoxic alternatives like gentamicin. & Correct &  \\
\hline
scot\_09\_20 & 8 & add | (new) → Memantine - 5 mg once daily @ Alzheimer’S Dementia & Patient has Alzheimer's dementia with suboptimal control on donepezil monotherapy, presenting with cognitive and behavioral decline (MMSE 23/30). NICE and international guidelines recommend memantine as an adjunct to donepezil for moderate-to-severe disease or progressive impairment. With CrCl 25 mL/min, initiating memantine at a renal-adjusted starting dose of 5 mg once daily provides dual-mechanism NMDA-receptor neuroprotection alongside acetylcholinesterase inhibition without displacing donepezil, aiming to stabilize cognition and functional capacity. & Debatable &  \\
\hline
scot\_09\_20 & 9 & increase\_dosage | Furosemide - 20mg tablets: one tablet in the morning @ Heart Failure With Preserved Ejection Fraction → Furosemide - 40 mg once daily @ Heart Failure With Preserved Ejection Fraction & Patient has persistent fluid retention with marked pitting oedema to both knees, ankle swelling, and worsening breathlessness in the setting of HFpEF and renal impairment (CrCl 25 mL/min). The current dose of Furosemide 20mg once daily is inadequate to achieve effective diuresis given reduced renal clearance. Increasing Furosemide to 40mg once daily is recommended to promote decongestion and relieve symptoms. Monitor renal function, blood pressure, and electrolytes closely. & Incorrect &  \\
\hline
scot\_09\_20 & 10 & replace | Trelegy Ellipta - ® (fluticasone 92 micrograms/ vilanterol 22 micrograms/ umeclidinium 65 micrograms) dry powder inhaler, one inhalation in the morning @ Copd → Relvar Ellipta - 92 mcg once daily @ Copd & Trelegy Ellipta contains umeclidinium, which exerts anticholinergic effects that antagonize donepezil and worsen delirium. Replacing it with Relvar Ellipta (fluticasone furoate/vilanterol 92/22 mcg) maintains ICS/LABA maintenance therapy for COPD while eliminating the anticholinergic component to break the interaction. Dosage specified as 92 mcg. Programmatic validation failed: recommended\_dosage: Dosage format unclear - missing units or frequency: '1 inhalation'. & Debatable &  \\
\hline
scot\_09\_20 & 11 & reduce\_dosage | Paracetamol - 500mg caplets: two caplets four times daily @ Upper Urinary Tract Infection → Paracetamol - 500 mg four times daily @ Upper Urinary Tract Infection & Paracetamol is currently prescribed at 1000 mg four times daily (4 g/day) for pain and fever associated with upper urinary tract infection. The patient weighs 47 kg (<50 kg), is 85 years old, frail, and acutely unwell with reduced oral intake. In low body weight (<50 kg) adults, paracetamol dosing should not exceed 60 mg/kg/day (approximately 2–3 g/day) to prevent hepatotoxicity. Reducing the dose to 500 mg four times daily (total 2 g daily) maintains effective analgesia and antipyresis while mitigating cumulative toxicity risks. & Correct &  \\
\hline
scot\_09\_21 & 1 & kept | Aspirin - 75mg dispersible tablets: one tablet daily @ Ischaemic Heart Disease &  & Incorrect & paciente em fase final de vida. Sem indicação. Como fazer o algoritmo reconhecer? \\
\hline
scot\_09\_21 & 2 & kept | Iv Pamidronate - for hypercalcaemia of malignancy @ Non-Small Cell Lung Cancer With Bone, Liver And Adrenal Metastases, Hypercalcaemia Of Malignancy &  & Incorrect &  \\
\hline
scot\_09\_21 & 3 & discontinue | Amlodipine - 5mg tablets: one tablet daily @ Hypertension, Ischaemic Heart Disease → (discontinued) & Amlodipine is indicated for hypertension and ischaemic heart disease, but the patient presents with acute hypotension (BP 100/55 mmHg, pulse 110 bpm), dehydration, and acute renal impairment. Continuing a vasodilator risks exacerbating hypoperfusion and renal injury. A temporary safety discontinuation is required until acute illness resolves and hemodynamics stabilize. & Correct &  \\
\hline
scot\_09\_21 & 4 & discontinue | Macrogol 3350 - sachets: one sachet twice daily @ No health problem indication → (discontinued) & Macrogol 3350 is prescribed without an active documented indication for constipation. The patient is currently experiencing incontinence of loose stools and dehydration with impaired renal function (creatinine clearance 21 mL/min, urea 22.0 mmol/L). Discontinuing macrogol\_3350\_sachets\_one\_ directly resolves the unnecessary laxative therapy, preventing further fluid loss, worsening dehydration, and fecal incontinence. & Correct &  \\
\hline
scot\_09\_21 & 5 & discontinue | Quinine Sulfate - 200mg tablets: one tablet at night @ No health problem indication → (discontinued) & Quinine sulfate has no documented clinical indication (e.g., malaria, leg cramps) in this severely frail patient with renal impairment and acute deterioration. Discontinuing it avoids unnecessary risks, including thrombocytopenia and CYP3A4 inhibition with charted simvastatin. & Correct &  \\
\hline
scot\_09\_21 & 6 & discontinue | Senna - 7.5mg tablets: two tablets at night @ No health problem indication → (discontinued) & Senna 7.5mg is charted without an active indication for constipation. The patient is severely frail, dehydrated (urea 22.0 mmol/L), and actively incontinent of loose stools, which Senna exacerbates. Discontinuing Senna directly addresses the drug therapy problem of unnecessary medication contributing to diarrhea and worsening dehydration. & Correct &  \\
\hline
scot\_09\_21 & 7 & discontinue | Thiamine - 100mg tablets: one tablet three times daily @ Acute Confusion → (discontinued) & Thiamine was historically started for alcohol excess, but the patient has been abstinent for 5 years. The current acute confusion is multifactorial delirium driven by community-acquired pneumonia, hypercalcaemia of malignancy, and dehydration, rather than active Wernicke encephalopathy or acute alcohol withdrawal. Additionally, the patient has difficulty swallowing. Continuing oral thiamine lacks a valid medical indication at this time and adds unnecessary pill burden. Deprescribing thiamine\_100mg resolves this unnecessary therapy without compromising delirium management. & Correct &  \\
\hline
scot\_09\_21 & 8 & add | (new) → Clarithromycin - 250 mg twice daily @ Presumed Pneumonia & The patient displays worsening severe pneumonia (tachypnea RR 26, tachycardia 110 bpm, coarse crackles, hypothermia 35.5°C, CRP 75 mg/L) despite 24 hours of intravenous antibiotics. For severe community-acquired pneumonia meeting refractory criteria, guidelines recommend adjunct macrolide therapy (clarithromycin) co-prescribed with beta-lactam IV antibiotics for atypical coverage and immunomodulatory benefit. Clarithromycin 500 mg is adjusted for renal clearance (CrCl 21 mL/min) to 250 mg twice daily. High risk due to renal impairment and frailty; monitor renal function and ECG. & Debatable &  \\
\hline
scot\_09\_21 & 9 & add | (new) → Sodium Chloride 0.9\% - 1000 ml every 8 hours @ Dehydration & Patient has severe dehydration, prerenal renal impairment (urea 22 mmol/L, Cr 180 umol/L), hypotension (100/55 mmHg), and hypercalcaemia of malignancy with difficulty swallowing. Intravenous fluid resuscitation with sodium chloride 0.9\% is guideline-directed to restore intravascular volume, improve renal perfusion, and facilitate calciuresis. Risk is moderate given frailty and crackles; rate must be tailored to hemodynamic status and fluid balance monitored closely. & Acceptable &  \\
\hline
scot\_09\_21 & 10 & add | (new) → Nystatin - 100,000 units four times daily @ Oral Candidiasis & Patient has severe oral candidiasis contributing to dysphagia and dehydration, with no antifungal therapy prescribed. Adding nystatin oral suspension provides localized topical antifungal coverage without systemic CYP3A4 interactions with charted simvastatin or renal toxicity given the patient's low CrCl (21 mL/min). Chosen over systemic fluconazole, which is not on the regimen, to avoid significant CYP-mediated drug interactions. Treating oral candidiasis will relieve oral discomfort and improve oral intake. & Correct &  \\
\hline
scot\_09\_21 & 11 & add | (new) → Denosumab - 120 mg every 4 weeks @ Non-Small Cell Lung Cancer With Bone, Liver And Adrenal Metastases & Adjusted calcium remains elevated at 3.21 mmol/L secondary to metastatic NSCLC despite IV pamidronate and IV fluids. The patient has renal impairment (CrCl 21 mL/min). Subcutaneous denosumab provides a distinct mechanism (RANKL inhibition) recommended for refractory hypercalcaemia of malignancy, avoiding bisphosphonate nephrotoxicity. Adding denosumab as adjunct therapy addresses severe hypercalcaemia alongside current supportive care. Given renal dysfunction and frailty, moderate risk is managed with calcium and phosphate monitoring. & Incorrect & paciente em fase final de vida. Não precisa \\
\hline
scot\_09\_21 & 12 & replace | Co-Beneldopa - 100/25mg capsules: one capsule every six hours @ Parkinson’S Disease → Rotigotine - 2 mg/24hr once daily @ Parkinson’S Disease & Co-Beneldopa oral capsules cannot be safely swallowed due to acute dysphagia and confusion, posing risks of aspiration pneumonia and missed doses. Rotigotine transdermal patch replaces oral dopaminergic therapy non-orally, ensuring continuous Parkinson's symptom management and avoiding neuroleptic malignant-like withdrawal while oral intake is impaired. & Acceptable &  \\
\hline
scot\_09\_21 & 13 & replace | Paracetamol - 500mg caplets: two caplets every six hours @ Presumed Pneumonia, Non-Small Cell Lung Cancer With Bone, Liver And Adrenal Metastases, Community Acquired Pneumonia → Paracetamol IV - 500 mg every six hours @ Presumed Pneumonia, Non-Small Cell Lung Cancer With Bone, Liver And Adrenal Metastases, Community Acquired Pneumonia & Paracetamol oral caplets are inappropriate due to acute dysphagia and aspiration risk. Switching to Paracetamol IV preserves analgesia and antipyresis. However, 1000 mg IV q6h is excessive given low body weight (52 kg), liver metastases, CKD (CrCl 21 mL/min), and severe frailty, risking toxicity; renal/hepatic and low body weight dosing guidelines recommend reducing paracetamol IV to 500 mg every 6 hours (max 2-3 g/day, <=15 mg/kg per dose, interval >=6 hours). & Acceptable &  \\
\hline
scot\_09\_21 & 14 & replace | Intravenous (Iv) Antibiotics @ Presumed Pneumonia, Community Acquired Pneumonia → Piperacillin/Tazobactam - 2.25 g every 8 hours @ Presumed Pneumonia, Community Acquired Pneumonia & Patient is deteriorating with severe pneumonia (RR 26, coarse crepitations, hypothermia, CRP 75 mg/L) despite 24 hours of unspecified IV antibiotics. Clarithromycin was added, but the primary broad-spectrum IV coverage requires escalation/replacement. Replacing the refractory IV antibiotics with IV Piperacillin/Tazobactam provides necessary antipseudomonal and broad-spectrum coverage for severe hospital/healthcare-associated aspiration or community-acquired pneumonia, adjusted for renal impairment (CrCl 21 mL/min). & Debatable &  \\
\hline
scot\_09\_21 & 15 & replace | Simvastatin - 40mg tablets: one tablet at night @ Ischaemic Heart Disease → Rosuvastatin - 5 mg once daily @ Ischaemic Heart Disease & Simvastatin and clarithromycin interact via CYP3A4 inhibition, risking rhabdomyolysis. Atorvastatin also undergoes extensive CYP3A4 metabolism, maintaining high interaction risk with clarithromycin. Rosuvastatin is minimally metabolized by CYP3A4, safely breaking this interaction while preserving secondary prevention for ischaemic heart disease. A conservative dose of 5 mg once daily is selected given renal impairment (CrCl 21 mL/min) and advanced age to ensure safety and cardiovascular efficacy. & Incorrect &  \\
\hline
scot\_09\_22 & 1 & kept | Atorvastatin - 80mg tablets: one tablet daily @ Ischaemic Heart Disease, Type 2 Diabetes Mellitus &  & Acceptable &  \\
\hline
scot\_09\_22 & 2 & kept | Finasteride - 5mg tablets: one tablet daily @ Benign Prostatic Hyperplasia &  & Correct &  \\
\hline
scot\_09\_22 & 3 & kept | Metformin - 500mg tablets: one tablet twice daily @ Type 2 Diabetes Mellitus &  & Acceptable &  \\
\hline
scot\_09\_22 & 4 & kept | Mirtazapine - 30mg tablets: one tablet daily @ Depression &  & Correct &  \\
\hline
scot\_09\_22 & 5 & kept | Ramipril - 10mg capsules: one capsule daily @ Ischaemic Heart Disease, Type 2 Diabetes Mellitus, Hypertension &  & Acceptable &  \\
\hline
scot\_09\_22 & 6 & kept | Tamsulosin - 400 micrograms modified-release capsules: one capsule daily @ Benign Prostatic Hyperplasia &  & Correct &  \\
\hline
scot\_09\_22 & 7 & discontinue | Co-Codamol - 30mg/500mg tablets: one or two tablets up to four times a day @ No health problem indication → (discontinued) & Co-codamol 30mg/500mg lacks an active documented pain indication or pain diagnosis in the chart. In an 83-year-old moderately frail resident (Rockwood 6) already receiving mirtazapine and prochlorperazine, unnecessary opioid therapy confers substantial risks of CNS depression, falls, confusion, and worsening constipation without therapeutic benefit. Discontinuing co-codamol resolves this unnecessary drug therapy and mitigates fall and sedative risks. & Correct &  \\
\hline
scot\_09\_22 & 8 & discontinue | Dapagliflozin - 10mg tablets: one tablet daily @ Ischaemic Heart Disease, Type 2 Diabetes Mellitus → (discontinued) & Dapagliflozin remains indicated long-term for type 2 diabetes mellitus and ischaemic heart disease risk reduction. However, during an acute illness with nausea and vomiting treated with prochlorperazine, continuing dapagliflozin alongside furosemide poses a severe risk of volume depletion, acute kidney injury, and euglycemic DKA. Sick day guidance mandates a temporary pause. & Correct &  \\
\hline
scot\_09\_22 & 9 & discontinue | Lansoprazole - 30mg capsules: one capsule daily @ No health problem indication → (discontinued) & Lansoprazole 30mg has no documented gastrointestinal indication (such as GERD or peptic ulcer disease) and the patient is not prescribed gastro-toxic medications requiring gastroprotection. Furthermore, ongoing use alongside furosemide increases the risk of hypomagnesemia, alongside potential infection and bone fracture risks in this frail elderly resident. Discontinuing lansoprazole resolves unnecessary drug therapy without compromising care. & Correct &  \\
\hline
scot\_09\_22 & 10 & discontinue | Macrogol 3350 - sachets: one or two sachets daily @ No health problem indication → (discontinued) & Macrogol 3350 was linked to co-codamol for opioid-induced constipation, but co-codamol has been removed from the patient's regimen. With no documented history or active diagnosis of chronic constipation, and senna remaining available, there is no valid ongoing indication for macrogol 3350 sachets. Deprescribing reduces unnecessary pill burden and polypharmacy in this moderately frail resident. & Correct &  \\
\hline
scot\_09\_22 & 11 & discontinue | Senna - 7.5mg tablets: one or two tablets at night @ No health problem indication → (discontinued) & Senna (senna\_7p5mg) has no valid medical indication following discontinuation of co-codamol and lacks documented ongoing constipation. Continued use poses unnecessary risk of electrolyte depletion and hypokalemia, particularly alongside furosemide. Discontinuing senna resolves this unnecessary therapy. & Correct &  \\
\hline
scot\_09\_22 & 12 & discontinue | Prochlorperazine - for duration of the acute illness of nausea and vomiting @ No health problem indication → (discontinued) & Prochlorperazine was prescribed for acute illness with nausea and vomiting without a continuing or documented medical indication. The patient is eating normally, and continuing this unnecessary therapy exposes this frail older resident to risks of additive CNS depression with mirtazapine, orthostatic hypotension with tamsulosin, and extrapyramidal adverse effects. Discontinuing prochlorperazine resolves the unnecessary therapy and mitigates fall and sedation risks. & Correct &  \\
\hline
scot\_09\_22 & 13 & add | (new) → Aspirin - 75 mg once daily @ Ischaemic Heart Disease & Patient has established Ischaemic Heart Disease receiving atorvastatin and ramipril, but lacks guideline-directed antiplatelet therapy for secondary prevention. Initiating low-dose aspirin 75mg daily provides essential antiplatelet coverage to reduce the risk of atherothrombotic events, recurrent myocardial infarction, and cardiovascular mortality, fulfilling an uncovered distinct therapeutic role alongside lipid-lowering and ACE inhibitor therapy. Clopidogrel was considered as an alternative non-charted option, but aspirin is standard first-line secondary prevention. & Debatable & Paciente 83 anos \\
\hline
scot\_09\_22 & 14 & add | (new) → Amlodipine - 5 mg once daily @ Hypertension & Blood pressure remains elevated at 145/85 mmHg despite ramipril 10mg daily and furosemide 40mg daily in a patient with diabetes and ischaemic heart disease. Per hypertension guidelines (NICE NG136), a dihydropyridine calcium channel blocker such as amlodipine is the preferred next-step adjunct alongside an ACE inhibitor. Amlodipine 5mg once daily is added to achieve target blood pressure control and cardiovascular risk reduction, chosen over non-dihydropyridine CCBs or thiazide diuretics given concurrent furosemide and pulse rate. & Incorrect & paciente com edema de MMII \\
\hline
scot\_09\_22 & 15 & replace | Furosemide - 40mg tablets: one tablet daily @ Hypertension → Indapamide - 2.5 mg once daily @ Hypertension & Furosemide is a loop diuretic and not a guideline-recommended first- or second-line therapy for essential hypertension in patients with normal renal function (eGFR >60 mL/min). Indapamide, a thiazide-like diuretic, provides superior cardiovascular outcomes and blood pressure efficacy without excessive volume depletion. Replacing furosemide with indapamide optimizes blood pressure control (current BP 145/85 mmHg) alongside ramipril and amlodipine while maintaining safety. & Acceptable &  \\
\hline
scot\_09\_23 & 1 & kept | Anoro Ellipta® - dry powder inhaler (DPI): one dose each morning @ Copd &  & Correct &  \\
\hline
scot\_09\_23 & 2 & kept | Atorvastatin - 40mg tablets: one tablet daily @ Cardiovascular Risk Reduction, Thromboembolism / Stroke Prophylaxis &  & Correct &  \\
\hline
scot\_09\_23 & 3 & kept | Diclofenac Gel - apply as needed @ Osteoarthritis &  & Correct &  \\
\hline
scot\_09\_23 & 4 & kept | Indapamide - 2.5mg tablets: one tablet daily @ Hypertension, Cardiovascular Risk Reduction, Thromboembolism / Stroke Prophylaxis &  & Debatable &  \\
\hline
scot\_09\_23 & 5 & kept | Lansoprazole - 15mg capsule: one capsule daily @ Dyspepsia &  & Debatable & possível de suspensão. \\
\hline
scot\_09\_23 & 6 & kept | Paracetamol - 500mg tablets: take two tablets when required for pain @ Osteoarthritis &  & Incorrect & Paciente abaixo de 50 kg. reduzir dose de paracetamol. \\
\hline
scot\_09\_23 & 7 & kept | Peptac® Liquid - 20ml twice daily @ Dyspepsia &  & Correct &  \\
\hline
scot\_09\_23 & 8 & kept | Ramipril - 10mg tablet: one tablet daily @ Hypertension, Cardiovascular Risk Reduction, Thromboembolism / Stroke Prophylaxis &  & Acceptable & É possível redução. Sem necessidade de controle rígido de pressão. \\
\hline
scot\_09\_23 & 9 & kept | Salbutamol - metered dose inhaler (MDI) 100 micrograms: two puffs when required for breathlessness @ Copd &  & Correct &  \\
\hline
scot\_09\_23 & 10 & kept | Warfarin - tablets 1mg tablets: three tablets daily (target INR 2.5) @ Non-Valvular Atrial Fibrillation, Cardiovascular Risk Reduction, Thromboembolism / Stroke Prophylaxis &  & Acceptable & Possível transicionar para DOAC \\
\hline
scot\_09\_23 & 11 & discontinue | Zopiclone - 7.5mg tablet: one tablet at night @ No health problem indication → (discontinued) & Zopiclone 7.5mg is prescribed without an established clinical indication or documented sleep disorder. In this 81-year-old mildly frail patient with osteoporosis and a history of falls/fracture, zopiclone significantly increases the risk of CNS depression, cognitive impairment, and further falls. Discontinuing this unindicated sedative directly addresses these safety hazards and aligns with fall prevention goals. & Correct &  \\
\hline
scot\_09\_23 & 12 & discontinue | Devil’S Claw Herbal Remedy @ Osteoarthritis, Dyspepsia → (discontinued) & Devil's Claw provides redundant analgesic therapy alongside paracetamol and topical diclofenac for osteoarthritis. Furthermore, it stimulates gastric acid production, directly antagonising lansoprazole and worsening active dyspepsia, while also interacting with warfarin to increase bleeding risk. Discontinuing Devil's Claw eliminates redundant therapy, reduces gastrointestinal irritation, and resolves significant drug-drug interaction risks without compromising pain control. & Correct &  \\
\hline
scot\_09\_23 & 13 & add | (new) → Capsaicin - 0.25 g four times daily @ Osteoarthritis & Persistent knee osteoarthritis pain refractory to paracetamol and diclofenac gel requires adjunct escalation. Oral NSAIDs are contraindicated due to warfarin, dyspepsia, and CrCl 36 mL/min. Topical capsaicin cream is guideline-recommended for knee osteoarthritis without systemic bleeding or renal risks, addressing sore knees. Dosage updated from 0.025\% to standardize format as a pea-sized amount (approx 0.25 g application) to resolve validation requirement for valid dosage units. & Correct &  \\
\hline
scot\_09\_23 & 14 & replace | Alendronic Acid - 70mg tablets: one tablet weekly @ Osteoporosis Fracture Prophylaxis, Bone Protection / Vitamin D And Calcium Supplementation → Denosumab - 60 mg every 6 months @ Osteoporosis Fracture Prophylaxis, Bone Protection / Vitamin D And Calcium Supplementation & Alendronic acid is causing upper gastrointestinal irritation and dyspepsia, necessitating lansoprazole and Peptac, which binds oral bisphosphonates. Furthermore, renal function (CrCl 36 mL/min) is near contraindication cutoffs. Denosumab provides potent anti-resorptive bone protection without gastrointestinal toxicity, bypasses renal clearance contraindications, and is administered subcutaneously every 6 months, avoiding oral absorption interactions with antacids. & Acceptable &  \\
\hline
scot\_09\_23 & 15 & replace | Calcium And Vitamin D - Calcium 1.5g and vitamin D 400iu chewable tablets: two tablets daily @ Osteoporosis Fracture Prophylaxis, Bone Protection / Vitamin D And Calcium Supplementation, Fall Prevention → Colecalciferol - 800 units once daily @ Osteoporosis Fracture Prophylaxis, Bone Protection / Vitamin D And Calcium Supplementation, Fall Prevention & Oral calcium carbonate exacerbates active dyspepsia and requires gastric acid for optimal absorption, which is compromised by lansoprazole. Switching to colecalciferol ensures adequate vitamin D repletion to support bone protection and denosumab therapy while removing the offending calcium salt contributing to dyspeptic symptoms. & Correct &  \\
\hline
scot\_09\_24 & 1 & kept | Alogliptin - 12.5mg tablets: one tablet daily @ Type 2 Diabetes Mellitus &  & Incorrect &  \\
\hline
scot\_09\_24 & 2 & kept | Amitriptyline - 25mg tablets: one tablet daily @ Bone Metastases, Depression, Urge Incontinence &  & Incorrect &  \\
\hline
scot\_09\_24 & 3 & kept | Atorvastatin - 20mg tablets: one tablet daily @ Type 2 Diabetes Mellitus &  & Incorrect &  \\
\hline
scot\_09\_24 & 4 & kept | Dexamethasone - 2mg tablets: three tablets daily @ Metastatic Non-Small Cell Lung Cancer, Brain Metastases, Liver Metastases, Bone Metastases, Chronic Obstructive Pulmonary Disease &  & Correct &  \\
\hline
scot\_09\_24 & 5 & kept | Lidocaine - 700mg plasters: Apply two plasters once daily for 12 hours, followed by a 12-hour plaster-free period @ Bone Metastases &  & Correct &  \\
\hline
scot\_09\_24 & 6 & kept | Metformin - 500mg tablets: one twice a day @ Type 2 Diabetes Mellitus &  & Incorrect &  \\
\hline
scot\_09\_24 & 7 & kept | Omeprazole - 20mg capsules: one capsule daily @ Gastroprotection &  & Debatable &  \\
\hline
scot\_09\_24 & 8 & kept | Paracetamol - 500mg tablets: two tablets four times a day @ Metastatic Non-Small Cell Lung Cancer, Liver Metastases, Bone Metastases &  & Correct &  \\
\hline
scot\_09\_24 & 9 & kept | Senna - 7.5mg tablets: two tablets at night @ Constipation Prophylaxis &  & Acceptable &  \\
\hline
scot\_09\_24 & 10 & kept | Salbutamol - 100 micrograms pMDI: two doses as required for breathlessness @ Metastatic Non-Small Cell Lung Cancer, Chronic Obstructive Pulmonary Disease &  & Acceptable &  \\
\hline
scot\_09\_24 & 11 & discontinue | Diazepam - 5mg tablets: one tablet three times a day @ No health problem indication → (discontinued) & Diazepam (False) has no documented clinical indication and carries a severe risk of additive central nervous system and respiratory depression, particularly in the context of underlying COPD and concurrent therapy with opioids (morphine) and pregabalin. Discontinuing diazepam eliminates this unnecessary risk without compromising therapy for any diagnosed condition. & Correct &  \\
\hline
scot\_09\_24 & 12 & discontinue | Folic Acid - 5mg tablets: one tablet daily @ No health problem indication → (discontinued) & Folic acid 5mg daily has no documented clinical indication (such as folate deficiency anemia or methotrexate therapy) in this patient with advanced metastatic lung cancer. Discontinuing it safely reduces pill burden, polypharmacy, and swallowing discomfort given her poor oral intake and dysphagia. & Correct &  \\
\hline
scot\_09\_24 & 13 & add | (new) → Denosumab - 120 mg every 4 weeks @ Metastatic Non-Small Cell Lung Cancer & Patient has metastatic NSCLC complicated by documented bone metastases and uncontrolled pain. Guidelines recommend a bone-modifying agent to prevent skeletal-related events and improve bone pain control. Denosumab 120 mg subcutaneously every 4 weeks is preferred over zoledronic acid given the patient's impaired renal function (eGFR 40 mL/min/1.73m2). Monitoring for hypocalcemia and osteonecrosis of the jaw is advised alongside baseline dental evaluation and calcium/vitamin D supplementation. & Debatable &  \\
\hline
scot\_09\_24 & 14 & increase\_dosage | Levothyroxine - 100microgram tablets: two tablets daily @ Hypothyroidism → Levothyroxine - 225 mcg once daily @ Hypothyroidism & Levothyroxine absorption and efficacy are reduced by concurrent omeprazole (decreased gastric acidity impairs dissolution) and sertraline. Current dose of 200 mcg daily is insufficient due to these interactions. Increasing dose to 225 mcg daily compensates for reduced bioavailability (resolving dose format validation issue). Monitor TSH and cardiac rhythm given paroxysmal atrial fibrillation history. & Debatable &  \\
\hline
scot\_09\_24 & 15 & increase\_dosage | Macrogol 3350 - sachets: one sachet daily @ Constipation Prophylaxis → Macrogol 3350 - 26.25 g once daily @ Constipation Prophylaxis & Macrogol 3350 at 1 sachet (13.125 g) daily is insufficient to prevent and treat constipation in this patient with active constipation receiving regular and PRN morphine, amitriptyline, and solifenacin. Dosage format adjusted to grams (26.25 g) to address validation failure: recommended\_dosage: Dosage format unclear - missing units or frequency: '2 sachets'. Increasing to 26.25 g (2 sachets) daily provides effective osmotic laxation alongside charted senna. & Correct &  \\
\hline
scot\_09\_24 & 16 & increase\_dosage | Morphine Mr - 30mg capsules: one capsule twice a day @ Metastatic Non-Small Cell Lung Cancer, Liver Metastases, Bone Metastases → Morphine Mr - 45 mg twice daily @ Metastatic Non-Small Cell Lung Cancer, Liver Metastases, Bone Metastases & Patient has uncontrolled cancer pain and limited mobility from metastatic NSCLC with bone and liver metastases despite Morphine MR 30mg twice daily and PRN immediate-release morphine. Palliative guidelines recommend titrating baseline modified-release opioids by 30-50\% for persistent background pain. Increasing Morphine MR to 45mg twice daily provides improved analgesia. Monitor for sedation, respiratory status, and constipation alongside current laxatives. Moderate risk reflects opioid escalation with renal impairment (eGFR 40 mL/min). & Debatable &  \\
\hline
scot\_09\_24 & 17 & increase\_dosage | Morphine Ir - 10mg tablets: One every four to six hours as required for pain relief @ Metastatic Non-Small Cell Lung Cancer, Liver Metastases, Bone Metastases, Brain Metastases → Morphine Ir - 15 mg every four to six hours as required @ Metastatic Non-Small Cell Lung Cancer, Liver Metastases, Bone Metastases, Brain Metastases & Patient has uncontrolled cancer pain on Morphine MR 45 mg twice daily (total 90 mg/24h). Breakthrough doses of immediate-release morphine should typically be 1/10th to 1/6th of total regular daily morphine (10-15 mg, standardly 15 mg). The current 10 mg PRN dose is inadequate for severe breakthrough pain episodes. Increasing Morphine Ir to 15 mg PRN aligns with standard palliative dosing proportions. Renal function (eGFR 40 mL/min) and CNS sedation should be routinely monitored, keeping risk moderate. & Correct &  \\
\hline
scot\_09\_24 & 18 & replace | Sertraline - 150mg tablets: one tablet daily @ Depression → Mirtazapine - 15 mg once daily @ Depression & Sertraline concurrently interacts with amitriptyline (additive serotonergic toxicity and elevated TCA levels) and with edoxaban plus dexamethasone (substantially increasing GI bleeding risk via platelet inhibition). Replacing sertraline with mirtazapine resolves the serotonergic reuptake inhibition and antiplatelet-related bleeding interactions, while safely maintaining antidepressant coverage for depression and offering beneficial appetite and sleep support in advanced malignancy. & Acceptable &  \\
\hline
scot\_09\_24 & 19 & replace | Solifenacin - 10mg tablets: one tablet daily @ Urge Incontinence → Mirabegron - 25 mg once daily @ Urge Incontinence & Solifenacin is an antimuscarinic agent aggravating dry mouth, sore mouth on swallowing, and constipation, especially alongside morphine, amitriptyline, and umeclidinium. Replacing solifenacin with mirabegron (a beta-3 adrenergic agonist) maintains effective control of urge incontinence while eliminating antimuscarinic adverse effects. The 25 mg dose is appropriate given her renal impairment (eGFR 40 mL/min/1.73 m2) and normotensive status. & Debatable &  \\
\hline
scot\_09\_24 & 20 & replace | Trelegy Ellipta - ® 22/92/55 micrograms DPI: one inhalation daily @ Chronic Obstructive Pulmonary Disease → Anoro Ellipta - 55/22 mcg once daily @ Chronic Obstructive Pulmonary Disease & Trelegy Ellipta contains fluticasone furoate, which increases pneumonia risk following a recent lower respiratory tract infection and duplicates steroid exposure with high-dose dexamethasone, worsening oral pain. Switching to Anoro Ellipta (umeclidinium/vilanterol) maintains optimal dual bronchodilation (LAMA/LABA) for COPD without the safety hazards of an inhaled corticosteroid. & Correct &  \\
\hline
scot\_09\_24 & 21 & reduce\_dosage | Edoxaban - 60mg tablets: one tablet daily @ Paroxysmal Atrial Fibrillation → Edoxaban - 30 mg once daily @ Paroxysmal Atrial Fibrillation & Edoxaban 60 mg once daily is excessive given the patient's moderate renal impairment (eGFR 40 mL/min/1.73 m2) and brain metastases, posing an elevated risk of major intracranial and systemic bleeding. Reducing edoxaban to 30 mg once daily aligns with renal dose adjustment criteria (CrCl/eGFR 15-50 mL/min) for stroke prevention in paroxysmal atrial fibrillation, preserving anticoagulant efficacy while mitigating hemorrhage risk. & Debatable &  \\
\hline
scot\_09\_24 & 22 & reduce\_dosage | Pregabalin - 300mg capsules: one capsule twice a day @ Bone Metastases → Pregabalin - 150 mg twice daily @ Bone Metastases & Pregabalin 300 mg twice daily (600 mg/day) exceeds recommended renal dosing limits for the patient's eGFR of 40 mL/min/1.73 m2, which restricts the maximum daily dose to 300 mg/day. This excessive dose increases the risk of severe CNS depression, cognitive impairment, and sedation, especially alongside opioids. Reducing pregabalin to 150 mg twice daily (300 mg/day) aligns with renal guidelines to maintain neuropathic analgesia for bone metastases while minimizing toxicity. & Correct &  \\
\hline
scot\_09\_24 & 23 & reduce\_dosage | Ramipril - 10mg capsules: one capsule daily @ Type 2 Diabetes Mellitus, Hypertension → Ramipril - 5 mg once daily @ Type 2 Diabetes Mellitus, Hypertension & Patient is receiving Ramipril 10mg daily despite borderline low blood pressure (110/60 mmHg), reduced eGFR (40 mL/min/1.73m2), poor oral intake, and advanced palliative metastatic disease. Continuing the maximum dose poses significant risks of symptomatic hypotension, hypoperfusion, and acute kidney injury. Reducing Ramipril to 5mg once daily maintains cardiovascular and renal protection for hypertension and type 2 diabetes mellitus while mitigating hypotension and renal decline in this frail palliative context. & Incorrect & Suspensão \\
\hline
scot\_09\_25 & 1 & kept | Fentanyl - patch 25micrograms/hour patch: one patch to be changed every 72 hours @ Locally Advanced Oesophageal Cancer, Lung Metastases, Bone Metastases &  & Correct &  \\
\hline
scot\_09\_25 & 2 & kept | Insulin Detemir - 100 unit per 1 ml injection (Levemir®): SC injection twice daily @ Type 1 Diabetes Mellitus &  & Incorrect &  \\
\hline
scot\_09\_25 & 3 & kept | Insulin Aspart - 100 unit per 1ml injection (NovoRapid®): SC injection 15mins before each meal @ Type 1 Diabetes Mellitus &  & Incorrect &  \\
\hline
scot\_09\_25 & 4 & kept | Laxido - ® sachets: one sachet twice daily @ Locally Advanced Oesophageal Cancer, Bowel Regimen / Constipation Prophylaxis &  & Debatable &  \\
\hline
scot\_09\_25 & 5 & kept | Metoclopramide - 10mg tablets: one tablet three times daily @ Locally Advanced Oesophageal Cancer &  & Incorrect &  \\
\hline
scot\_09\_25 & 6 & discontinue | Zopiclone - 7.5mg tablets: one tablet at night @ No health problem indication → (discontinued) & Zopiclone (zopiclone\_7p5mg) lacks a documented clinical indication or sleep disorder diagnosis. In this frail patient with poor performance status (AKPS 20\%), baseline confusion, hypoxemia (SpO2 85\%), and concomitant opioid therapy (fentanyl, morphine), continuing zopiclone poses severe risks of additive CNS and respiratory depression. Discontinuing zopiclone removes unnecessary pharmacotherapy risk without compromising palliative symptom management. & Correct &  \\
\hline
scot\_09\_25 & 7 & add | (new) → Dexamethasone - 4 mg once daily @ Locally Advanced Oesophageal Cancer & For locally advanced oesophageal cancer with dysphagia, fatigue, and severe debility (AKPS 20\%), corticosteroid therapy with dexamethasone provides essential palliative adjunct escalation. Per palliative oncology guidelines, short-term dexamethasone helps reduce peritumoral oedema to relieve malignant dysphagia and improves cancer-associated anorexia, fatigue, and pain, complementing existing analgesia and antiemetics. & Correct &  \\
\hline
scot\_09\_25 & 8 & add | (new) → Oxygen - 2 l/min continuously @ Lung Metastases & Patient with lung metastases has severe hypoxemia (SpO2 85\% on room air) and shortness of breath. Current analgesics (fentanyl, morphine) address pain, but supplemental oxygen therapy is guideline-indicated to relieve dyspnea and correct severe hypoxemic respiratory compromise in palliative advanced cancer, aligning with palliative home care preferences without requiring acute hospital transfer. & Correct &  \\
\hline
scot\_09\_25 & 9 & add | (new) → Senna - 15 mg once daily at bedtime @ Bowel Regimen / Constipation Prophylaxis & Patient has persistent constipation on dual opioids (fentanyl, morphine) and immobility (AKPS 20\%) despite Laxido. Palliative bowel management guidelines recommend combining an osmotic laxative with a stimulant laxative to effectively prevent and treat opioid-induced constipation. Adding senna provides necessary propulsive motility synergistic with Laxido's osmotic action. & Acceptable &  \\
\hline
scot\_09\_25 & 10 & replace | Fluconazole - 50mg capsules: one capsule daily (started two days ago) @ Oral Thrush → Nystatin - 100000 units four times daily @ Oral Thrush & Fluconazole 50mg solid capsules are inappropriate due to severe dysphagia, poor swallow, and bedbound status from advanced oesophageal cancer. Replacing solid fluconazole with Nystatin oral suspension allows topical contact, easier administration, and effective treatment of oral thrush without swallowing difficulties or major CYP3A4/CYP2C9 drug interactions. & Correct &  \\
\hline
scot\_09\_25 & 11 & replace | Fluoxetine - 40mg capsules: one capsule daily @ Anxiety, Depression → Mirtazapine - 15 mg once daily @ Anxiety, Depression & Patient has severe dysphagia from advanced esophageal cancer and oral thrush, making solid fluoxetine capsules inappropriate due to aspiration risk and impaired absorption. Fluoxetine liquid is difficult to swallow due to painful mouth and volume. Switching to mirtazapine orally disintegrating tablets (SolTab) provides effective treatment for depression and anxiety without swallowing whole capsules, while also addressing poor appetite and insomnia common in palliative care. Moderate risk due to potential additive sedation with existing opioids, requiring routine monitoring. & Correct &  \\
\hline
scot\_09\_25 & 12 & replace | Ibuprofen - 400mg tablets: one tablet three times a day @ Locally Advanced Oesophageal Cancer, Lung Metastases, Bone Metastases → Ketorolac - 30 mg once daily @ Locally Advanced Oesophageal Cancer, Lung Metastases, Bone Metastases & Ibuprofen 400 mg oral tablets cannot be swallowed reliably due to locally advanced oesophageal cancer, severe dysphagia, and painful oral thrush, rendering this solid dosage form ineffective for cancer-related bone pain. Replacing it with continuous subcutaneous ketorolac restores effective parenteral NSAID analgesia for metastatic bone pain while bypassing the oral route. Ketorolac is chosen over non-charted alternatives like oral liquid NSAIDs, which remain unsuitable given the patient's compromised swallow. Gastroprotection with charted omeprazole is maintained. & Acceptable &  \\
\hline
scot\_09\_25 & 13 & replace | Omeprazole - 20mg capsules: two capsules daily @ Locally Advanced Oesophageal Cancer, Gastroprotection → Pantoprazole - 40 mg once daily @ Locally Advanced Oesophageal Cancer, Gastroprotection & Patient has advanced oesophageal cancer with severe dysphagia, poor swallow, and painful oral thrush, rendering solid oral omeprazole capsules inappropriate and unsafe. In palliative dysphagia, replacing with subcutaneous or IV pantoprazole maintains acid suppression and gastroprotection (especially with dexamethasone/ketorolac therapy) while avoiding oral route failure. Subcutaneous pantoprazole 40 mg once daily provides equivalent gastroprotection without relying on swallowing intact capsules. & Acceptable &  \\
\hline
scot\_09\_25 & 14 & replace | Paracetamol - 500mg tablets: two tablets four times a day @ Locally Advanced Oesophageal Cancer, Lung Metastases, Bone Metastases → Paracetamol IV - 1000 mg four times daily @ Locally Advanced Oesophageal Cancer, Lung Metastases, Bone Metastases & Paracetamol oral tablets are an inappropriate dosage form due to the patient's locally advanced oesophageal cancer and severe dysphagia, creating risks of aspiration and undertreated pain. To maintain baseline non-opioid multimodal analgesia for bone and cancer-related pain without requiring oral tablet swallowing, replace oral paracetamol with intravenous paracetamol at the equivalent standard dose. & Correct &  \\
\hline
scot\_09\_25 & 15 & replace | Morphine Sulfate - 10mg/5ml oral solution, 10mg (5ml) every four hours if needed for pain relief @ Locally Advanced Oesophageal Cancer, Lung Metastases, Bone Metastases → Morphine Sulfate (SC injection) - 5 mg every four hours as needed @ Locally Advanced Oesophageal Cancer, Lung Metastases, Bone Metastases & Morphine Sulfate oral solution poses a significant aspiration risk due to the patient's dysphagia and advanced oesophageal cancer. Subcutaneous morphine provides a reliable, non-oral parenteral route for breakthrough pain without aspiration hazard, maintaining effective analgesia. Co-administered with transdermal fentanyl patch; IV paracetamol is also charted. & Acceptable &  \\
\hline
scot\_09\_25 & 16 & reduce\_dosage | Ramipril - 10mg capsules: one capsule daily @ Hypertension, Type 1 Diabetes Mellitus → Ramipril - 2.5 mg once daily @ Hypertension, Type 1 Diabetes Mellitus & Ramipril 10 mg daily is excessive given the patient's low blood pressure (105/55 mmHg), bedbound state, and palliative goals in advanced malignancy. Continuing the maximum dose risks symptomatic hypotension and organ hypoperfusion. Reducing the dose to 2.5 mg once daily safely de-escalates therapy, mitigating the risk of hypotension while preserving renal/cardiovascular protective benefits for hypertension and type 1 diabetes mellitus. & Incorrect &  \\
\hline
scot\_09\_26 & 1 & discontinue | Bisacodyl - 5mg tablets: one or two at night as required @ No health problem indication → (discontinued) & Bisacodyl is prescribed PRN without a documented diagnosis, symptom, or indication of constipation in the chart. Deprescribing bisacodyl eliminates unnecessary medication and reduces the risk of laxative-associated adverse effects such as abdominal cramping and electrolyte disturbances. Although the patient is receiving opioid analgesics, there is no documented constipation to justify continuing bisacodyl at this time. & Correct &  \\
\hline
scot\_09\_26 & 2 & discontinue | Gabapentin - 300mg capsules: one capsule three times daily @ Chronic Pain: Hip And Back → (discontinued) & Gabapentin represents duplicate gabapentinoid therapy alongside concurrent pregabalin for chronic hip and back pain with radiculopathy. Co-prescribing two distinct alpha-2-delta ligands provides redundant mechanism of action while substantially heightening risks of severe CNS and respiratory depression, particularly in the context of concurrent opioid therapy (methadone, co-codamol, tramadol). Discontinuing gabapentin eliminates therapeutic duplication while maintaining pregabalin for neuropathic pain control. & Correct &  \\
\hline
scot\_09\_26 & 3 & discontinue | Tramadol - 50mg capsules: one capsule three times daily @ Chronic Pain: Hip And Back → (discontinued) & Tramadol represents redundant opioid therapy alongside co-codamol and methadone, compounding risks of central nervous system and respiratory depression, serotonin syndrome (with sertraline and amitriptyline), and seizures. Chronic hip and back pain remains covered by ongoing analgesics. Discontinuing tramadol eliminates unnecessary opioid duplication and significant toxicological hazards without compromising baseline coverage. & Correct &  \\
\hline
scot\_09\_26 & 4 & replace | Amitriptyline - 50mg tablets: two tablets daily @ Depression, Chronic Pain: Hip And Back → Duloxetine - 30 mg once daily @ Depression, Chronic Pain: Hip And Back & Amitriptyline 100 mg daily has failed to control neuropathic pain symptoms, which are worsening alongside depression despite combination with pregabalin. Duloxetine is indicated for both neuropathic pain and depression and provides dual SNRI reuptake inhibition, offering an effective single-agent replacement for amitriptyline's dual indication role while avoiding tricyclic anticholinergic and cardiotoxic risks. Titrate sertraline down concurrently to mitigate serotonergic duplication. & Acceptable &  \\
\hline
scot\_09\_26 & 5 & replace | Co-Codamol - 30/500mg tablets: two tablets four times daily @ Chronic Pain: Hip And Back → Naproxen - 500 mg twice daily @ Chronic Pain: Hip And Back & Co-codamol has proven ineffective and refractory for chronic hip and back pain with radicular symptoms despite full doses, while presenting dangerous opioid duplication alongside methadone and drug dependency. Replacing it with naproxen provides targeted non-opioid anti-inflammatory analgesia for musculoskeletal back pain without compounding opioid toxicity or central nervous system depression. & Acceptable &  \\
\hline
scot\_09\_26 & 6 & increase\_dosage | Pregabalin - 150mg capsules: one capsule twice daily @ Chronic Pain: Hip And Back → Pregabalin - 300 mg twice daily @ Chronic Pain: Hip And Back & Pregabalin 150 mg twice daily (300 mg/day) provides incomplete control for chronic hip and back pain with radicular neuropathic symptoms (tingling, shooting pain). Titrating to 300 mg twice daily (600 mg/day maximum guideline dose) addresses persistent symptoms. Given concurrent methadone, CNS depression and sedation risk are increased, warranting a moderate risk rating and close monitoring during titration. & Acceptable &  \\
\hline
scot\_09\_26 & 7 & replace | Methadone - 1mg/ml oral solution: 30ml daily, twice weekly pick-up @ Chronic Pain: Hip And Back, Drug Dependency → Buprenorphine/naloxone - 8/2 mg once daily @ Chronic Pain: Hip And Back, Drug Dependency & Methadone combined with sertraline increases the risk of additive QT prolongation and serotonin toxicity. Transitioning to buprenorphine/naloxone maintains effective opioid agonist therapy for drug dependency and chronic pain while substantially lowering the risk of QT prolongation and fatal cardiac arrhythmias. & Correct &  \\
\hline
scot\_09\_26 & 8 & replace | Sertraline - 100mg tablets: one tablet daily @ Depression → Bupropion - 150 mg once daily @ Depression & Sertraline co-administered with newly added duloxetine creates duplicate serotonergic exposure and elevates the risk of serotonin syndrome. Replacing sertraline with bupropion effectively maintains pharmacotherapy for depression via a non-serotonergic mechanism (norepinephrine-dopamine reuptake inhibition), eliminating the serotonergic drug interaction while ensuring active management of depression. & Acceptable &  \\
\hline
\end{longtable}
\begin{longtable}{|p{1.45cm}|c|p{10.2cm}|p{1.55cm}|p{2.3cm}|}
\caption{Per-medicine control recommendations and the specialist's recommendation label.}\label{tab:appsheetscontrol}\\
\hline
Case & No. & Recommendation & Label & Note \\
\hline
\endfirsthead
\multicolumn{5}{l}{\tablename\ \thetable, continued}\\
\hline
Case & No. & Recommendation & Label & Note \\
\hline
\endhead
\hline
\endfoot
scot\_09\_01 & 1 & Target Medication: Amisulpride 100mg tablets Recommendation: Reduce dose to 50mg twice daily with a view to tapering to the lowest effective dose to manage verbal aggression while reducing fall risk and sedation. & Debatable &  \\
\hline
scot\_09\_01 & 2 & Target Medication: Trazodone 150mg capsules Recommendation: Reduce dose to 50mg to 100mg at night to decrease excessive daytime somnolence and fall risk. & Incorrect & correto seria suspensão \\
\hline
scot\_09\_01 & 3 & Target Medication: Tramadol 50mg capsules Recommendation: Gradually taper and discontinue, assessing for ongoing pain, to reduce sedation, confusion, and fall risk. & Correct &  \\
\hline
scot\_09\_01 & 4 & Target Medication: Paracetamol 500mg tablets Recommendation: Initiate 500mg to 1000mg up to four times daily as needed for chronic pain management as a safer alternative to tramadol. & Correct &  \\
\hline
scot\_09\_01 & 5 & Target Medication: Bendroflumethiazide 2.5mg tablets Recommendation: Discontinue, as it is ineffective for symptomatic peripheral oedema and may cause electrolyte disturbances. & Correct &  \\
\hline
scot\_09\_01 & 6 & Target Medication: Furosemide 20mg-40mg tablets Recommendation: Initiate 20mg to 40mg once daily in the morning to treat fluid overload, leg oedema, and breathlessness, subject to clinical evaluation. & Correct &  \\
\hline
scot\_09\_01 & 7 & Target Medication: Cetirizine 10mg tablets Recommendation: Discontinue unless there is an active, documented allergic indication, to reduce anticholinergic and sedative burden. & Correct &  \\
\hline
scot\_09\_01 & 8 & Target Medication: Fucibet® cream Recommendation: Discontinue or limit to a short defined course (typically max 2 weeks) to avoid steroid-induced skin atrophy and antimicrobial resistance; switch to emollient maintenance. & Correct &  \\
\hline
scot\_09\_01 & 9 & Target Medication: Diprobase® cream Recommendation: Keep the medication in the regimen. & Debatable &  \\
\hline
scot\_09\_01 & 10 & Target Medication: Thiamine 50mg tablets Recommendation: Keep the medication in the regimen. & Incorrect & suspender \\
\hline
scot\_09\_01 & 11 & Target Medication: Colecalciferol (Vitamin D3) / Calcium supplementation Recommendation: Initiate appropriate bone protection (e.g., Colecalciferol 800 IU daily with or without calcium depending on dietary intake) for secondary prevention of fragility fractures. & Acceptable &  \\
\hline
scot\_09\_02 & 1 & Target Medication: Adcal-D3 1500mg/400iu tablets Recommendation: Keep the medication in the regimen. & Acceptable &  \\
\hline
scot\_09\_02 & 2 & Target Medication: Alendronic acid 70mg tablets Recommendation: Discontinue the medication due to severe renal impairment (CrCl <35 mL/min). & Correct &  \\
\hline
scot\_09\_02 & 3 & Target Medication: Atorvastatin 80mg tablets Recommendation: Reduce dose to Atorvastatin 20mg tablets, one tablet daily at night. & Correct &  \\
\hline
scot\_09\_02 & 4 & Target Medication: Clopidogrel 75mg tablets Recommendation: Discontinue the medication to minimize bleeding risk, continuing anticoagulation monotherapy for atrial fibrillation and secondary stroke prevention. & Correct &  \\
\hline
scot\_09\_02 & 5 & Target Medication: Digoxin 125microgram tablets Recommendation: Discontinue the medication (or withhold pending serum digoxin level and ECG) due to bradycardia (HR 56 bpm), low weight, and impaired renal clearance. & Correct &  \\
\hline
scot\_09\_02 & 6 & Target Medication: Edoxaban 60mg tablets Recommendation: Reduce dose to Edoxaban 30mg tablets, one tablet daily, due to body weight ≤60 kg and CrCl 15–50 mL/min. & Correct &  \\
\hline
scot\_09\_02 & 7 & Target Medication: Gliclazide 80mg tablets Recommendation: Discontinue the medication to prevent severe hypoglycaemia, given the overtight HbA1c (34 mmol/mol) and high fall risk. & Correct &  \\
\hline
scot\_09\_02 & 8 & Target Medication: Indapamide 2.5mg tablets Recommendation: Discontinue the medication to address profound postural hypotension and reduce fall risk. & Correct &  \\
\hline
scot\_09\_02 & 9 & Target Medication: Levothyroxine 100microgram tablets Recommendation: Consolidate levothyroxine regimen into a single daily formulation or adjust based on a recent thyroid function test. & Correct &  \\
\hline
scot\_09\_02 & 10 & Target Medication: Levothyroxine 50microgram tablets Recommendation: Discontinue this separate strength and supply the total daily dose as a single tablet/strength. & Correct &  \\
\hline
scot\_09\_02 & 11 & Target Medication: Metformin 500mg tablets Recommendation: Discontinue the medication due to severe renal impairment (CrCl 20 mL/min, risk of lactic acidosis) and appetite suppression. & Correct &  \\
\hline
scot\_09\_02 & 12 & Target Medication: Mirtazapine 30mg tablets Recommendation: Keep the medication in the regimen. & Incorrect &  \\
\hline
scot\_09\_02 & 13 & Target Medication: Omeprazole 20mg capsules Recommendation: Keep the medication in the regimen. & Incorrect &  \\
\hline
scot\_09\_02 & 14 & Target Medication: Paracetamol 500mg tablets Recommendation: Keep the medication in the regimen. & Correct &  \\
\hline
scot\_09\_02 & 15 & Target Medication: Perindopril 4mg tablets Recommendation: Reduce dose to Perindopril 2mg tablets once daily (or temporarily withhold) and monitor sitting/standing blood pressure and renal function. & Correct &  \\
\hline
scot\_09\_02 & 16 & Target Medication: Tolterodine MR 4mg capsules Recommendation: Discontinue the medication to reduce anticholinergic burden, improve cognition, and re-evaluate baseline urinary function. & Correct &  \\
\hline
scot\_09\_02 & 17 & Target Medication: Zopiclone 7.5mg tablets Recommendation: Taper and discontinue the medication to reduce sedation, delirium, and nocturnal fall risk. & Correct &  \\
\hline
scot\_09\_03 & 1 & Target Medication: Gliclazide Recommendation: Discontinue gliclazide to prevent severe hypoglycemia given an HbA1c of 38 mmol/mol. & Correct &  \\
\hline
scot\_09\_03 & 2 & Target Medication: Lisinopril Recommendation: Consolidate the duplicate orders (10mg and 20mg) and reduce dose to 10mg once daily (or discontinue if hypotension persists). & Correct &  \\
\hline
scot\_09\_03 & 3 & Target Medication: Amlodipine Recommendation: Reduce amlodipine dose to 5mg once daily to prevent over-treatment of blood pressure and reduce fall risk. & Correct &  \\
\hline
scot\_09\_03 & 4 & Target Medication: Gabapentin Recommendation: Taper and discontinue gabapentin, as the patient reports no current pain, the current dose is excessively high for his renal function (CrCl 39 mL/min), and it carries a high risk of sedation and falls. & Correct &  \\
\hline
scot\_09\_03 & 5 & Target Medication: Lorazepam Recommendation: Gradually taper and discontinue lorazepam to reduce sedation, cognitive worsening, and fall risk. & Correct &  \\
\hline
scot\_09\_03 & 6 & Target Medication: Atorvastatin Recommendation: Reduce atorvastatin from 80mg to 20mg once daily to minimize the risk of adverse effects in severe frailty and low body weight. & Debatable &  \\
\hline
scot\_09\_03 & 7 & Target Medication: Folic acid Recommendation: Discontinue folic acid if recent full blood count and serum folate levels are within normal limits. & Correct &  \\
\hline
scot\_09\_03 & 8 & Target Medication: Paracetamol Recommendation: Convert paracetamol from a regular schedule (1g four times daily) to "as needed" (PRN) for pain, as he currently reports no pain and is low body weight (54 kg). & Acceptable &  \\
\hline
scot\_09\_03 & 9 & Target Medication: Adcal-D3 Recommendation: Keep the medication in the regimen. & Acceptable &  \\
\hline
scot\_09\_03 & 10 & Target Medication: Aspirin Recommendation: Keep the medication in the regimen. & Debatable &  \\
\hline
scot\_09\_03 & 11 & Target Medication: Omeprazole Recommendation: Keep the medication in the regimen. & Debatable &  \\
\hline
scot\_09\_03 & 12 & Target Medication: Sertraline Recommendation: Keep the medication in the regimen. & Correct &  \\
\hline
scot\_09\_03 & 13 & Target Medication: Zoledronic acid Recommendation: Consider initiating an anti-resorptive agent (such as annual intravenous zoledronic acid 5mg or an oral bisphosphonate if suitable) for secondary fracture prevention following his hip fracture. & Acceptable &  \\
\hline
scot\_09\_04 & 1 & Target Medication: Metoclopramide 10mg tablets Recommendation: Discontinue metoclopramide due to the risk of extrapyramidal side effects and exceeding the recommended 5-day duration of therapy. & Correct &  \\
\hline
scot\_09\_04 & 2 & Target Medication: Tramadol 50mg capsules Recommendation: Taper and discontinue tramadol to reduce nausea, cognitive burden, and fall risk, transitioning pain management to a 7-day low-dose transdermal buprenorphine patch (e.g., 5 micrograms/hour). & Correct &  \\
\hline
scot\_09\_04 & 3 & Target Medication: Buprenorphine 5 micrograms/hour transdermal patch Recommendation: Initiate one 5 micrograms/hour transdermal patch applied once every 7 days to provide continuous pain management without oral pill burden. & Acceptable &  \\
\hline
scot\_09\_04 & 4 & Target Medication: Paracetamol 500mg tablets Recommendation: Switch paracetamol to soluble/dispersible tablets (500mg two tablets four times daily dissolved in water) or paracetamol oral suspension (1000mg four times daily) to alleviate difficulties with swallowing and chewing halved tablets. & Correct &  \\
\hline
scot\_09\_04 & 5 & Target Medication: Lorazepam 500 micrograms tablets Recommendation: Plan a gradual dose reduction and discontinuation of regular and PRN lorazepam to minimize fall risk, sedation, and cognitive worsening, utilizing non-pharmacological approaches for agitation. & Correct &  \\
\hline
scot\_09\_04 & 6 & Target Medication: Metformin 500mg tablets Recommendation: Withhold or discontinue metformin temporarily while acute nausea resolves; if glycemic management remains necessary, consider switching to a well-tolerated once-daily agent with minimal GI burden and low hypoglycemia risk (e.g., linagliptin 5mg daily). & Debatable &  \\
\hline
scot\_09\_04 & 7 & Target Medication: Furosemide 40mg tablets (regular) Recommendation: Keep the medication in the regimen. & Incorrect &  \\
\hline
scot\_09\_04 & 8 & Target Medication: Furosemide 40mg tablets (PRN) Recommendation: Discontinue the PRN furosemide order to prevent dehydration, electrolyte derangements, and dosing ambiguity. & Correct &  \\
\hline
scot\_09\_04 & 9 & Target Medication: Omeprazole 20mg capsules Recommendation: Discontinue omeprazole if there is no confirmed history of peptic ulcer disease or upper GI bleed, as nausea was related to tablet burden and formulation intolerance. & Correct &  \\
\hline
scot\_09\_04 & 10 & Target Medication: Senna 7.5mg tablets Recommendation: Keep the medication in the regimen. & Incorrect &  \\
\hline
scot\_09\_04 & 11 & Target Medication: Laxido sachets Recommendation: Keep the medication in the regimen. & Incorrect &  \\
\hline
scot\_09\_04 & 12 & Target Medication: Fexofenadine 120mg tablets Recommendation: Keep the medication in the regimen. & Acceptable &  \\
\hline
scot\_09\_05 & 1 & Target Medication: Phenytoin 30mg/5ml solution Recommendation: Check urgent total and free serum phenytoin levels and immediately withhold/taper off the medication under specialist guidance, as it is non-standard in LGS and a primary suspect for neurotoxicity and paradoxical seizure exacerbation. & Correct &  \\
\hline
scot\_09\_05 & 2 & Target Medication: Metoclopramide 5mg/5ml solution Recommendation: Discontinue metoclopramide due to the heightened risk of extrapyramidal adverse effects with prolonged use and its potential to lower the seizure threshold. & Correct &  \\
\hline
scot\_09\_05 & 3 & Target Medication: Sodium valproate (Epilim®) 200mg/5ml liquid Recommendation: Check serum valproate and plasma ammonia levels to rule out valproate-induced hyperammonemic encephalopathy; temporarily reduce or adjust dose guided by levels and clinical response. & Acceptable &  \\
\hline
scot\_09\_05 & 4 & Target Medication: Lamotrigine 100mg dispersible tablets Recommendation: Check serum lamotrigine concentration to assess for toxicity resulting from valproate co-administration, and adjust or simplify dosing in consultation with neurology. & Acceptable &  \\
\hline
scot\_09\_05 & 5 & Target Medication: Lamotrigine 25mg dispersible tablets Recommendation: Consolidate tablet strengths once an optimal maintenance dose is re-established to reduce gastrostomy tube administration burden. & Correct &  \\
\hline
scot\_09\_05 & 6 & Target Medication: Clobazam 5mg/5ml suspension Recommendation: Keep the medication in the regimen. & Correct &  \\
\hline
scot\_09\_05 & 7 & Target Medication: Rufinamide 200mg/5ml suspension Recommendation: Keep the medication in the regimen. & Debatable &  \\
\hline
scot\_09\_05 & 8 & Target Medication: Mirtazapine 15mg/5ml solution Recommendation: Temporarily withhold or reduce dose to 15mg at night to minimize additive CNS depression while acute encephalopathy/toxicity is investigated. & Acceptable &  \\
\hline
scot\_09\_06 & 1 & Target Medication: Diclofenac 50mg tablets Recommendation: Discontinue medication due to contraindication in ischemic heart disease and contribution to dyspepsia, fluid retention, and renal impairment. & Correct &  \\
\hline
scot\_09\_06 & 2 & Target Medication: Pioglitazone 30mg tablets Recommendation: Discontinue medication due to significant peripheral edema and fluid retention. & Correct &  \\
\hline
scot\_09\_06 & 3 & Target Medication: Aspirin 75mg dispersible tablets Recommendation: Discontinue medication; 12 months post-NSTEMI has elapsed and patient requires full anticoagulation for atrial fibrillation. & Correct &  \\
\hline
scot\_09\_06 & 4 & Target Medication: Apixaban 5mg tablets Recommendation: Initiate apixaban 5mg orally twice daily for stroke prevention in atrial fibrillation (CHA₂DS₂-VASc = 5). & Correct &  \\
\hline
scot\_09\_06 & 5 & Target Medication: Atorvastatin 80mg tablets Recommendation: Initiate atorvastatin 80mg orally once daily at bedtime for secondary prevention of coronary heart disease. & Correct &  \\
\hline
scot\_09\_06 & 6 & Target Medication: Empagliflozin 10mg tablets Recommendation: Initiate empagliflozin 10mg orally once daily to improve glycemic control and provide cardiovascular and renal protection. & Correct &  \\
\hline
scot\_09\_06 & 7 & Target Medication: Metformin 500mg tablets Recommendation: Reduce dose to 1,000mg twice daily with meals (maximum 2,000mg daily) in view of renal function (eGFR 55 mL/min/1.73m²). & Correct &  \\
\hline
scot\_09\_06 & 8 & Target Medication: Lisinopril 10mg tablets Recommendation: Discontinue the 10mg tablet and consolidate with the 20mg tablet into a single daily formulation. & Correct &  \\
\hline
scot\_09\_06 & 9 & Target Medication: Lisinopril 20mg tablets Recommendation: Increase dose to 30mg orally once daily (single tablet) to simplify regimen and assist blood pressure control. & Correct &  \\
\hline
scot\_09\_06 & 10 & Target Medication: Atenolol 50mg tablets Recommendation: Switch atenolol to bisoprolol 2.5mg orally once daily (titrating according to heart rate and blood pressure) for optimized rate control in AF and post-MI management. & Correct &  \\
\hline
scot\_09\_06 & 11 & Target Medication: Bisoprolol tablets Recommendation: Initiate bisoprolol 2.5mg orally once daily to replace atenolol, titrating as tolerated for rate control. & Correct & Recomendação duplicada? \\
\hline
scot\_09\_06 & 12 & Target Medication: Bendroflumethiazide 2.5mg tablets Recommendation: Discontinue medication; re-evaluate fluid status and blood pressure once NSAID, amlodipine/pioglitazone are addressed. & Correct &  \\
\hline
scot\_09\_06 & 13 & Target Medication: Furosemide 40mg tablets Recommendation: Taper and discontinue as ankle swelling is likely medication-induced (pioglitazone, amlodipine, NSAID) and LV function is normal. & Correct &  \\
\hline
scot\_09\_06 & 14 & Target Medication: Amlodipine 10mg tablets Recommendation: Reduce dose to 5mg daily (or consider switching to an alternative agent if ankle edema persists post-discontinuation of pioglitazone and diclofenac). & Debatable & Suspensão é o correto \\
\hline
scot\_09\_06 & 15 & Target Medication: Levothyroxine 50micrograms/5ml liquid Recommendation: Discontinue liquid presentation. & Correct &  \\
\hline
scot\_09\_06 & 16 & Target Medication: Levothyroxine 25micrograms/5ml liquid Recommendation: Discontinue liquid presentation. & Correct &  \\
\hline
scot\_09\_06 & 17 & Target Medication: Levothyroxine 75micrograms tablets Recommendation: Initiate levothyroxine 75micrograms tablets, one tablet once daily in the morning on an empty stomach. & Correct &  \\
\hline
scot\_09\_06 & 18 & Target Medication: Omeprazole 40mg capsules Recommendation: Reduce dose to 20mg daily or switch to famotidine 20mg twice daily to minimize interaction with citalopram. & Correct &  \\
\hline
scot\_09\_06 & 19 & Target Medication: Citalopram 20mg tablets Recommendation: Keep the medication in the regimen. & Acceptable &  \\
\hline
scot\_09\_06 & 20 & Target Medication: Clenil® 100mcg metered dose inhaler Recommendation: Keep the medication in the regimen. & Debatable &  \\
\hline
scot\_09\_06 & 21 & Target Medication: Co-codamol 8/500mg tablets Recommendation: Keep the medication in the regimen. & Incorrect &  \\
\hline
scot\_09\_06 & 22 & Target Medication: Gabapentin 400mg capsules Recommendation: Keep the medication in the regimen. & Incorrect &  \\
\hline
scot\_09\_06 & 23 & Target Medication: Gliclazide 80mg tablets Recommendation: Keep the medication in the regimen. & Debatable &  \\
\hline
scot\_09\_06 & 24 & Target Medication: Salbutamol 100micrograms metered dose inhaler Recommendation: Keep the medication in the regimen. & Correct &  \\
\hline
scot\_09\_07 & 1 & Target Medication: Edoxaban Recommendation: Reduce dose to 30mg orally once daily due to renal impairment (CrCl 27mL/min) and low body weight (52kg). & Correct &  \\
\hline
scot\_09\_07 & 2 & Target Medication: Clopidogrel Recommendation: Discontinue clopidogrel to reduce the risk of major bleeding, as edoxaban monotherapy is sufficient for stroke prevention in AF and stable IHD. & Correct &  \\
\hline
scot\_09\_07 & 3 & Target Medication: Citalopram Recommendation: Reduce total daily dose to 20mg orally once daily (discontinue the 10mg tablet) to align with maximum dosing guidelines for patients over 65 years and minimize QT prolongation risk. & Acceptable &  \\
\hline
scot\_09\_07 & 4 & Target Medication: Bendroflumethiazide Recommendation: Discontinue bendroflumethiazide due to reduced efficacy at CrCl <30mL/min, low diastolic blood pressure readings, and fall risk. & Correct &  \\
\hline
scot\_09\_07 & 5 & Target Medication: Zopiclone Recommendation: Formulate a deprescribing plan to gradually taper and discontinue zopiclone to reduce sedation, confusion, and fall risk. & Correct &  \\
\hline
scot\_09\_07 & 6 & Target Medication: Co-codamol 30/500 Recommendation: Review ongoing pain management; discontinue or restrict strictly to short-term PRN use, avoiding routine use due to fall, confusion, and constipation risks. & Acceptable & Melhor seria suspender \\
\hline
scot\_09\_07 & 7 & Target Medication: Levothyroxine Recommendation: Consolidate the two separate strengths (100mcg and 25mcg) into a single 125mcg daily dose or explore alternate dosing formulations to reduce pill count. & Correct &  \\
\hline
scot\_09\_07 & 8 & Target Medication: Lisinopril Recommendation: Keep the medication in the regimen. & Debatable &  \\
\hline
scot\_09\_07 & 9 & Target Medication: Atorvastatin Recommendation: Keep the medication in the regimen. & Correct &  \\
\hline
scot\_09\_07 & 10 & Target Medication: Omeprazole Recommendation: Keep the medication in the regimen. & Acceptable &  \\
\hline
scot\_09\_08 & 1 & Target Medication: Longtec® (Oxycodone modified-release) Recommendation: Initiate a slow, patient-agreed gradual taper (e.g., reduce daily dose by \textasciitilde{}10\% every 2–4 weeks) with close monitoring for withdrawal and pain levels. & Correct &  \\
\hline
scot\_09\_08 & 2 & Target Medication: Shortec® (Oxycodone immediate-release) Recommendation: Gradually reduce frequency and limit use to severe breakthrough pain during the Longtec taper, with the ultimate goal of discontinuation. & Correct &  \\
\hline
scot\_09\_08 & 3 & Target Medication: Diazepam Recommendation: Gradually taper dose to discontinuation (e.g., reduce by 0.5 mg to 1 mg every 2–4 weeks) to decrease sedative burden and fall risk. & Correct &  \\
\hline
scot\_09\_08 & 4 & Target Medication: Zopiclone Recommendation: Taper and discontinue slowly to reduce hypnotic burden, confusion, and fall risk; reinforce sleep hygiene. & Correct &  \\
\hline
scot\_09\_08 & 5 & Target Medication: Oxybutynin Recommendation: Trial deprescribing/discontinuation to reduce anticholinergic cognitive burden; re-evaluate urinary symptoms. & Correct &  \\
\hline
scot\_09\_08 & 6 & Target Medication: Nitrofurantoin Recommendation: Discontinue long-term prophylaxis due to reduced renal clearance (CrCl \textasciitilde{}33 mL/min) and risk of long-term toxicities; manage UTIs reactively. & Correct &  \\
\hline
scot\_09\_08 & 7 & Target Medication: Quinine sulphate Recommendation: Discontinue medication after reviewing indication for leg cramps, given risk of toxicity and lack of sustained efficacy. & Correct &  \\
\hline
scot\_09\_08 & 8 & Target Medication: Bisphosphonate (Oral or IV) Recommendation: Re-evaluate and re-initiate osteoporosis treatment (e.g., oral bisphosphonate or denosumab, subject to renal function and DXA review). & Correct &  \\
\hline
scot\_09\_08 & 9 & Target Medication: Amlodipine Recommendation: Keep the medication in the regimen. & Correct &  \\
\hline
scot\_09\_08 & 10 & Target Medication: Clopidogrel Recommendation: Keep the medication in the regimen. & Debatable &  \\
\hline
scot\_09\_08 & 11 & Target Medication: Mirtazapine Recommendation: Keep the medication in the regimen. & Acceptable &  \\
\hline
scot\_09\_08 & 12 & Target Medication: Omeprazole Recommendation: Keep the medication in the regimen. & Correct &  \\
\hline
scot\_09\_08 & 13 & Target Medication: Paracetamol Recommendation: Keep the medication in the regimen. & Acceptable &  \\
\hline
scot\_09\_08 & 14 & Target Medication: Ramipril Recommendation: Keep the medication in the regimen. & Correct &  \\
\hline
scot\_09\_08 & 15 & Target Medication: Senna Recommendation: Keep the medication in the regimen. & Acceptable &  \\
\hline
scot\_09\_08 & 16 & Target Medication: Vitamin D3 Recommendation: Keep the medication in the regimen. & Acceptable &  \\
\hline
scot\_09\_09 & 1 & Target Medication: Femoston-conti 1mg/5mg tablets Recommendation: Keep the medication in the regimen. & Correct &  \\
\hline
scot\_09\_09 & 2 & Target Medication: Paracetamol 500mg caplets Recommendation: Keep the medication in the regimen, encouraging regular or PRN use (up to 1g four times daily) as a non-sedating baseline analgesic during the weaning period. & Acceptable &  \\
\hline
scot\_09\_09 & 3 & Target Medication: Tramadol 50mg capsules Recommendation: Initiate a gradual, individualized taper (e.g., reduce daily dose by 50–100mg every 1 to 2 weeks, for instance reducing to 50mg four times daily, then tapering further) with the goal of complete discontinuation to avoid withdrawal symptoms and alleviate sedation. & Correct &  \\
\hline
scot\_09\_09 & 4 & Target Medication: Gabapentin 300mg capsules Recommendation: Maintain current dose while tramadol is being actively tapered, then initiate a slow, stepwise reduction (e.g., reduce by 300mg every 1 to 2 weeks) toward discontinuation or the lowest effective dose once opioid cessation is achieved. & Acceptable &  \\
\hline
scot\_09\_10 & 1 & Target Medication: Tramadol 50mg capsules Recommendation: Taper and discontinue gradually to avoid opioid withdrawal and minimize the risk of serotonin syndrome. & Correct &  \\
\hline
scot\_09\_10 & 2 & Target Medication: Amitriptyline 50mg tablets Recommendation: Reduce dose to 25mg at night with a plan to taper and discontinue to reduce anticholinergic load and serotonergic drug interactions. & Correct &  \\
\hline
scot\_09\_10 & 3 & Target Medication: Duloxetine 60mg capsules Recommendation: Keep the medication in the regimen. & Acceptable &  \\
\hline
scot\_09\_10 & 4 & Target Medication: Pregabalin 75mg capsules Recommendation: Initiate at 75mg twice daily (or 75mg at night initially) and titrate according to response and tolerability for diabetic peripheral neuropathic pain. & Acceptable &  \\
\hline
scot\_09\_10 & 5 & Target Medication: Metformin 500mg tablets Recommendation: Keep the medication in the regimen. & Correct &  \\
\hline
scot\_09\_10 & 6 & Target Medication: Atorvastatin 10mg tablets Recommendation: Keep the medication in the regimen. & Incorrect &  \\
\hline
scot\_09\_10 & 7 & Target Medication: Allopurinol 100mg tablets Recommendation: Keep the medication in the regimen. & Acceptable &  \\
\hline
scot\_09\_10 & 8 & kept | cannabis & Correct &  \\
\hline
scot\_09\_11 & 1 & Target Medication: Doxazosin Recommendation: Discontinue doxazosin to reduce the risk of orthostatic hypotension and falls. & Correct &  \\
\hline
scot\_09\_11 & 2 & Target Medication: Furosemide Recommendation: Reduce dose to 40 mg once daily in the morning and review ongoing indication for loop diuretic therapy. & Acceptable &  \\
\hline
scot\_09\_11 & 3 & Target Medication: Folic acid Recommendation: Discontinue folic acid as serum folate levels are elevated (>20 ng/mL) and FBC is normal. & Correct &  \\
\hline
scot\_09\_11 & 4 & Target Medication: Ascorbic acid Recommendation: Discontinue ascorbic acid as there is no clear clinical indication, reducing pill burden. & Correct &  \\
\hline
scot\_09\_11 & 5 & Target Medication: Solifenacin Recommendation: Deprescribe and discontinue solifenacin to reduce anticholinergic burden and fall risk. & Correct &  \\
\hline
scot\_09\_11 & 6 & Target Medication: Venlafaxine Recommendation: Gradually taper dose downwards (e.g., reduce to 75 mg twice daily, then taper further) with an aim to rationalise antidepressant therapy and decrease fall/sedation risk. & Correct &  \\
\hline
scot\_09\_11 & 7 & Target Medication: Mirtazapine Recommendation: Reduce dose from 45 mg to 15–30 mg at bedtime to minimise excessive CNS depression, orthostasis, and fall risk. & Correct &  \\
\hline
scot\_09\_11 & 8 & Target Medication: Oxycodone Recommendation: Review pain indication and initiate a structured downward titration to deprescribe if chronic pain is stable. & Correct &  \\
\hline
scot\_09\_11 & 9 & Target Medication: Pregabalin Recommendation: Review indication and consider tapering to discontinue to reduce dizziness, sedation, and fall risk. & Correct &  \\
\hline
scot\_09\_11 & 10 & Target Medication: Insulin lispro Recommendation: Deprescribe short-acting insulin lispro in coordination with the diabetes specialist team to relax glycaemic targets and minimise hypoglycaemic fall risk. & Acceptable &  \\
\hline
scot\_09\_11 & 11 & Target Medication: Insulin glargine Recommendation: Titrate and simplify basal insulin therapy, targeting a less stringent HbA1c (58–64 mmol/mol) in frailty. & Acceptable &  \\
\hline
scot\_09\_11 & 12 & Target Medication: Liraglutide Recommendation: Keep the medication in the regimen. & Acceptable &  \\
\hline
scot\_09\_11 & 13 & Target Medication: Metformin Recommendation: Keep the medication in the regimen. & Incorrect &  \\
\hline
scot\_09\_11 & 14 & Target Medication: Ferrous fumarate Recommendation: Review ongoing indication; discontinue if ferritin/iron stores are normal and FBC remains stable to reduce GI side effects and pill burden. & Correct &  \\
\hline
scot\_09\_11 & 15 & Target Medication: Lactulose Recommendation: Keep the medication in the regimen. & Acceptable &  \\
\hline
scot\_09\_11 & 16 & Target Medication: Levothyroxine sodium Recommendation: Keep the medication in the regimen. & Correct &  \\
\hline
scot\_09\_11 & 17 & Target Medication: Carbocisteine Recommendation: Keep the medication in the regimen. & Debatable &  \\
\hline
scot\_09\_11 & 18 & Target Medication: Trelegy Ellipta Recommendation: Keep the medication in the regimen. & Correct &  \\
\hline
scot\_09\_11 & 19 & Target Medication: Salbutamol inhaler CFC-free Recommendation: Keep the medication in the regimen. & Correct &  \\
\hline
scot\_09\_11 & 20 & Target Medication: Salbutamol nebuliser liquid Recommendation: Keep the medication in the regimen. & Debatable &  \\
\hline
scot\_09\_11 & 21 & Target Medication: Glyceryl trinitrate sublingual spray Recommendation: Keep the medication in the regimen. & Acceptable &  \\
\hline
scot\_09\_11 & 22 & Target Medication: Blood glucose testing strips Recommendation: Keep the medication in the regimen. & Correct &  \\
\hline
scot\_09\_11 & 23 & Target Medication: Lancets Recommendation: Keep the medication in the regimen. & Correct &  \\
\hline
scot\_09\_11 & 24 & Target Medication: Insulin needles Recommendation: Keep the medication in the regimen. & Correct &  \\
\hline
scot\_09\_11 & 25 & Target Medication: Aspirin Recommendation: Initiate aspirin 75 mg once daily for secondary prevention of ischaemic heart disease and previous CABG/MI, taking into account bleeding versus thrombotic risk in frailty. & Debatable &  \\
\hline
scot\_09\_11 & 26 & Target Medication: Atorvastatin Recommendation: Initiate atorvastatin 20 mg once daily for secondary cardiovascular prevention, subject to shared decision-making regarding goals of care. & Acceptable &  \\
\hline
scot\_09\_12 & 1 & Target Medication: Oxypro 40mg modified release tablets Recommendation: Gradually taper dose under close clinical supervision with the goal of significant reduction or discontinuation to mitigate respiratory depression risk and reduce constipation. & Correct &  \\
\hline
scot\_09\_12 & 2 & Target Medication: Shortec 10mg capsules Recommendation: Restrict to strictly defined breakthrough pain during the tapering phase, then discontinue. & Correct &  \\
\hline
scot\_09\_12 & 3 & Target Medication: Aspirin 75mg enteric coated tablets Recommendation: Discontinue to reduce gastrointestinal and systemic bleeding risk, as there is no compelling indication for concurrent therapy with apixaban. & Correct &  \\
\hline
scot\_09\_12 & 4 & Target Medication: Ferrous fumarate 210mg tablets Recommendation: Discontinue, as FBC is currently within normal limits and the medication is exacerbating chronic constipation. & Correct &  \\
\hline
scot\_09\_12 & 5 & Target Medication: Gabapentin 100mg capsules Recommendation: Review indication and consider tapering or discontinuing to reduce combined CNS/respiratory depression with opioids. & Correct &  \\
\hline
scot\_09\_12 & 6 & Target Medication: Senna 7.5mg tablets Recommendation: Switch to or add an osmotic laxative (e.g., macrogol sachets) for opioid-induced constipation, titrating down as opioids are reduced. & Correct &  \\
\hline
scot\_09\_12 & 7 & Target Medication: Solifenacin 5mg tablets Recommendation: Review efficacy and consider a trial withdrawal to alleviate anticholinergic burden and constipation. & Correct &  \\
\hline
scot\_09\_12 & 8 & Target Medication: Amiodarone 200mg tablets Recommendation: Refer to cardiology to review the ongoing indication for rhythm control versus rate control with bisoprolol alone, ensuring regular TFT, LFT, and pulmonary monitoring while continued. & Correct &  \\
\hline
scot\_09\_12 & 9 & Target Medication: Atorvastatin 20mg tablets Recommendation: Initiate 20mg once daily for primary cardiovascular risk reduction in type 2 diabetes and hypertension. & Correct &  \\
\hline
scot\_09\_12 & 10 & Target Medication: Apixaban 5mg tablets Recommendation: Keep the medication in the regimen. & Correct &  \\
\hline
scot\_09\_12 & 11 & Target Medication: Bisoprolol 10mg tablets Recommendation: Keep the medication in the regimen. & Correct &  \\
\hline
scot\_09\_12 & 12 & Target Medication: Bumetanide 1mg tablets Recommendation: Keep the medication in the regimen. & Acceptable &  \\
\hline
scot\_09\_12 & 13 & Target Medication: Salbutamol 100microgram metered dose inhaler Recommendation: Keep the medication in the regimen. & Correct &  \\
\hline
scot\_09\_12 & 14 & Target Medication: Spironolactone 25mg tablets Recommendation: Keep the medication in the regimen. & Correct &  \\
\hline
scot\_09\_12 & 15 & Target Medication: Trelegy dry powder inhaler 92/55/22micrograms/dose Recommendation: Keep the medication in the regimen. & Correct &  \\
\hline
scot\_09\_13 & 1 & Target Medication: Co-beneldopa 25mg/100mg tablets Recommendation: Switch to Co-beneldopa 25mg/100mg dispersible tablets one tablet four times daily to facilitate administration and improve morning absorption. & Correct &  \\
\hline
scot\_09\_13 & 2 & Target Medication: Co-beneldopa 12.5mg/50mg dispersible tablet Recommendation: Keep the medication in the regimen. & Correct &  \\
\hline
scot\_09\_13 & 3 & Target Medication: Co-codamol 30mg/500 tablets Recommendation: Reduce dose to one tablet four times daily with a view to tapering down further or switching to regular Paracetamol 500mg tablets/soluble tablets (two tablets QDS) to reduce sedation, constipation, and fall risk. & Correct &  \\
\hline
scot\_09\_13 & 4 & Target Medication: Pregabalin 25mg capsules Recommendation: Review indication for neuropathic pain and consider a gradual deprescribing trial to alleviate daytime somnolence if pain is controlled. & Correct &  \\
\hline
scot\_09\_13 & 5 & Target Medication: Colecalciferol 400 unit tablet Recommendation: Increase dose to 800–1,000 units daily (e.g., Colecalciferol 800 unit tablet/capsule or liquid formulation once daily) for fall and fracture prevention in a housebound, frail adult. & Correct &  \\
\hline
scot\_09\_13 & 6 & Target Medication: Sertraline 100mg tablets Recommendation: Keep the medication in the regimen. & Debatable &  \\
\hline
scot\_09\_13 & 7 & Target Medication: Laxido sachets Recommendation: Keep the medication in the regimen. & Acceptable &  \\
\hline
scot\_09\_13 & 8 & Target Medication: Peptac suspension Recommendation: Keep the medication in the regimen. & Debatable &  \\
\hline
scot\_09\_14 & 1 & Target Medication: Citalopram 20mg tablets Recommendation: Increase dose to 30mg once daily for 1 to 2 weeks, then increase to 40mg once daily if tolerated; issue short supply (e.g., 14 days at a time) due to overdose risk. & Debatable &  \\
\hline
scot\_09\_14 & 2 & Target Medication: Folic acid 5mg tablets Recommendation: Initiate 5mg orally once daily for 4 months to treat documented folate deficiency. & Acceptable &  \\
\hline
scot\_09\_14 & 3 & Target Medication: Clenil (beclomethasone dipropionate) 100microgram meter dose inhaler Recommendation: Keep the medication in the regimen. & Correct &  \\
\hline
scot\_09\_14 & 4 & Target Medication: Salbutamol 100mcg meter dose inhaler Recommendation: Keep the medication in the regimen. & Correct &  \\
\hline
scot\_09\_15 & 1 & Target Medication: None Recommendation: Do not initiate pharmacotherapy at this time in alignment with patient preference and clinical guidelines for moderate anxiety. & Correct &  \\
\hline
scot\_09\_16 & 1 & Target Medication: Solpadeine Max® (Paracetamol 500mg / Codeine 12.8mg) Recommendation: Discontinue entirely; stop purchasing over-the-counter and provide a structured detox plan to avoid paracetamol-induced hepatotoxicity and opioid dependence. & Correct &  \\
\hline
scot\_09\_16 & 2 & Target Medication: Street benzodiazepines Recommendation: Discontinue illicit supply under medical supervision with a planned, gradual substitution and tapering regimen (e.g., titrated diazepam reduction schedule) via specialized addiction services. & Correct &  \\
\hline
scot\_09\_16 & 3 & Target Medication: Paracetamol Recommendation: Initiate paracetamol 500mg to 1000mg orally up to four times daily as needed for chronic ankle pain, not to exceed a reduced daily maximum (e.g., 2g to 3g daily) given regular alcohol consumption. & Correct &  \\
\hline
scot\_09\_16 & 4 & Target Medication: Sertraline Recommendation: Initiate sertraline 50mg orally once daily, titrating according to response up to 150–200mg daily, to treat underlying PTSD, anxiety, and depressive symptoms. & Correct &  \\
\hline
scot\_09\_16 & 5 & Target Medication: Take-home Naloxone Recommendation: Initiate and supply a take-home naloxone kit (e.g., intranasal spray 1.8mg/0.1mL or intramuscular injection) with comprehensive overdose response training. & Correct &  \\
\hline
scot\_09\_17 & 1 & Target Medication: Lithium (Priadel®) 400mg tablets Recommendation: Reduce dose to Priadel® 600mg (e.g., one 400mg tablet and one 200mg tablet) once daily at night. & Correct &  \\
\hline
scot\_09\_17 & 2 & Target Medication: Quinine sulfate 200mg tablets Recommendation: Deprescribe and discontinue with a trial withdrawal, switching to non-pharmacological management for leg cramps. & Correct &  \\
\hline
scot\_09\_17 & 3 & Target Medication: Ramipril 10mg capsules Recommendation: Keep the medication in the regimen. & Acceptable &  \\
\hline
scot\_09\_17 & 4 & Target Medication: Atorvastatin 20mg tablets Recommendation: Keep the medication in the regimen. & Correct &  \\
\hline
scot\_09\_17 & 5 & kept | ibuprofen & Incorrect &  \\
\hline
scot\_09\_18 & 1 & Target Medication: Hydroxyzine Recommendation: Discontinue immediately due to high anticholinergic burden, fall risk, QTc prolongation risk, and exacerbation of delirium. & Correct &  \\
\hline
scot\_09\_18 & 2 & Target Medication: Solifenacin Recommendation: Discontinue immediately due to anticholinergic adverse effects, antagonism of donepezil, and contribution to delirium, falls, and constipation. & Correct &  \\
\hline
scot\_09\_18 & 3 & Target Medication: Risperidone Recommendation: Discontinue and do not administer; avoid in dementia-related delirium until reversible causes are managed and non-pharmacological de-escalation measures have been exhausted. & Correct &  \\
\hline
scot\_09\_18 & 4 & Target Medication: Donepezil Recommendation: Withhold temporarily due to symptomatic bradycardia (HR 50 bpm) and reassess continuation/dose once pulse rate normalizes and delirium resolves. & Correct &  \\
\hline
scot\_09\_18 & 5 & Target Medication: Hydroxocobalamin Recommendation: Initiate 1 mg intramuscularly on alternate days for up to two weeks (or per local protocol for neurological involvement), followed by maintenance dosing to treat severe vitamin B12 deficiency. & Correct &  \\
\hline
scot\_09\_18 & 6 & Target Medication: Folic acid Recommendation: Initiate 5 mg orally once daily to treat folate deficiency, ensuring vitamin B12 therapy has been commenced. & Correct &  \\
\hline
scot\_09\_18 & 7 & Target Medication: Senna Recommendation: Initiate 15 mg orally at night as needed to treat constipation. & Acceptable &  \\
\hline
scot\_09\_18 & 8 & Target Medication: Aspirin Recommendation: Keep the medication in the regimen. & Debatable &  \\
\hline
scot\_09\_18 & 9 & Target Medication: Atorvastatin Recommendation: Keep the medication in the regimen. & Incorrect &  \\
\hline
scot\_09\_18 & 10 & Target Medication: Perindopril Recommendation: Keep the medication in the regimen. & Incorrect &  \\
\hline
scot\_09\_19 & 1 & Target Medication: Risperidone 500 microgram tablets Recommendation: Deprescribe and taper off over 1 to 2 weeks (e.g., reduce to 250 micrograms once daily for a few days, then discontinue), as the acute infection and acute delirium have resolved. & Correct &  \\
\hline
scot\_09\_19 & 2 & Target Medication: Amitriptyline 50mg tablets Recommendation: Taper and discontinue gradually (e.g., reduce to 25mg at night for 1–2 weeks, then 10mg, then stop) to reduce anticholinergic burden and sedation; explore alternative non-pharmacological measures or safer alternatives if prescribed for neuropathic pain or sleep. & Correct &  \\
\hline
scot\_09\_19 & 3 & Target Medication: Oxybutynin 5mg tablets Recommendation: Discontinue to minimize anticholinergic burden and cognitive impairment; re-evaluate urinary symptoms post-cessation and consider non-pharmacological bladder training or a uroselective agent (e.g., mirabegron) if urgency recurs. & Correct &  \\
\hline
scot\_09\_19 & 4 & Target Medication: Amlodipine 5mg tablets Recommendation: Re-initiate at 5mg once daily to manage uncontrolled systolic hypertension (BP 176/84 mmHg), with blood pressure monitoring. & Correct &  \\
\hline
scot\_09\_19 & 5 & Target Medication: Bendroflumethiazide 2.5mg tablets Recommendation: Withhold at present while amlodipine is re-established and fluid/nutritional intake stabilizes, to avoid volume depletion, hyponatraemia, and falls. & Correct &  \\
\hline
scot\_09\_19 & 6 & Target Medication: Metformin 500mg tablets Recommendation: Keep the medication in the regimen. & Correct &  \\
\hline
scot\_09\_19 & 7 & Target Medication: Amoxicillin 500mg capsules Recommendation: Ensure medication is confirmed discontinued and removed from the active regimen as the course is complete. & Correct &  \\
\hline
scot\_09\_19 & 8 & Target Medication: Trimethoprim 200mg tablets Recommendation: Ensure medication is confirmed discontinued and removed from the active regimen as the course is complete. & Correct &  \\
\hline
scot\_09\_20 & 1 & Target Medication: Amoxicillin Recommendation: Initiate amoxicillin oral capsules 500 mg every 12 hours (dose-adjusted for eCrCl <30 mL/min) to complete a 7- to 14-day total course for upper UTI, after updating allergy status. & Correct &  \\
\hline
scot\_09\_20 & 2 & Target Medication: Furosemide Recommendation: Increase dose to 40 mg orally once daily in the morning and monitor fluid balance, weight, renal function, and electrolytes. & Correct &  \\
\hline
scot\_09\_20 & 3 & Target Medication: Lisinopril Recommendation: Continue to withhold until blood pressure, renal function, and acute delirium stabilize. & Acceptable & correto seria suspender \\
\hline
scot\_09\_20 & 4 & Target Medication: Atorvastatin Recommendation: Discontinue permanently due to lack of clear indication, advanced age, low BMI, and frailty. & Correct &  \\
\hline
scot\_09\_20 & 5 & Target Medication: Donepezil Recommendation: Keep the medication in the regimen. & Correct &  \\
\hline
scot\_09\_20 & 6 & Target Medication: Paracetamol Recommendation: Keep the medication in the regimen. & Incorrect &  \\
\hline
scot\_09\_20 & 7 & Target Medication: Trelegy Ellipta (fluticasone/vilanterol/umeclidinium) Recommendation: Keep the medication in the regimen. & Debatable &  \\
\hline
scot\_09\_20 & 8 & Target Medication: Salbutamol inhaler Recommendation: Keep the medication in the regimen. & Acceptable &  \\
\hline
scot\_09\_20 & 9 & Target Medication: Macrogol 3350 Recommendation: Keep the medication in the regimen. & Acceptable &  \\
\hline
scot\_09\_20 & 10 & Target Medication: Senna Recommendation: Keep the medication in the regimen. & Acceptable &  \\
\hline
scot\_09\_21 & 1 & Target Medication: Amlodipine Recommendation: Discontinue the medication due to hypotension and transition to comfort-focused care. & Correct &  \\
\hline
scot\_09\_21 & 2 & Target Medication: Aspirin Recommendation: Discontinue the medication as it is non-essential and oral route is compromised. & Correct &  \\
\hline
scot\_09\_21 & 3 & Target Medication: Co-beneldopa Recommendation: Discontinue oral capsules and convert to a transdermal rotigotine patch to prevent acute dopamine withdrawal. & Correct &  \\
\hline
scot\_09\_21 & 4 & Target Medication: Macrogol 3350 Recommendation: Discontinue the medication due to current loose stools and swallowing impairment. & Correct &  \\
\hline
scot\_09\_21 & 5 & Target Medication: Paracetamol Recommendation: Discontinue oral caplets due to dysphagia; use subcutaneous analgesia as needed for comfort. & Correct &  \\
\hline
scot\_09\_21 & 6 & Target Medication: Quinine sulfate Recommendation: Discontinue the medication due to unfavourable risk-benefit profile and lack of acute indication. & Correct &  \\
\hline
scot\_09\_21 & 7 & Target Medication: Senna Recommendation: Discontinue the medication due to active loose stools and swallowing impairment. & Correct &  \\
\hline
scot\_09\_21 & 8 & Target Medication: Simvastatin Recommendation: Discontinue the medication as preventative therapy is no longer appropriate. & Correct &  \\
\hline
scot\_09\_21 & 9 & Target Medication: Thiamine Recommendation: Discontinue the medication as it is non-essential in this acute/comfort-care setting. & Correct &  \\
\hline
scot\_09\_21 & 10 & Target Medication: Rotigotine transdermal patch Recommendation: Initiate rotigotine patch 2 mg/24 hours applied topically once daily, titrating as needed to cover baseline dopaminergic requirements. & Correct &  \\
\hline
scot\_09\_21 & 11 & Target Medication: Midazolam injection Recommendation: Prescribe midazolam 2.5 mg to 5 mg subcutaneously PRN (hourly as needed) for terminal agitation and restlessness. & Acceptable &  \\
\hline
scot\_09\_21 & 12 & Target Medication: Oxycodone injection Recommendation: Prescribe oxycodone 1.25 mg to 2.5 mg subcutaneously PRN (every 2 to 4 hours as needed) for pain or dyspnoea, preferred over morphine due to renal impairment (CrCl 21 mL/min). & Acceptable &  \\
\hline
scot\_09\_21 & 13 & Target Medication: Glycopyrronium injection Recommendation: Prescribe glycopyrronium 200 micrograms subcutaneously PRN (every 4 hours as needed) for respiratory secretions. & Acceptable &  \\
\hline
scot\_09\_21 & 14 & Target Medication: Nystatin oral suspension Recommendation: Prescribe nystatin 100,000 units/mL oral suspension, 1 mL applied gently to the oral cavity four times daily for oral candidiasis. & Correct &  \\
\hline
scot\_09\_22 & 1 & Target Medication: Aspirin Recommendation: Initiate aspirin 75mg dispersible/gastro-resistant tablets one tablet daily for secondary prevention of ischaemic heart disease. & Debatable &  \\
\hline
scot\_09\_22 & 2 & Target Medication: Prochlorperazine Recommendation: Discontinue prochlorperazine as the acute episode of nausea and vomiting has resolved. & Correct &  \\
\hline
scot\_09\_22 & 3 & Target Medication: Co-codamol 30mg/500mg tablets Recommendation: Review ongoing pain requirement; consider reducing to paracetamol 500mg–1g PRN/regular to minimize constipation, sedation, and fall risk. & Correct &  \\
\hline
scot\_09\_22 & 4 & Target Medication: Lansoprazole 30mg capsules Recommendation: Reduce dose to 15mg once daily for gastroprotection (co-prescribed with aspirin) and review ongoing indication. & Correct &  \\
\hline
scot\_09\_22 & 5 & Target Medication: Macrogol 3350 sachets Recommendation: Continue while on opioid therapy, but step down/titrate according to bowel frequency and potential opioid dose reduction. & Debatable &  \\
\hline
scot\_09\_22 & 6 & Target Medication: Senna 7.5mg tablets Recommendation: Continue while on opioid therapy, but step down/titrate according to bowel frequency and potential opioid dose reduction. & Debatable &  \\
\hline
scot\_09\_22 & 7 & Target Medication: Atorvastatin 80mg tablets Recommendation: Keep the medication in the regimen. & Acceptable &  \\
\hline
scot\_09\_22 & 8 & Target Medication: Dapagliflozin 10mg tablets Recommendation: Keep the medication in the regimen. & Acceptable &  \\
\hline
scot\_09\_22 & 9 & Target Medication: Finasteride 5mg tablets Recommendation: Keep the medication in the regimen. & Correct &  \\
\hline
scot\_09\_22 & 10 & Target Medication: Furosemide 40mg tablets Recommendation: Keep the medication in the regimen. & Acceptable &  \\
\hline
scot\_09\_22 & 11 & Target Medication: Metformin 500mg tablets Recommendation: Keep the medication in the regimen. & Acceptable &  \\
\hline
scot\_09\_22 & 12 & Target Medication: Mirtazapine 30mg tablets Recommendation: Keep the medication in the regimen. & Correct &  \\
\hline
scot\_09\_22 & 13 & Target Medication: Ramipril 10mg capsules Recommendation: Keep the medication in the regimen. & Acceptable &  \\
\hline
scot\_09\_22 & 14 & Target Medication: Tamsulosin 400 micrograms modified-release capsules Recommendation: Keep the medication in the regimen. & Correct &  \\
\hline
scot\_09\_23 & 1 & Target Medication: Zopiclone 7.5mg tablet Recommendation: Reduce dose to 3.75mg at night with a plan to gradually taper and discontinue to reduce fall risk. & Correct &  \\
\hline
scot\_09\_23 & 2 & Target Medication: Devil’s claw herbal remedy Recommendation: Discontinue use immediately due to potential drug interactions and bleeding risks with warfarin. & Correct &  \\
\hline
scot\_09\_23 & 3 & Target Medication: Diclofenac gel Recommendation: Discontinue as patient reports it is ineffective for knee pain. & Correct &  \\
\hline
scot\_09\_23 & 4 & Target Medication: Paracetamol 500mg tablets Recommendation: Increase to 1g (two tablets) three times daily regularly for osteoarthritis knee pain instead of PRN use. & Correct &  \\
\hline
scot\_09\_23 & 5 & Target Medication: Alendronic acid 70mg tablets Recommendation: Pause therapy pending repeat DXA scan and fracture risk reassessment to consider a bisphosphonate holiday. & Correct &  \\
\hline
scot\_09\_23 & 6 & Target Medication: Warfarin tablets 1mg tablets Recommendation: Keep the medication in the regimen. & Acceptable &  \\
\hline
scot\_09\_23 & 7 & Target Medication: Anoro Ellipta® dry powder inhaler Recommendation: Keep the medication in the regimen. & Acceptable &  \\
\hline
scot\_09\_23 & 8 & Target Medication: Salbutamol metered dose inhaler 100 micrograms Recommendation: Keep the medication in the regimen. & Acceptable &  \\
\hline
scot\_09\_23 & 9 & Target Medication: Atorvastatin 40mg tablets Recommendation: Keep the medication in the regimen. & Correct &  \\
\hline
scot\_09\_23 & 10 & Target Medication: Ramipril 10mg tablet Recommendation: Keep the medication in the regimen. & Acceptable &  \\
\hline
scot\_09\_23 & 11 & Target Medication: Indapamide 2.5mg tablets Recommendation: Keep the medication in the regimen. & Acceptable &  \\
\hline
scot\_09\_23 & 12 & Target Medication: Lansoprazole 15mg capsule Recommendation: Keep the medication in the regimen. & Debatable &  \\
\hline
scot\_09\_23 & 13 & Target Medication: Calcium 1.5g and vitamin D 400iu chewable tablets Recommendation: Keep the medication in the regimen. & Correct &  \\
\hline
scot\_09\_23 & 14 & Target Medication: Peptac® liquid Recommendation: Keep the medication in the regimen. & Incorrect &  \\
\hline
scot\_09\_24 & 1 & Target Medication: Alogliptin 12.5mg tablets Recommendation: Discontinue the medication to reduce pill burden as strict glycemic control is no longer required in palliative care. & Correct &  \\
\hline
scot\_09\_24 & 2 & Target Medication: Amitriptyline 25mg tablets Recommendation: Discontinue the medication (or taper if needed) to reduce anticholinergic burden, dry mouth, constipation, and fall risk. & Correct &  \\
\hline
scot\_09\_24 & 3 & Target Medication: Atorvastatin 20mg tablets Recommendation: Discontinue the medication as primary/secondary lipid prevention offers no benefit in the context of limited life expectancy. & Correct &  \\
\hline
scot\_09\_24 & 4 & Target Medication: Dexamethasone 2mg tablets Recommendation: Keep the medication in the regimen, ensuring the daily dose is taken as a single morning dose to avoid sleep disturbance. & Correct &  \\
\hline
scot\_09\_24 & 5 & Target Medication: Diazepam 5mg tablets Recommendation: Taper and discontinue regular daily dosing due to excessive sedation and fall risk when combined with escalating opioids; consider low-dose PRN if indicated for anxiety. & Correct &  \\
\hline
scot\_09\_24 & 6 & Target Medication: Edoxaban 60mg tablets Recommendation: Keep the medication in the regimen, with ongoing monitoring of renal function and bleeding risk. & Debatable &  \\
\hline
scot\_09\_24 & 7 & Target Medication: Folic acid 5mg tablets Recommendation: Discontinue the medication due to lack of ongoing indication and to reduce pill burden. & Correct &  \\
\hline
scot\_09\_24 & 8 & Target Medication: Levothyroxine 100microgram tablets Recommendation: Keep the medication in the regimen. & Correct &  \\
\hline
scot\_09\_24 & 9 & Target Medication: Lidocaine 700mg plasters Recommendation: Keep the medication in the regimen. & Correct &  \\
\hline
scot\_09\_24 & 10 & Target Medication: Macrogol 3350 sachets Recommendation: Increase the dose to one sachet twice daily to address ongoing constipation. & Correct &  \\
\hline
scot\_09\_24 & 11 & Target Medication: Metformin 500mg tablets Recommendation: Discontinue the medication to reduce pill burden and avoid risk of adverse effects in declining renal function. & Correct &  \\
\hline
scot\_09\_24 & 12 & Target Medication: Morphine MR 30mg capsules Recommendation: Discontinue the oral capsules and rotate to Transdermal Fentanyl (e.g., 25 mcg/hr patch changed every 72 hours) or increase oral morphine dose using liquid formulation to alleviate swallowing difficulty and achieve pain control. & Acceptable &  \\
\hline
scot\_09\_24 & 13 & Target Medication: Morphine IR 10mg tablets Recommendation: Change formulation to Morphine oral solution (10 mg/5 mL) 10–15 mg every 2 to 4 hours as required for breakthrough pain. & Correct &  \\
\hline
scot\_09\_24 & 14 & Target Medication: Omeprazole 20mg capsules Recommendation: Keep the medication in the regimen for gastroprotection while patient remains on systemic dexamethasone. & Debatable &  \\
\hline
scot\_09\_24 & 15 & Target Medication: Paracetamol 500mg tablets Recommendation: Keep the medication in the regimen (can consider soluble or suspension formulations if swallowing tablets is difficult). & Correct &  \\
\hline
scot\_09\_24 & 16 & Target Medication: Pregabalin 300mg capsules Recommendation: Reduce the dose to 100 mg twice daily (or 150 mg twice daily) to account for renal impairment (eGFR 40 mL/min) and minimize excessive sedation. & Correct &  \\
\hline
scot\_09\_24 & 17 & Target Medication: Ramipril 10mg capsules Recommendation: Discontinue the medication given blood pressure of 110/60 mmHg, to prevent hypotension, falls, and pill burden. & Correct &  \\
\hline
scot\_09\_24 & 18 & Target Medication: Salbutamol 100 micrograms pMDI Recommendation: Keep the medication in the regimen. & Acceptable &  \\
\hline
scot\_09\_24 & 19 & Target Medication: Senna 7.5mg tablets Recommendation: Keep the medication in the regimen, titrating up to two tablets twice daily if required alongside increased Macrogol. & Acceptable &  \\
\hline
scot\_09\_24 & 20 & Target Medication: Sertraline 150mg tablets Recommendation: Keep the medication in the regimen. & Correct &  \\
\hline
scot\_09\_24 & 21 & Target Medication: Solifenacin 10mg tablets Recommendation: Discontinue the medication to relieve severe dry mouth and constipation. & Correct &  \\
\hline
scot\_09\_24 & 22 & Target Medication: Trelegy Ellipta® 22/92/55 micrograms DPI Recommendation: Keep the medication in the regimen, ensuring proper mouth rinsing after inhalation. & Debatable &  \\
\hline
scot\_09\_24 & 23 & Target Medication: Nystatin 100,000 units/mL oral suspension Recommendation: Initiate 1 mL four times daily after food for suspected oral candidiasis secondary to inhaled/systemic corticosteroids, following clinical mouth inspection. & Correct &  \\
\hline
scot\_09\_24 & 24 & Target Medication: Artificial saliva spray / gel Recommendation: Initiate as required for symptomatic management of dry mouth. & Correct &  \\
\hline
scot\_09\_25 & 1 & Target Medication: Fentanyl patch 25micrograms/hour Recommendation: Keep the medication in the regimen. & Correct &  \\
\hline
scot\_09\_25 & 2 & Target Medication: Morphine sulfate 10mg/5ml oral solution Recommendation: Discontinue and switch to subcutaneous morphine 2.5mg to 5mg every 2 to 4 hours as required for breakthrough pain and breathlessness. & Correct &  \\
\hline
scot\_09\_25 & 3 & Target Medication: Ibuprofen 400mg tablets Recommendation: Discontinue the medication. & Correct &  \\
\hline
scot\_09\_25 & 4 & Target Medication: Paracetamol 500mg tablets Recommendation: Discontinue the medication. & Acceptable &  \\
\hline
scot\_09\_25 & 5 & Target Medication: Ramipril 10mg capsules Recommendation: Discontinue the medication. & Correct &  \\
\hline
scot\_09\_25 & 6 & Target Medication: Fluoxetine 40mg capsules Recommendation: Discontinue the medication. & Correct &  \\
\hline
scot\_09\_25 & 7 & Target Medication: Zopiclone 7.5mg tablets Recommendation: Discontinue the medication. & Correct &  \\
\hline
scot\_09\_25 & 8 & Target Medication: Laxido sachets Recommendation: Discontinue the medication. & Correct &  \\
\hline
scot\_09\_25 & 9 & Target Medication: Omeprazole 20mg capsules Recommendation: Discontinue the medication. & Correct &  \\
\hline
scot\_09\_25 & 10 & Target Medication: Fluconazole 50mg capsules Recommendation: Discontinue the medication. & Correct &  \\
\hline
scot\_09\_25 & 11 & Target Medication: Metoclopramide 10mg tablets Recommendation: Discontinue oral tablets and switch to metoclopramide 10mg subcutaneously every 8 hours as required (or 30mg/24 hours via continuous subcutaneous infusion). & Correct &  \\
\hline
scot\_09\_25 & 12 & Target Medication: Insulin aspart (NovoRapid) Recommendation: Discontinue prandial insulin injections. & Correct &  \\
\hline
scot\_09\_25 & 13 & Target Medication: Insulin detemir (Levemir) Recommendation: Reduce and consolidate to a simplified, reduced once-daily basal dose (e.g., 50\% to 75\% of baseline basal requirements once daily in the morning) with relaxed glycemic targets (8 to 15 mmol/L) to prevent symptomatic hypoglycemia and ketoacidosis. & Correct &  \\
\hline
scot\_09\_25 & 14 & Target Medication: Midazolam injection Recommendation: Initiate midazolam 2.5mg to 5mg subcutaneously every 2 to 4 hours as required for terminal restlessness, anxiety, or breathlessness. & Correct &  \\
\hline
scot\_09\_25 & 15 & Target Medication: Glycopyrronium injection Recommendation: Initiate glycopyrronium 200micrograms subcutaneously every 4 hours as required for moist respiratory tract secretions. & Correct &  \\
\hline
scot\_09\_25 & 16 & Target Medication: Nystatin 100,000 units/ml oral suspension Recommendation: Initiate 1ml topically to the oral cavity four times daily after gentle mouth care for oral candidiasis. & Correct &  \\
\hline
scot\_09\_26 & 1 & Target Medication: Gabapentin 300mg capsules Recommendation: Discontinue therapy (therapeutic duplication with pregabalin; high risk of toxicity). & Correct &  \\
\hline
scot\_09\_26 & 2 & Target Medication: Pregabalin 150mg capsules Recommendation: Keep the medication in the regimen. & Acceptable &  \\
\hline
scot\_09\_26 & 3 & Target Medication: Tramadol 50mg capsules Recommendation: Discontinue therapy via a planned taper to reduce additive opioid burden, respiratory depression risk, and serotonin syndrome risk. & Correct &  \\
\hline
scot\_09\_26 & 4 & Target Medication: Co-codamol 30/500mg tablets Recommendation: Review and formulate a reduction plan to taper off the codeine component, transitioning to plain paracetamol for baseline analgesia. & Correct &  \\
\hline
scot\_09\_26 & 5 & Target Medication: Methadone 1mg/ml oral solution Recommendation: Keep the medication in the regimen. & Debatable &  \\
\hline
scot\_09\_26 & 6 & Target Medication: Amitriptyline 50mg tablets Recommendation: Reduce dose to 25–50 mg once daily at night to mitigate excessive sedation, anticholinergic burden, QTc prolongation, and serotonin toxicity risk. & Debatable &  \\
\hline
scot\_09\_26 & 7 & Target Medication: Sertraline 100mg tablets Recommendation: Keep the medication in the regimen. & Correct &  \\
\hline
scot\_09\_26 & 8 & Target Medication: Bisacodyl 5mg tablets Recommendation: Keep the medication in the regimen. & Debatable &  \\
\hline
scot\_09\_26 & 9 & Target Medication: Naloxone (take-home naloxone kit) Recommendation: Prescribe a take-home naloxone kit (nasal spray or pre-filled syringe) and provide overdose education to the patient. & Correct &  \\
\hline
\end{longtable}
\begin{table}[!ht]
\caption{Gap rows on the returned sheets.}
\label{tab:appgaps}
\centering
\footnotesize
\begin{tabular}{|l|l|p{5.5cm}|p{6.5cm}|}
\hline
Arm & Case & Missed item & What should have been done \\
\hline
Staged & scot\_09\_01 & Tramadol & Descontinuado, deveria ter sido acrescentado algum analgésico fixo. \\
\hline
Staged & scot\_09\_01 & Melatonina & a adição é controvérsia \\
\hline
Staged & scot\_09\_02 & Clopidogrel - PRM1 & suspensão \\
\hline
Staged & scot\_09\_02 & Levotiroxina & precisa de exame para saber se essa dose já não está muito alta. A dose inicial deve ser 1.6 mcg/kg. \\
\hline
Staged & scot\_09\_19 & PRM 2 & Deveria ser iniciado tratamento para HAS \\
\hline
Baseline & scot\_09\_12 & Falta medicamentos para adição de novos medicamentos. & --- \\
\hline
\end{tabular}
\end{table}
\endgroup
